% Les foyers moteurs de l’économie mondiale — Géographie Tle C — Cours signé OnBuch+
% Programme officiel MINESEC. Compiler avec : tectonic N13-foyers-moteurs-economie-mondiale.tex
\newcommand{\DOCMATIERE}{Géographie}
\newcommand{\DOCNIVEAU}{Tle C}
\newcommand{\DOCMODULE}{Module 2 : La libéralisation des échanges (la mondialisation)}
\newcommand{\DOCLECON}{N13}
\newcommand{\DOCTITRE}{Les foyers moteurs de l’économie mondiale}
% OnBuch+ — préambule « cours signés » ENRICHI (compilé avec tectonic / XeLaTeX)
% Système visuel « Pop » d'OnBuch+ : fond crème (jamais blanc), bordures encre
% épaisses, ombres « dures » décalées, titres Archivo Black en capitales.
% Version MATHÉMATIQUES : graphes (pgfplots/tikz), boîtes pédagogiques
% (définition, propriété, méthode, attention, le savais-tu, exercice résolu,
% exercice, TP, bilan), page de garde et signature OnBuch+ ancrée en bas.
%
% Macros attendues dans main.tex AVANT \input{preamble} :
%   \DOCMATIERE  \DOCNIVEAU  \DOCMODULE  \DOCLECON (numéro)  \DOCTITRE
\documentclass[11pt]{article}

\usepackage{fontspec}
\usepackage[french]{babel}
\usepackage[a4paper,margin=2cm,top=3.4cm,bottom=2.8cm,headheight=1.4cm,footskip=1.3cm]{geometry}
\usepackage{amsmath,amssymb,amsfonts}
\usepackage{mathtools} % \xmapsto, \coloneqq... souvent utilisés par les modèles
\usepackage{cancel} % simplifications barrées (bilans de forces)
% \abs{z} (module, valeur absolue) et \norm{u} (norme), utilisés par les modèles
\providecommand{\abs}{}\let\abs\relax\DeclarePairedDelimiter{\abs}{\lvert}{\rvert}
\providecommand{\norm}{}\let\norm\relax\DeclarePairedDelimiter{\norm}{\lVert}{\rVert}
\usepackage[table]{xcolor} % [table] : fournit aussi \rowcolors (alternance de lignes)
\usepackage{pagecolor}
\usepackage{tikz}
\usetikzlibrary{arrows.meta,positioning,calc,shapes.geometric,shapes.misc,shapes.arrows,decorations.pathmorphing,decorations.pathreplacing,decorations.markings,patterns,shadows,mindmap,trees,fit,backgrounds,matrix,intersections,angles,quotes}
\usepackage{pgfplots}
\usepackage[version=4]{mhchem} % équations nucléaires \ce{^{235}_{92}U}
\usepackage[european,siunitx]{circuitikz} % schémas électriques
\pgfplotsset{compat=1.18}
\usepgfplotslibrary{fillbetween,groupplots} % aires entre courbes (calcul intégral)
\usepackage{enumitem}
\usepackage{fancyhdr}
% [most] charge toutes les bibliothèques tcolorbox (listings, minted, raster,
% documentation...) et gonfle inutilement la mémoire TeX sur un document long
% (~300 boîtes) : on ne charge que ce qui est réellement utilisé.
\usepackage[skins,breakable]{tcolorbox}
\usepackage{graphicx}
\graphicspath{{/home/user/t-t/geo-onbuch/images/}}
\usepackage{colortbl}
\usepackage{booktabs}
\usepackage{tabularx}
\usepackage{multirow}
\usepackage{diagbox} % case d'en-tête diagonale des tableaux à double entrée
% Colonnes extensibles centrée / gauche / droite (C, L, R), souvent utilisées par les modèles
\newcolumntype{C}{>{\centering\arraybackslash}X}
\newcolumntype{L}{>{\raggedright\arraybackslash}X}
\newcolumntype{R}{>{\raggedleft\arraybackslash}X}
\usepackage{array}
\usepackage{multicol}
\usepackage{siunitx}
% parse-numbers=false : accepte aussi \SI{2,5 \times 10^{-3}}{...}
\sisetup{locale=FR,output-decimal-marker={,},group-minimum-digits=4,parse-numbers=false}
\DeclareSIUnit\ohm{\ensuremath{\symup{\Omega}}} % U+2126 absent de la police
\DeclareSIUnit\division{div} % réglages d'oscilloscope (ms/div, V/div)
\DeclareSIUnit\tr{tr} % tours (vitesse de rotation, ex. \SI{1500}{\tr\per\minute})
\DeclareSIUnit\molar{M} % concentration molaire (ex. \SI{3}{\molar}, \SI{1}{\micro\molar})
\DeclareSIUnit\an{an} % année (le modèle confond parfois avec \year, un compteur TeX) — ex. \SI{5}{\kilo\meter\per\an}
\DeclareSIUnit\FCFA{FCFA} % franc CFA (ex. \SI{500}{\FCFA\per\kilo\gram})
\DeclareSIUnit\degreeBrix{\textdegree Brix} % teneur en sucre d'un sirop/jus (réfractométrie)
\DeclareSIUnit\month{mois} % le modèle utilise parfois \month, un compteur TeX, comme unité de durée
\DeclareSIUnit\microliter{\micro\liter} % le modèle écrit parfois \microliter en un seul token
\DeclareSIUnit\million{millions} % dénombrement (ex. \SI{15}{\million\per\mL} spermatozoïdes)
\usepackage{qrcode}
% \degree : commande fréquemment utilisée par les modèles pour un angle ou une
% température en dehors de \SI{}{} (non fournie par les paquets chargés).
\providecommand{\degree}{^{\circ}}
% \qed (fin de démonstration), utilisé par les modèles sans amsthm
\providecommand{\qed}{\hfill$\square$}
% \barre : double barre des valeurs interdites dans un tableau de variations (tkz-tab)
\providecommand{\barre}{\ensuremath{\|}}
% environnement proof (amsthm non chargé) utilisé par les modèles
\ifdefined\proof\else\newenvironment{proof}[1][Démonstration]{\par\noindent\textit{#1.}\ }{\qed\par}\fi
\usepackage{lastpage}
\usepackage{hyperref}
\hypersetup{hidelinks}

% ============================================================
% Palette « Pop » OnBuch+ — mêmes hex que la classe P (plus_ui.dart)
% ============================================================
\definecolor{poporange}{HTML}{F59321}
\definecolor{popdark}{HTML}{E07A0C}
\definecolor{popcream}{HTML}{FFF6E8}
\definecolor{popink}{HTML}{1C1714}
\definecolor{poppurple}{HTML}{7A5AE0}
\definecolor{popgreen}{HTML}{2E9E6A}
\definecolor{poppink}{HTML}{E0457B}
\definecolor{popblue}{HTML}{2F7DE1}
\definecolor{popgold}{HTML}{C98A12}
\definecolor{popmuted}{HTML}{8A7F76}
\definecolor{popline}{HTML}{EADFD2}
\definecolor{popgray}{HTML}{8A7F76}
\colorlet{popgrey}{popgray}
% Alias tolérants (le modèle nomme parfois les couleurs différemment)
\colorlet{popwhite}{white}
\colorlet{popblack}{popink}
\colorlet{popred}{poppink}
\colorlet{popyellow}{popgold}
% Teintes claires (fonds de tableaux/encadrés) — le modèle nomme parfois
% les couleurs avec un suffixe L (ex. popblueL) sans les avoir définies.
\colorlet{poporangeL}{poporange!20}
\colorlet{popdarkL}{popdark!20}
\colorlet{poppurpleL}{poppurple!20}
\colorlet{popgreenL}{popgreen!20}
\colorlet{poppinkL}{poppink!20}
\colorlet{popblueL}{popblue!20}
\colorlet{popgoldL}{popgold!20}
\colorlet{popredL}{popred!20}
\colorlet{popgrayL}{popgray!20}
% Teintes claires pour fonds de tableaux / zones
\colorlet{popblueL}{popblue!10!white}
\colorlet{poporangeL}{poporange!14!white}
\colorlet{popgreenL}{popgreen!12!white}
\colorlet{poppurpleL}{poppurple!12!white}
\colorlet{poppinkL}{poppink!12!white}
\colorlet{popgoldL}{popgold!14!white}

\pagecolor{popcream}

% ============================================================
% Typographie
% ============================================================
\defaultfontfeatures{Ligatures=TeX}
\setmainfont{PlusJakartaSans-Medium.ttf}[Path=../../../fonts/,BoldFont=PlusJakartaSans-Bold.ttf]
% Maths en sans-serif (Fira Math) avec chiffres et lettres droites de Plus
% Jakarta, pour que les formules se fondent dans le texte.
\usepackage[math-style=ISO,bold-style=ISO]{unicode-math}
\setmathfont{FiraMath-Regular.otf}[Path=../../../fonts/]
\setmathfont{PlusJakartaSans-Medium.ttf}[Path=../../../fonts/,range={up/num,up/{Latin,latin},\mathup}]
\usepackage{icomma} % virgule décimale sans espace en mode math
% Caractères mathématiques Unicode absents des polices de texte (Plus Jakarta,
% Archivo Black) : rendus par leur équivalent mathématique, en texte comme en
% formule — sinon ils s'affichent en carré vide (ex. « Arithmétique dans ℤ »).
\usepackage{newunicodechar}
\newunicodechar{ℝ}{\ensuremath{\mathbb{R}}}
\newunicodechar{ℤ}{\ensuremath{\mathbb{Z}}}
\newunicodechar{ℂ}{\ensuremath{\mathbb{C}}}
\newunicodechar{ℕ}{\ensuremath{\mathbb{N}}}
\newunicodechar{ℚ}{\ensuremath{\mathbb{Q}}}
\newunicodechar{𝔻}{\ensuremath{\mathbb{D}}}
\newunicodechar{α}{\ensuremath{\alpha}}
\newunicodechar{β}{\ensuremath{\beta}}
\newunicodechar{γ}{\ensuremath{\gamma}}
\newunicodechar{δ}{\ensuremath{\delta}}
\newunicodechar{ε}{\ensuremath{\varepsilon}}
\newunicodechar{θ}{\ensuremath{\theta}}
\newunicodechar{λ}{\ensuremath{\lambda}}
\newunicodechar{μ}{\ensuremath{\mu}}
\newunicodechar{σ}{\ensuremath{\sigma}}
\newunicodechar{τ}{\ensuremath{\tau}}
\newunicodechar{φ}{\ensuremath{\varphi}}
\newunicodechar{ψ}{\ensuremath{\psi}}
\newunicodechar{ω}{\ensuremath{\omega}}
\newunicodechar{Σ}{\ensuremath{\Sigma}}
% majuscules produites par \MakeUppercase dans les titres (α-aminés -> Α)
\newunicodechar{Α}{\ensuremath{\alpha}}
\newunicodechar{Β}{\ensuremath{\beta}}
\newunicodechar{Γ}{\ensuremath{\Gamma}}
\newunicodechar{Δ}{\ensuremath{\Delta}}
\newunicodechar{Ε}{\ensuremath{\varepsilon}}
\newunicodechar{Θ}{\ensuremath{\Theta}}
\newunicodechar{Λ}{\ensuremath{\Lambda}}
\newunicodechar{Μ}{\ensuremath{\mu}}
\newunicodechar{Τ}{\ensuremath{\tau}}
\newunicodechar{Φ}{\ensuremath{\Phi}}
\newunicodechar{Ψ}{\ensuremath{\Psi}}
\newunicodechar{Ω}{\ensuremath{\Omega}}
\newunicodechar{↦}{\ensuremath{\mapsto}}
\newunicodechar{∈}{\ensuremath{\in}}
\newunicodechar{∉}{\ensuremath{\notin}}
\newunicodechar{∧}{\ensuremath{\wedge}}
\newunicodechar{⇔}{\ensuremath{\Leftrightarrow}}
\newunicodechar{⟺}{\ensuremath{\Longleftrightarrow}}
\newunicodechar{⇒}{\ensuremath{\Rightarrow}}
\newunicodechar{⟹}{\ensuremath{\Longrightarrow}}
\newunicodechar{∘}{\ensuremath{\circ}}
\newunicodechar{∪}{\ensuremath{\cup}}
\newunicodechar{∩}{\ensuremath{\cap}}
\newunicodechar{≡}{\ensuremath{\equiv}}
\newunicodechar{⊂}{\ensuremath{\subset}}
\newunicodechar{⊥}{\ensuremath{\perp}}
\newunicodechar{∃}{\ensuremath{\exists}}
\newunicodechar{∀}{\ensuremath{\forall}}
\newunicodechar{∛}{\ensuremath{\sqrt[3]{\,}}}
\newunicodechar{∜}{\ensuremath{\sqrt[4]{\,}}}
\newunicodechar{⃗}{\ensuremath{{}^{\to}}}
\newunicodechar{ⁿ}{\ifmmode^{n}\else\textsuperscript{\ensuremath{n}}\fi}
\newunicodechar{ˣ}{\ifmmode^{x}\else\textsuperscript{\ensuremath{x}}\fi}
\newunicodechar{ᵇ}{\ifmmode^{b}\else\textsuperscript{\ensuremath{b}}\fi}
\newunicodechar{ʳ}{\ifmmode^{r}\else\textsuperscript{\ensuremath{r}}\fi}
\newunicodechar{ᵘ}{\ifmmode^{u}\else\textsuperscript{\ensuremath{u}}\fi}
\newunicodechar{ʸ}{\ifmmode^{y}\else\textsuperscript{\ensuremath{y}}\fi}
\newunicodechar{ᵃ}{\ifmmode^{a}\else\textsuperscript{\ensuremath{a}}\fi}
\newunicodechar{ᵉ}{\ifmmode^{e}\else\textsuperscript{\ensuremath{e}}\fi}
\newunicodechar{ᵏ}{\ifmmode^{k}\else\textsuperscript{\ensuremath{k}}\fi}
\newunicodechar{ᵖ}{\ifmmode^{p}\else\textsuperscript{\ensuremath{p}}\fi}
\newunicodechar{ᴺ}{\ifmmode^{N}\else\textsuperscript{\ensuremath{N}}\fi}
\newunicodechar{ᶜ}{\ifmmode^{c}\else\textsuperscript{\ensuremath{c}}\fi}
\newunicodechar{ʰ}{\ifmmode^{h}\else\textsuperscript{\ensuremath{h}}\fi}
\newunicodechar{ᵠ}{\ifmmode^{q}\else\textsuperscript{\ensuremath{q}}\fi}
\newunicodechar{ₙ}{\ifmmode_{n}\else\textsubscript{\ensuremath{n}}\fi}
\newunicodechar{ₐ}{\ifmmode_{a}\else\textsubscript{\ensuremath{a}}\fi}
\newunicodechar{ᵢ}{\ifmmode_{i}\else\textsubscript{\ensuremath{i}}\fi}
\newunicodechar{ᵣ}{\ifmmode_{r}\else\textsubscript{\ensuremath{r}}\fi}
\newunicodechar{ₖ}{\ifmmode_{k}\else\textsubscript{\ensuremath{k}}\fi}
\newunicodechar{ᵦ}{\ifmmode_{\beta}\else\textsubscript{\ensuremath{\beta}}\fi}
\newunicodechar{ₓ}{\ifmmode_{x}\else\textsubscript{\ensuremath{x}}\fi}
\newfontfamily\popdisplay{ArchivoBlack-Regular.ttf}[Path=../../../fonts/]
\newfontfamily\poplogo{SpaceGrotesk-Bold.ttf}[Path=../../../fonts/]
\newfontfamily\popsemi{PlusJakartaSans-SemiBold.ttf}[Path=../../../fonts/]
\newfontfamily\popxbold{PlusJakartaSans-ExtraBold.ttf}[Path=../../../fonts/]
\newfontfamily\popxboldls{PlusJakartaSans-ExtraBold.ttf}[Path=../../../fonts/,LetterSpace=3.0]

\setlist[enumerate]{leftmargin=1.7em,itemsep=5pt,topsep=4pt}
\setlist[itemize]{leftmargin=1.7em,itemsep=4pt,topsep=4pt,label={\color{poporange}\textbullet}}
\setlength{\parindent}{0pt}
\setlength{\parskip}{5pt}
\renewcommand{\arraystretch}{1.25}

% pgfplots : style Pop par défaut
\pgfplotsset{
  every axis/.append style={axis line style={popink,line width=1pt},tick style={popink},
    grid style={popline,line width=0.5pt},label style={font=\small},tick label style={font=\footnotesize}},
  popcurve/.style={poporange,line width=1.8pt,no markers,smooth},
}
% Même style disponible dans un \draw TikZ simple (hors pgfplots)
\tikzset{popcurve/.style={poporange,line width=1.8pt}}
\tikzset{popgrid/.style={popline,very thin}}
\colorlet{popcurve}{poporange} % le nom du style est parfois utilisé comme couleur

% ============================================================
% Pill / squiggle
% ============================================================
\newcommand{\poppill}[3]{%
  \tikz[baseline=-0.55ex]\node[draw=popink,line width=1.2pt,rounded corners=2.6pt,%
    fill=#2,inner xsep=6.5pt,inner ysep=3pt]{%
    \popxboldls\fontsize{7}{7}\selectfont\color{#3}\MakeUppercase{#1}};%
}
\newcommand{\popsquiggle}[1]{%
  \tikz[baseline=-0.4ex]{
    \draw[#1,line width=2.6pt,line cap=round]
      plot [smooth] coordinates {(0,0.09) (0.34,0.01) (0.68,0.21) (1.02,0.01) (1.36,0.10)};
  }%
}
% Mot-clé surligné (marqueur orange sous le texte)
\newcommand{\cle}[1]{{\popxbold\color{popdark}#1}}

% ============================================================
% En-tête / pied
% ============================================================
\pagestyle{fancy}
\fancyhf{}
\renewcommand{\headrulewidth}{0pt}
\renewcommand{\footrulewidth}{0pt}
\lhead{\raisebox{-0.15ex}{\poplogo\fontsize{13}{13}\selectfont\color{popink}OnBuch\color{poporange}+}}
\rhead{\poppill{\DOCMATIERE\ \textbullet\ \DOCNIVEAU\ \textbullet\ Leçon \DOCLECON}{white}{popink}}
\lfoot{\popxbold\fontsize{7.6}{7.6}\selectfont\color{popmuted}\MakeUppercase{Cours signé OnBuch+}\ \textbullet\ {\popsemi\DOCTITRE}}
\rfoot{\tikz[baseline=-0.6ex]\node[rounded corners=8.5pt,fill=popink,minimum height=17pt,minimum width=17pt,inner xsep=5pt,inner sep=0]{\popxbold\fontsize{8}{8}\selectfont\color{popcream}\thepage\,/\,\pageref*{LastPage}};}

% ============================================================
% Titres
% ============================================================
\newcommand{\chaptitre}[2]{%
  \begin{tcolorbox}[enhanced,colback=poporange,colframe=popink,boxrule=2.6pt,arc=11pt,
      drop shadow=popink,left=15pt,right=15pt,top=12pt,bottom=11pt,before skip=3pt,after skip=18pt]
    \poppill{#2}{popink}{popcream}\par\vspace{9pt}
    {\raggedright\hyphenpenalty=10000\exhyphenpenalty=10000\popdisplay\fontsize{22}{24}\selectfont\color{popcream}\MakeUppercase{#1}\par}
  \end{tcolorbox}
}
% Section numérotée : pastille orange avec le numéro + titre Archivo
\newcounter{popsec}
\newcommand{\coursec}[1]{%
  \stepcounter{popsec}\setcounter{popsub}{0}%
  \par\vspace{16pt}\noindent
  \begin{minipage}{\linewidth}
  \tikz[baseline=-0.7ex]\node[draw=popink,line width=1.6pt,fill=poporange,rounded corners=4pt,minimum size=22pt,inner sep=0]{\popdisplay\fontsize{12}{12}\selectfont\color{popcream}\thepopsec};%
  \hspace{8pt}{\popdisplay\fontsize{15}{17}\selectfont\color{popink}\MakeUppercase{#1}}\par
  \vspace{4pt}\noindent{\color{poporange}\rule{36pt}{3.6pt}}
  \end{minipage}\par\nopagebreak\vspace{8pt}
}
\newcounter{popsub}
\newcommand{\courssub}[1]{%
  \stepcounter{popsub}%
  \par\vspace{9pt}\noindent{\popxbold\fontsize{11.5}{13}\selectfont\color{popink}{\color{poporange}\thepopsec.\thepopsub}\hspace{6pt}#1}\par\nopagebreak\vspace{4pt}
}

% ============================================================
% Boîtes pédagogiques — même recette : fond blanc, bordure épaisse colorée,
% ombre dure encre, badge plein. Titre optionnel [#1].
% ============================================================
\tcbset{popbase/.style={breakable,enhanced,colback=white,boxrule=2.3pt,arc=9pt,
  shadow={3pt}{-3pt}{0pt}{popink},left=14pt,right=14pt,top=11pt,bottom=11pt,before skip=11pt,after skip=15pt,
  fontupper=\popsemi\fontsize{10.6}{14.5}\selectfont\color{popink},
  fontlower=\popsemi\fontsize{10.6}{14.5}\selectfont\color{popink}}}
% Coche dessinée (le glyphe ✓ n'existe pas dans les polices du document)
\newcommand{\popcheck}[1]{\tikz[baseline=-0.5ex]\draw[#1,line width=1.6pt,line cap=round,line join=round](-0.13,0.0)--(-0.04,-0.09)--(0.13,0.1);}
\providecommand{\coche}{\popcheck{popgreen}}
\providecommand{\resultat}[1]{\par\smallskip\noindent\fcolorbox{popgreen}{popgreenL}{\parbox{\dimexpr\linewidth-2\fboxsep-2\fboxrule\relax}{\textbf{Résultat.} #1}}\par\smallskip}
\newcommand{\popboxhead}[3]{\poppill{#1}{#2}{white}\ifx\relax#3\relax\else\hspace{6pt}{\popxbold\fontsize{10.6}{12}\selectfont\color{popink}#3}\fi\par\vspace{7pt}}

\newtcolorbox{aretenir}[1][]{popbase,colframe=popgreen,before upper={\popboxhead{À retenir}{popgreen}{#1}}}
\newtcolorbox{exemplebox}[1][]{popbase,colframe=poporange,before upper={\popboxhead{Exemple}{poporange}{#1}}}
\newtcolorbox{definition}[1][]{popbase,colframe=popblue,before upper={\popboxhead{Définition}{popblue}{#1}}}
\newtcolorbox{propriete}[1][]{popbase,colframe=poppurple,before upper={\popboxhead{Propriété}{poppurple}{#1}}}
\newtcolorbox{methode}[1][]{popbase,colframe=popgold,before upper={\popboxhead{Méthode}{popgold}{#1}}}
\newtcolorbox{attention}[1][]{popbase,colframe=poppink,before upper={\popboxhead{Attention}{poppink}{#1}}}
\newtcolorbox{experience}[1][]{popbase,colframe=popblue,colback=popblueL,before upper={\popboxhead{Expérience}{popblue}{#1}}}
% « Le savais-tu ? » : carte violette pleine (contexte, culture, Cameroun)
\newtcolorbox{savaistu}[1][]{popbase,colback=poppurple,colframe=popink,
  fontupper=\popsemi\fontsize{10.4}{14.2}\selectfont\color{white},
  before upper={\poppill{Le savais-tu\,?}{white}{poppurple}\ifx\relax#1\relax\else\hspace{6pt}{\popxbold\color{white}#1}\fi\par\vspace{7pt}}}
% Exercice résolu : énoncé en haut, \tcblower puis la solution sur fond crème
\newcounter{popexo}\newcounter{popexr}
\newtcolorbox{exoresolu}[1][]{popbase,colframe=poporange,colbacklower=popcream,
  segmentation style={popink,line width=1.2pt,dashed},
  before upper={\stepcounter{popexr}\popboxhead{Exercice résolu \thepopexr}{poporange}{#1}},
  before lower={\poppill{Solution}{popink}{popcream}\par\vspace{6pt}}}
% Exercice (non résolu en place) — la correction est dans la partie Corrigés
\newtcolorbox{exercice}[1][]{popbase,colframe=popink,
  before upper={\stepcounter{popexo}\popboxhead{Exercice \thepopexo}{popink}{#1}}}
% Corrigé
\newtcolorbox{corrige}[1][]{popbase,colframe=popgreen,colback=popgreenL,
  before upper={\popboxhead{Corrigé}{popgreen}{#1}}}
% Objectifs / prérequis / bilan
\newtcolorbox{objectifs}{popbase,colframe=popink,colback=poporangeL,
  before upper={\popboxhead{Objectifs de la leçon}{popink}{}}}
\newtcolorbox{prerequis}{popbase,colframe=popink,colback=popblueL,
  before upper={\popboxhead{Prérequis}{popblue}{}}}
\newtcolorbox{bilan}{popbase,colframe=popink,colback=white,boxrule=2.8pt,
  before upper={\popboxhead{Fiche bilan}{poporange}{L'essentiel à connaître pour l'examen}}}
% Figure encadrée avec légende
\newtcolorbox{popfigure}[1][]{enhanced,breakable=false,colback=white,colframe=popline,boxrule=1.4pt,arc=8pt,
  left=10pt,right=10pt,top=10pt,bottom=8pt,before skip=10pt,after skip=12pt,halign=center,
  lower separated=false,halign lower=center,
  fontlower=\popsemi\fontsize{9}{11.5}\selectfont\color{popmuted},
  #1}
\newcommand{\legende}[1]{\par\vspace{6pt}{\popsemi\fontsize{9}{11.5}\selectfont\color{popmuted}#1}}

% ============================================================
% Signature OnBuch+
% ============================================================
\newcommand{\poplogoinline}[1]{{\poplogo\fontsize{#1}{#1}\selectfont\color{popink}OnBuch\color{poporange}+}}
\newcommand{\signatureonbuch}{%
  \begin{tcolorbox}[enhanced,colback=popink,colframe=popink,boxrule=2.2pt,arc=10pt,
      drop shadow=poporange,left=14pt,right=14pt,top=10pt,bottom=10pt,before skip=0pt,after skip=0pt]
    \tikz[baseline=-0.7ex]\node[circle,fill=popgreen,draw=popcream,line width=1.3pt,minimum size=22pt,inner sep=0]{\popcheck{white}};%
    \hspace{9pt}\begin{minipage}[c]{0.62\linewidth}
      {\popxbold\fontsize{12}{13}\selectfont\color{popcream}Signé\ \textbullet\ {\poplogo OnBuch\color{poporange}+}}\par\vspace{2pt}
      {\raggedright\popsemi\fontsize{8}{10.5}\selectfont\color{popline}\MakeUppercase{Équipe pédagogique OnBuch+}\\\MakeUppercase{Conforme au programme officiel MINESEC}\par}
    \end{minipage}\hfill
    {\poplogo\fontsize{22}{22}\selectfont\color{popcream}OnBuch\color{poporange}+}
  \end{tcolorbox}%
}

% ============================================================
% Page de garde
% #1 titre · #2 matière · #3 niveau · #4 module · #5 n° leçon · #6 durée ·
% #7 description / objectifs (texte ou itemize)
% ============================================================
\newcommand{\pagegarde}[7]{%
  \thispagestyle{empty}
  \newgeometry{a4paper,margin=1.7cm,top=1.6cm,bottom=1.4cm}%
  \begin{tikzpicture}[remember picture,overlay]
    \fill[poporange!18] ([xshift=-3.2cm,yshift=-3.6cm]current page.north east) circle (4.6cm);
    \fill[poppurple!14] ([xshift=2.4cm,yshift=7.6cm]current page.south west) circle (3.8cm);
  \end{tikzpicture}%
  {\poplogo\fontsize{30}{30}\selectfont\color{popink}OnBuch\color{poporange}+}\hfill
  \raisebox{0.6ex}{\poppill{Cours signé \textbullet\ Premium}{popink}{popcream}}\par
  \vspace{1pt}\popsquiggle{poppurple}\par
  \vspace{8pt}
  \begin{tcolorbox}[enhanced,colback=poporange,colframe=popink,boxrule=2.8pt,arc=13pt,
      drop shadow=popink,left=18pt,right=18pt,top=12pt,bottom=13pt,after skip=10pt]
    \poppill{#2 \textbullet\ #3}{popink}{popcream}\hspace{5pt}\poppill{#4}{white}{popink}\par\vspace{8pt}
    {\popdisplay\fontsize{14}{14}\selectfont\color{popink}LEÇON #5}\par\vspace{4pt}
    {\raggedright\hyphenpenalty=10000\exhyphenpenalty=10000\popdisplay\fontsize{25}{27}\selectfont\color{popcream}\MakeUppercase{#1}\par}
    \vspace{8pt}
    {\popxbold\fontsize{9.5}{9.5}\selectfont\color{popink}\MakeUppercase{Durée conseillée : #6}}
  \end{tcolorbox}
  \begin{tcolorbox}[enhanced,colback=white,colframe=popink,boxrule=2.2pt,arc=10pt,
      drop shadow=popink,left=16pt,right=16pt,top=10pt,bottom=10pt,after skip=8pt]
    \poppill{Dans cette leçon}{poppurple}{white}\par\vspace{5pt}
    \popsemi\fontsize{9.3}{12.6}\selectfont\color{popink}#7
  \end{tcolorbox}
  \noindent\begin{minipage}[t]{0.487\linewidth}
    \begin{tcolorbox}[enhanced,colback=popgreen,colframe=popink,boxrule=2.2pt,arc=10pt,
        drop shadow=popink,left=11pt,right=11pt,top=10pt,bottom=10pt,height=3.1cm,valign=center]
      \tcbox[colback=white,colframe=popink,boxrule=1.3pt,arc=5pt,boxsep=0pt,nobeforeafter,
        left=3pt,right=3pt,top=3pt,bottom=3pt,valign=center]{\qrcode[height=1.9cm]{https://play.google.com/store/apps/details?id=com.luvvix.onbuch&hl=fr}}%
      \hspace{8pt}\begin{minipage}[c]{0.5\linewidth}
        {\popdisplay\fontsize{11}{12.5}\selectfont\color{white}\MakeUppercase{Télécharge OnBuch}}\par\vspace{4pt}
        {\raggedright\popsemi\fontsize{8.4}{11}\selectfont\color{white}Gratuit \textbullet\ Google Play\\Tuteur IA, annales, concours, résultats\par}
      \end{minipage}
    \end{tcolorbox}
  \end{minipage}\hfill
  \begin{minipage}[t]{0.487\linewidth}
    \begin{tcolorbox}[enhanced,colback=poppurple,colframe=popink,boxrule=2.2pt,arc=10pt,
        drop shadow=popink,left=11pt,right=11pt,top=10pt,bottom=10pt,height=3.1cm,valign=center]
      \tikz[baseline=-0.7ex]\node[circle,fill=white,draw=popink,line width=1.3pt,minimum size=1.6cm,inner sep=0]{\popdisplay\fontsize{24}{24}\selectfont\color{poppurple}+};%
      \hspace{8pt}\begin{minipage}[c]{0.6\linewidth}
        {\popdisplay\fontsize{11}{12.5}\selectfont\color{white}\MakeUppercase{Passe à OnBuch+}}\par\vspace{4pt}
        {\raggedright\popsemi\fontsize{8.4}{11}\selectfont\color{white}Tous les cours signés, Tuteur IA illimité, fiches PDF — onglet OnBuch+ de l'app\par}
      \end{minipage}
    \end{tcolorbox}
  \end{minipage}
  \vspace{8pt}
  \popcontenu
  \vfill
  \signatureonbuch
  \restoregeometry
  \newpage
}

% Rangée « Ce cours contient » (page de garde)
\newcommand{\poptile}[3]{% #1 couleur #2 gros texte #3 légende
  \begin{tcolorbox}[enhanced,colback=#1,colframe=popink,boxrule=2pt,arc=9pt,shadow={2.5pt}{-2.5pt}{0pt}{popink},
      left=6pt,right=6pt,top=9pt,bottom=9pt,halign=center,valign=center,height=1.75cm,width=\linewidth,nobeforeafter]
    {\popdisplay\fontsize{12.5}{13}\selectfont\color{white}\MakeUppercase{#2}}\par\vspace{3pt}
    {\centering\popsemi\fontsize{7.8}{9.5}\selectfont\color{white}#3\par}
  \end{tcolorbox}}
\newcommand{\popcontenu}{%
  \par\noindent{\popxboldls\fontsize{7.5}{7.5}\selectfont\color{popmuted}CE COURS SIGNÉ CONTIENT}\par\vspace{7pt}
  \noindent\begin{minipage}[t]{0.235\linewidth}\poptile{popblue}{Cours}{complet, illustré, pas à pas}\end{minipage}\hfill
  \begin{minipage}[t]{0.235\linewidth}\poptile{poporange}{Méthodes}{et exercices résolus}\end{minipage}\hfill
  \begin{minipage}[t]{0.235\linewidth}\poptile{poppink}{Exercices}{type examen + corrigés}\end{minipage}\hfill
  \begin{minipage}[t]{0.235\linewidth}\poptile{popgreen}{Fiche bilan}{l'essentiel à retenir}\end{minipage}\par}

% Sommaire
\newtcolorbox{sommaire}{enhanced,colback=white,colframe=popink,boxrule=2.2pt,arc=10pt,
  drop shadow=popink,left=18pt,right=18pt,top=13pt,bottom=13pt,before skip=6pt,after skip=20pt,
  before upper={\poppill{Sommaire}{poppurple}{white}\par\vspace{9pt}\popsemi\fontsize{10.6}{14.5}\selectfont\color{popink}}}

% Signature de fin de cours — poussée tout en bas de la dernière page
\newcommand{\findecours}{%
  \par\vfill\nopagebreak
  \begin{center}\popsquiggle{poporange}\end{center}\vspace{4pt}
  \signatureonbuch
}
\AtBeginDocument{\def\llbracket{\mathopen{[\mkern-3.5mu[}}\def\rrbracket{\mathclose{]\mkern-3.5mu]}}} % intervalles d'entiers (glyphe ⟦ absent de la police)
% Glyphes absents de Fira Math : redéfinis à partir de glyphes disponibles
\AtBeginDocument{%
  \def\setminus{\mathbin{\backslash}}\let\smallsetminus\setminus
  \def\vdots{\mathord{\vbox{\baselineskip3.2pt\lineskiplimit0pt\kern4pt\hbox{.}\hbox{.}\hbox{.}}}}%
  \def\ddots{\mathinner{\mkern1mu\raise6.5pt\hbox{.}\mkern2mu\raise3.6pt\hbox{.}\mkern2mu\raise0.7pt\hbox{.}\mkern1mu}}%
  \def\triangle{\mathord{\scalebox{0.9}{$\symup{\Delta}$}}}%
  \def\nabla{\mathord{\raisebox{\depth}{\scalebox{1}[-1]{$\symup{\Delta}$}}}}%
  \def\checkmark{\popcheck{popdark}}%
  \def\star{\mathbin{\ast}}\def\bigstar{\mathord{\ast}}%
  \def\diamond{\mathbin{\scalebox{0.6}{$\Diamond$}}}%
  \def\bigcup{\mathop{\mathchoice{\raisebox{-0.3em}{\scalebox{1.7}{$\cup$}}}{\scalebox{1.25}{$\cup$}}{\cup}{\cup}}\displaylimits}%
  \def\bigcap{\mathop{\mathchoice{\raisebox{-0.3em}{\scalebox{1.7}{$\cap$}}}{\scalebox{1.25}{$\cap$}}{\cap}{\cap}}\displaylimits}%
  \def\lesseqqgtr{\mathrel{\lessgtr}}\def\bigoplus{\mathop{\oplus}\displaylimits}%
}
\providecommand{\ssi}{si et seulement si} % abréviation fréquente
\AtBeginDocument{\providecommand{\wideparen}{\overparen}} % arc de cercle

\begin{document}

\pagegarde{Les foyers moteurs de l’économie mondiale}{Géographie}{Tle C}{Module 2 : La libéralisation des échanges (la mondialisation)}{N13}{2 h}{Cette leçon identifie et caractérise les grands foyers moteurs de l'économie mondiale (Triade, Asie de l'Est, Russie), analyse leurs stratégies de puissance et leurs interdépendances, pour comprendre l'architecture spatiale de la mondialisation libérale.
\vspace{4pt}
\begin{multicols}{2}\small
\begin{itemize}
\item À la fin de cette leçon, je sais définir et localiser la Triade et les nouveaux foyers moteurs (Asie de l'Est, Russie).
\item Je sais comparer les modèles de développement de l'UE, de l'Amérique du Nord, de la Chine, du Japon, de la Corée du Sud et de la Russie à l'aide d'indicateurs (PIB, IDH, commerce, IDE).
\item Je sais expliquer la complémentarité et la concurrence entre ces pôles dans la division internationale du travail (DIT).
\item Je sais construire et légender un croquis de synthèse mondialisant les flux, les pôles et les relais de croissance.
\item Je sais interpréter des documents statistiques (tableaux, graphiques) et cartographiques pour hiérarchiser les puissances économiques.
\item Je sais analyser la place du Cameroun et de l'Afrique centrale dans les stratégies d'approvisionnement et de débouchés de ces foyers.
\end{itemize}\end{multicols}}

\begin{sommaire}
\begin{multicols}{2}
\begin{enumerate}
\item 1. Cadre conceptuel : la Triade et l'émergence de nouveaux moteurs
\item 2. L'Union Européenne : le premier foyer intégré par le droit et le marché unique
\item 3. L'Amérique du Nord : la puissance hégémonique (USA) et ses partenaires intégrés (USMCA)
\item 4. L'Asie de l'Est : l'atelier du monde et l'innovation (Chine, Japon, Corée du Sud)
\item 5. La Russie : foyer énergétique, militaire et pivot eurasiatique
\item 6. Synthèse : Hiérarchie, interdépendances et croquis de synthèse mondiale
\item Activité d'intégration
\item Méthodes et exercices résolus
\item Exercices
\item Corrigés des exercices
\item Fiche bilan
\end{enumerate}
\end{multicols}
\end{sommaire}

\begin{objectifs}
\begin{itemize}
\item À la fin de cette leçon, je sais définir et localiser la Triade et les nouveaux foyers moteurs (Asie de l'Est, Russie).
\item Je sais comparer les modèles de développement de l'UE, de l'Amérique du Nord, de la Chine, du Japon, de la Corée du Sud et de la Russie à l'aide d'indicateurs (PIB, IDH, commerce, IDE).
\item Je sais expliquer la complémentarité et la concurrence entre ces pôles dans la division internationale du travail (DIT).
\item Je sais construire et légender un croquis de synthèse mondialisant les flux, les pôles et les relais de croissance.
\item Je sais interpréter des documents statistiques (tableaux, graphiques) et cartographiques pour hiérarchiser les puissances économiques.
\item Je sais analyser la place du Cameroun et de l'Afrique centrale dans les stratégies d'approvisionnement et de débouchés de ces foyers.
\item Je sais mobiliser le vocabulaire spécifique : IDE, flux, DIT, Triade, émergence, puissance, interface maritime.
\end{itemize}
\end{objectifs}

\begin{prerequis}
\begin{itemize}
\item Notion de mondialisation (Première) : définition, acteurs, flux.
\item Repérage des grands ensembles régionaux (UE, ASEAN, USMCA, BRICS).
\item Indicateurs de développement (PIB, PNB/hab, IDH, Taux d'ouverture).
\item Lecture de cartes choroplèthes et de diagrammes (camemberts, histogrammes).
\item Notion d'IDE (Investissements Directs à l'Étranger) et de firme multinationale.
\end{itemize}
\end{prerequis}

\newpage

\coursec{Cadre conceptuel : la Triade et l'émergence de nouveaux moteurs}

\courssub{Situation d'accroche et problématique}

Imagine-toi un samedi matin au marché \textbf{Mokolo} à Yaoundé. Un commerçant déballe des cartons de smartphones : l'écran vient de Corée du Sud (Samsung/LG Display), le processeur est conçu aux États-Unis (Qualcomm/Apple) mais gravé à Taïwan (TSMC), la batterie est chinoise (CATL/ATL), le design finalisé en Europe (Allemagne/Suède), l'assemblage final fait en Chine (Foxconn/Zhengzhou) ou en Inde (Tata/Pegatron). Le navire qui a apporté le conteneur bat pavillon libérien, est assuré à \textbf{Lloyd's de Londres}, financé par une banque new-yorkaise (JPMorgan, Citi) et suit la route du cap de Bonne-Espérance ou le canal de Suez.

\begin{aretenir}[Problématique de la section]
Comment ces territoires lointains, organisés en pôles de puissance inégaux, structurent-ils l'économie mondiale contemporaine ? Qu'est-ce qu'un « foyer moteur » ? La vieille « Triade » (Amérique du Nord, Europe, Japon) domine-t-elle encore face à l'ascension de la Chine et des pays émergents (BRICS+) ? Et où se situe le Cameroun dans ces flux ?
\end{aretenir}

\courssub{Qu'est-ce qu'un foyer moteur ? Définition et critères}

Un \cle{foyer moteur} n'est pas simplement un pays riche. C'est un espace qui \textbf{commande}, \textbf{innove} et \textbf{irradie}. Pour l'identifier, le géographe croise quatre critères quantitatifs et qualitatifs :

\begin{enumerate}
    \item \textbf{Concentration de richesses :} Part élevée dans le PIB mondial (nominal et PPA), revenu par tête haut, capacité d'investissement (IDE sortants massifs).
    \item \textbf{Pouvoir de décision :} Sièges des \cle{FMN} (Firmes Multinationales) du Fortune Global 500, places boursières de référence (NYSE, NASDAQ, LSE, Euronext, Tokyo, Shanghai, Hong Kong, Shenzhen), sièges des organisations internationales (FMI, Banque Mondiale, OMC, ONU, BAD).
    \item \textbf{Innovation et R\&D :} Dépenses intérieures brutes de R\&D (DIRD) massives (\textgreater 2,5\% du PIB), nombre de brevets déposés (OMPI/PCT), universités de rang mondial (Shanghai Ranking, THE), clusters technologiques (Silicon Valley, Shenzhen, Île-de-France, Bavière, région de Séoul, Boston).
    \item \textbf{Connectivité et flux :} Hubs portuaires (Shanghai, Singapour, Ningbo-Zhoushan, Rotterdam, Los Angeles/Long Beach), aéroports hubs (Dubaï, Londres, Hong Kong, Memphis pour le fret), câbles sous-marins (points d'atterrage), centres de commande logistique (Maersk, MSC, CMA CGM, Hapag-Lloyd, DHL, FedEx, UPS).
\end{enumerate}

\begin{definition}[Foyer moteur]
Espace (pays, région métropolitaine, mégalopole) concentrant des fonctions de \textbf{commande} (sièges de FMN, bourses, centres de R\&D), de \textbf{production à haute valeur ajoutée} (industrie de pointe, services financiers, numériques, pharmaceutiques) et de \textbf{rayonnement mondial} (normes techniques, langues, culture, monnaie de réserve). Il structure les réseaux de la mondialisation.
\end{definition}

\begin{attention}[Piège fréquent : Foyer moteur $\neq$ Grande puissance démographique]
L'Inde (1,44 Md d'hab., 2024) ou le Nigeria (220 M d'hab.) sont des géants démographiques et des marchés immenses, mais ne sont pas encore des foyers moteurs \emph{au sens strict} car leur part dans l'innovation mondiale (brevets PCT), le siège des FMN du Top 500 et la finance internationale reste faible. Ils sont des « relais de croissance » ou « marchés cibles ». La Russie est un cas à part : foyer énergétique/militaire majeur, mais foyer économique secondaire (PIB nominal $\approx$ 1 900 Md \$ en 2023, proche de l'Italie ou du Canada, loin de l'Allemagne ou du Japon).
\end{attention}

\courssub{La Triade historique (années 1980--2000) : architecture de la domination}

Le concept de \cle{Triade} a été popularisé par le stratège japonais \textbf{Kenichi Ohmae} (1985, \emph{Triade Power}). Il désigne les trois pôles qui, de 1945 aux années 2000, ont concentré l'essentiel de l'activité économique mondiale : \textbf{Amérique du Nord} (USA + Canada), \textbf{Europe occidentale} (devenue UE), \textbf{Japon}.

\begin{propriete}[Caractéristiques de la Triade classique (années 1990)]
\begin{itemize}
    \item \textbf{80\% du PIB mondial} (nominal, taux de change).
    \item \textbf{70\% des échanges mondiaux} de marchandises et services.
    \item \textbf{Siège de 85\% des 500 plus grandes FMN} (Fortune Global 500).
    \item Monnaies de réserve : Dollar US (60\%), Deutsche Mark/€ (20\%), Yen (5\%).
    \item Standardisation technique (ISO, IEEE) et normative (OMC, Codex Alimentarius, Bâle) pilotée par ces trois pôles.
\end{itemize}
\end{propriete}

\begin{savaistu}[Genèse du concept]
Kenichi Ohmae, dans \emph{Triade Power} (1985), prédit un monde tripolaire géré par des « états-régions » (Californie, Rhénanie, Kanto) plus que par des nations. Le terme a structuré la géographie économique scolaire pendant 25 ans. Aujourd'hui, on parle de \textbf{G2 (USA--Chine)} pour le leadership stratégique, mais la Triade reste l'outil pertinent pour analyser les \textbf{flux d'IDE} (Investissements Directs à l'Étranger) et de \textbf{brevets} : l'UE et le Japon restent des investisseurs nets massifs vers le Sud.
\end{savaistu}

\courssub{L'érosion relative de la Triade et l'ascension de la Chine : les chiffres}

Observons l'évolution de la part du PIB mondial (nominal, FMI/BM) pour mesurer le basculement. Le graphique ci-dessous utilise des données arrondies pour la lisibilité pédagogique (parts en \%). Le périmètre de l'UE évolue (CEE $\to$ UE15 $\to$ UE27) ; les parts sont rétropolées pour l'UE27 actuelle.

\begin{popfigure}
\begin{tikzpicture}
\begin{axis}[
    width=\linewidth,
    height=8cm,
    ybar stacked,
    bar width=12pt,
    legend style={at={(0.5,-0.28)},anchor=north,legend columns=-1,font=\footnotesize, cells={anchor=west}},
    ylabel={Part du PIB mondial (\% nominal)},
    symbolic x coords={1980,1990,2000,2010,2023},
    xtick=data,
    x tick label style={rotate=45,anchor=east,font=\footnotesize},
    ymin=0,ymax=100,
    ymajorgrids=true,
    grid style=dashed,
    enlarge x limits=0.15,
]
\addplot[fill=popblue!70] coordinates {(1980,26) (1990,26) (2000,30) (2010,23) (2023,26)};
\addlegendentry{USA}
\addplot[fill=popblue!30] coordinates {(1980,28) (1990,25) (2000,24) (2010,22) (2023,17)};
\addlegendentry{UE (UE27 rétropolée)}
\addplot[fill=poporange!70] coordinates {(1980,9) (1990,13) (2000,8) (2010,6) (2023,4)};
\addlegendentry{Japon}
\addplot[fill=popred!70] coordinates {(1980,2) (1990,3) (2000,4) (2010,9) (2023,18)};
\addlegendentry{Chine}
\addplot[fill=popgreen!30] coordinates {(1980,35) (1990,33) (2000,34) (2010,40) (2023,35)};
\addlegendentry{Reste du monde}
\end{axis}
\end{tikzpicture}
\legende{Évolution de la part du PIB nominal mondial (estimations FMI World Economic Outlook / Banque Mondiale, arrondies). La Triade (USA+UE+Japon) passe de ~63\% (1980) à ~47\% (2023). La Chine passe de 2\% à 18\%. Le « Reste du monde » (Inde, ASEAN, Brésil, Golfe, Afrique) dépasse 35\%.}
\end{popfigure}

\begin{exemplebox}[Lecture guidée du graphique]
\begin{enumerate}
    \item \textbf{1980--2000 :} L'hégémonie de la Triade est maximale (pic à 62\% en 2000). Le Japon pèse 8--13\% (miracle économique), l'UE s'élargit.
    \item \textbf{2000--2010 :} Choc de 2008, entrée de la Chine à l'OMC (2001). La part chinoise triple (4\% $\to$ 9\%). La Triade tombe sous les 50\%.
    \item \textbf{2010--2023 :} La Chine (18\%) talonne l'UE (17\%). Les USA résistent (~26\%) grâce à la tech (GAFAM/Big Tech), à l'énergie (pétrole/gaz de schiste) et au dollar. Le Japon poursuit son déclin relatif (vieillissement, déflation, délocalisations). Le « Reste du monde » (Inde ~3,5\%, ASEAN ~3,5\%, Brésil, Golfe, Afrique) dépasse 35\%.
\end{enumerate}
\end{exemplebox}

\begin{methode}[Construire un diagramme en barres empilées (évolution de parts)]
\textbf{Objectif :} Visualiser la redistribution de la richesse mondiale.
\begin{enumerate}
    \item \textbf{Choisir l'indicateur :} PIB nominal (taux de change) pour le poids financier/échange ; PIB PPA (parité de pouvoir d'achat) pour le volume réel d'activité. Préciser le choix en légende.
    \item \textbf{Sélectionner les aires et dates :} Minimum 3 dates (ex. 1980, 2000, 2023). Aires : USA, UE, Japon, Chine, Inde, Reste du monde.
    \item \textbf{Calculer les parts relatives :} $\frac{\text{PIB aire}}{\text{PIB monde}} \times 100$. Vérifier que la somme = 100\% par date.
    \item \textbf{Tracer :} Axe X = dates. Axe Y = 0 à 100\%. Barres cumulatives (stacked). Une couleur par aire (cohérence : bleu USA/UE, orange Japon, rouge Chine, vert Reste).
    \item \textbf{Légender rigoureusement :} Titre informatif, source (FMI, BM, OCDE), année de la donnée, unité (\%), note si PPA ou nominal.
\end{enumerate}
\end{methode}

\courssub{Nouvelle géométrie : multipolarité économique vs unipolarité financière/militaire}

Le monde n'est pas « post-occidental », il est \textbf{multipolaire économiquement} mais reste \textbf{unipolaire financièrement et militairement}.

\begin{center}
\begin{tabularx}{\linewidth}{|l|X|X|}
\hline
\rowcolor{popblueL} \textbf{Dimension} & \textbf{Multipolarité (Diffusion)} & \textbf{Unipolarité / Bipolarité (Persistance)} \\ \hline
\textbf{Production / PIB (PPA)} & Chine (19\%), Inde (7,5\%), ASEAN (5\%), BRICS+ > G7 & USA (15\%), UE (15\%) dominent encore le nominal \\ \hline
\textbf{Commerce} & Chine = 1er partenaire commercial de 120+ pays (dont Cameroun) & Règles OMC encore largement occidentales (Tribunaux, ADPIC) \\ \hline
\textbf{Finance / Monnaie} & Dé-dollarisation partielle (yuan, roupie, or, crypto, accords bilatéraux) & \textbf{Dollar : 58\% réserves, 88\% change, 50\% dette ext.} (FMI 2023) \\ \hline
\textbf{Innovation / Brevets} & Chine = 1er déposant PCT (OMPI 2023, ~70 000) & USA/UE/Japon/Corée = haute valeur (citations, semi-conducteurs, IA) \\ \hline
\textbf{Militaire} & Modernisation Chine, Inde, Turquie, Iran & USA = 40\% dépenses mondiales (SIPRI), 750+ bases, OTAN \\ \hline
\end{tabularx}
\end{center}

\begin{aretenir}[Synthèse géométrique]
\begin{itemize}
    \item \textbf{Triade élargie (Cœur / Core) :} USA -- UE -- Asie de l'Est développée (Japon + Corée du Sud + Taïwan + Chine côtière). Concentre la R\&D, la finance, les sièges FMN, les normes.
    \item \textbf{Pays Émergents / BRICS+ (Relais de croissance) :} Chine (moteur industriel/technologique), Inde (services/démographie), Russie (énergie/armement/géopolitique), Brésil/Indonésie (matières premières), Golfe (capitaux/énergie/hubs logistiques). Nouveaux pôles de décision partielle (Nouvelles Banques de Développement, BRICS Pay).
    \item \textbf{Périphéries actives :} Afrique (matières, démographie, marché), Amérique latine. Le Cameroun est une \textbf{périphérie ressource} (bois, cacao, coton, pétrole, bauxite, minerais) en quête d'industrialisation (ZLECAf, corridors, transformation locale).
\end{itemize}
\end{aretenir}

\courssub{Les interfaces maritimes : les façades de la puissance}

Les foyers moteurs sont des \textbf{puissances maritimes} (sauf Russie, puissance tellurique avec accès Arctique/Pacifique). La maîtrise des \textbf{détroits} et des \textbf{façades} conditionne la projection de puissance.

\begin{popfigure}
\includegraphics[width=\linewidth,height=9.5cm,keepaspectratio]{N13/shipping_routes.jpg}
\legende{Le trafic maritime commercial mondial, qui relie les foyers moteurs (plus le bleu est foncé, plus le trafic est dense ; carte sans noms de lieux). On y lit les flux dominants : Atlantique Nord entre l'Amérique du Nord et l'Europe, Pacifique Nord, axe Méditerranée--Suez--océan Indien--Asie de l'Est. Les trafics se resserrent dans les passages obligés (Panama, Gibraltar, Suez, Malacca), où l'accès aux détroits devient un enjeu de puissance. L'Arctique est très peu fréquenté sur cette carte : la route du Nord reste marginale. \\ Source : Wikimedia Commons, « Shipping routes » (densité du trafic maritime commercial), T. Hengl, CC BY-SA 3.0.}
\end{popfigure}

\begin{exoresolu}[Analyse de la nouvelle géométrie des foyers]
\textbf{Document :} Extrait du rapport \emph{World Investment Report 2023} (CNUCED) -- Flux d'IDE entrants (Md \$).
\begin{center}
\begin{tabular}{|l|c|c|c|}\hline
\textbf{Économie} & \textbf{2015} & \textbf{2020} & \textbf{2022} \\ \hline
USA & 380 & 156 & 285 \\ \hline
UE & 420 & 220 & 180 \\ \hline
Chine (continentale) & 135 & 149 & 189 \\ \hline
Hong Kong & 181 & 119 & 118 \\ \hline
Singapour & 65 & 91 & 141 \\ \hline
Inde & 44 & 64 & 49 \\ \hline
Brésil & 75 & 25 & 86 \\ \hline
Russie & 28 & -10 & -18 \\ \hline
Afrique (total) & 54 & 40 & 45 \\ \hline
\end{tabular}
\end{center}
\textbf{Questions :}
\begin{enumerate}
    \item Construis un diagramme en barres groupées (non empilées) comparant les IDE entrants 2015 vs 2022 pour USA, UE, Chine, Singapour, Afrique.
    \item Interprète : que révèlent ces chiffres sur l'attractivité relative des foyers moteurs ?
    \item Pourquoi la Russie a-t-elle des IDE négatifs en 2022 ?
\end{enumerate}

\textbf{Solution détaillée :}
\begin{enumerate}
    \item \textit{Graphique :} Axe X : 5 économies. Axe Y : Md \$. Deux barres par économie (2015 bleu clair, 2022 bleu foncé). Légende : Source CNUCED 2023.
    \item \textit{Interprétation :}
        \begin{itemize}
            \item \textbf{USA :} Volatilité forte (choc Trump 2017, Covid 2020, reprise 2022 via tech/énergie/biotech). Restent \#1 mondial.
            \item \textbf{UE :} Déclin structurel (fragmentation fiscale/réglementaire, Brexit, guerre Ukraine, vieillissement). Perte d'attractivité relative.
            \item \textbf{Chine + HK :} Stabilité haute (~300 Md\$ combinés). Hub manufacturier mondial + marché intérieur de 1,4 Md consommateurs + montée en gamme (véhicules électriques, batteries, PV).
            \item \textbf{Singapour :} Explosion (+117\% sur la période). Hub financier/tech/logistique de l'ASEAN, neutralité relative, fiscalité attractive, État de droit.
            \item \textbf{Afrique :} Stagnation (45 Md\$ pour 1,4 Md hab. = 32 \$/hab. vs 850 \$/hab. USA). Déficit d'infrastructures, risque perçu élevé, gouvernance perfectible, marchés fragmentés (hors ZLECAf).
        \end{itemize}
    \item \textit{Russie :} Sanctions massives post-invasion Ukraine (février 2022). Gel des avoirs de la banque centrale, retrait des FMN occidentales (Shell, BP, TotalEnergies, McDonald's, IKEA, Renault, Volkswagen). IDE négatifs = désinvestissements (ventes d'actifs bradées à des locaux ou États amis) $>$ nouveaux entrants (Chine, Inde, Turquie, Émirats).
\end{enumerate}
\end{exoresolu}

\courssub{Le Cameroun face aux foyers moteurs : dépendance et marges de manœuvre}

\begin{exemplebox}[Le Cameroun dans les flux de la Triade et de la Chine (Données 2022-2023, arrondies, sources : INS, BCEAO, Douanes, Comtrade)]
\begin{itemize}
    \item \textbf{Exportations (≈ 3 500 Md FCFA / 5,3 Md \$ US) :} Chine (bois, coton, pétrole brut, minerais) $\approx$ 20\% ; UE (Pays-Bas, Italie, Espagne, Belgique : bois, cacao, banane, aluminium, caoutchouc) $\approx$ 35\% ; Inde (pétrole, coton, bois) $\approx$ 10\% ; USA (pétrole, bois) $\approx$ 5\%.
    \item \textbf{Importations (≈ 5 500 Md FCFA / 8,4 Md \$ US) :} Chine (machines, électronique, textiles, acier, véhicules, céramique) $\approx$ 22\% ; UE (France, Allemagne, Belgique : machines, médicaments, véhicules, blé, produits chimiques) $\approx$ 30\% ; Nigeria (pétrole raffiné, informel, ciment) $\approx$ 10\% ; Inde/Thaïlande/USA/ Turquie $\approx$ 15\%.
    \item \textbf{IDE (Stock entrant, estimations) :} France (Bolloré Logistics/Africa Global Logistics, TotalEnergies, Société Générale, Orange, Castel/Brasseries, CFAO) historiquement \#1. Chine (Sinohydro, CHEC, China Railway, Huawei, ZTE, mines de fer Mbalam/Nabeba -- projet, barrages) en forte hausse. USA (Perenco, Kosmos Energy, Google/Starlink naissant, Cargill).
    \item \textbf{Dette extérieure publique (≈ 12 000 Md FCFA, 2023) :} Chine (prêts concessionnels/export buyers credit : barrages, routes, port Kribi, autoroutes) $\approx$ 40\% de la dette bilatérale ; France/AfD ; Marché Eurobonds (dollar/euro, coupons élevés) ; Multilatéraux (FMI, BM, BAD).
\end{itemize}
\end{exemplebox}

\begin{attention}[Erreur de copie : « La Chine est le 1er partenaire du Cameroun »]
Faux en \textbf{valeur totale échangée} (Import + Export) si on compte l'\textbf{UE comme un bloc unique} (souvent 1er partenaire global ~30-35\%). Vrai si on prend les \textbf{pays pris individuellement} (Chine > France > Nigeria > Inde > Pays-Bas). Au Bac, précise toujours : « Pays par pays » ou « Bloc UE ». La \textbf{ZLECAf} (Zone de Libre-Échange Continentale Africaine) vise à faire de l'Afrique un bloc cohérent (1,4 Md hab., PIB ~3 000 Md \$) pour négocier d'égal à égal avec les foyers externes.
\end{attention}

\courssub{Exercices d'entraînement}

\begin{exercice}[Exercice 1 : Caractériser un foyer moteur]
À partir de la définition et des critères vus en 1.2, explique pourquoi \textbf{Singapour} (5,6 M hab., PIB nominal ~500 Md\$) est considéré comme un foyer moteur majeur, alors que le \textbf{Nigeria} (220 M hab., PIB nominal ~470 Md\$) ne l'est pas encore. Structure ta réponse en 4 arguments (Richesse, Décision, Innovation, Connectivité).
\end{exercice}

\begin{exercice}[Exercice 2 : Croquis de synthèse mondiale (TP)]
Sur une carte du monde en projection Robinson (fournie en annexe ou tracée au tableau), réalise un croquis légendé intitulé : \textbf{« Les foyers moteurs de l'économie mondiale vers 2023 »}.
\textbf{Consignes :}
\begin{enumerate}
    \item Représente les 3 pôles de la Triade historique (hachures bleues) + Chine (hachures rouges).
    \item Localise et nomme : 4 métropoles mondiales (New York, Londres, Tokyo, Shanghai), 2 hubs émergents (Singapour, Dubaï), 1 hub africain (Johannesburg/Lagos/Nairobi).
    \item Trace 3 flèches de flux majeurs (Atlantique Nord, Pacifique Nord, Europe--Asie via Suez) avec épaisseur proportionnelle au volume commercial.
    \item Ajoute la Route Arctique (pointillés) et le Cap de Bonne-Espérance (alternative).
    \item Légende organisée : Foyers / Métropoles / Flux / Interfaces (détroits).
\end{enumerate}
\end{exercice}

\begin{corrige}[Exercice 1]
\textbf{Singapour est un foyer moteur car :}
\begin{enumerate}
    \item \textbf{Richesse :} PIB/hab. ~90 000 \$ (nominal), ~130 000 \$ (PPA), réserves de change > 300 Md\$, IDE sortants massifs (Temasek, GIC - fonds souverains).
    \item \textbf{Décision :} Siège régional de 4 000+ FMN, place financière \#3 Asie (après Tokyo, Shanghai), hub juridiction arbitrage international (SIAC, ICC), droit des affaires anglosaxon.
    \item \textbf{Innovation :} R\&D > 2,5\% PIB, stratégie « Smart Nation », biotech (Biopolis), fintech, port autonome nouvelle génération (Tuas), hub câbles sous-marins (équidistant Europe-Asie).
    \item \textbf{Connectivité :} Port \#2 mondial conteneurs (transbordement \#1), Aéroport Changi \#1 connectivité/awards, détroit de Malacca (40\% commerce maritime mondial, 80\% pétrole Chine/Japon/Corée).
\end{enumerate}
\textbf{Le Nigeria n'est pas (encore) un foyer moteur car :}
\begin{enumerate}
    \item \textbf{Richesse :} PIB/hab. ~2 100 \$, dépendance pétrole (90\% export, 50\% recettes budgétaires), pauvreté monétaire 40\% (seuil 2,15 \$/j), inflation structurelle.
    \item \textbf{Décision :} Sièges FMN rares (Dangote = exception nationale, pas hub régional), places boursières secondaires (Lagos NSE), instabilité politique/sécuritaire (Boko Haram, banditisme, tensions ethniques).
    \item \textbf{Innovation :} R\&D < 0,2\% PIB, fuite des cerveaux massive (diaspora qualifiée), électricité instable (générateurs = coût majeur), éducation supérieure sous-financée.
    \item \textbf{Connectivité :} Ports saturés (Apapa/Lagos, temps d'attente semaines), infrastructures routières/ferroviaires déficientes (corridors Lagos-Kano, Port Harcourt-Maiduguri), bande passante coûteuse (câbles MainOne, Glo-1, mais dernier kilomètre faible).
\end{enumerate}
\textit{Nuance :} Nigeria = « Géant démographique », 1er marché de consommation Afrique, hub culturel mondial (Nollywood \#2 volume, Afrobeats), potentiel futur immense si gouvernance/énergie/éducation/sécurité s'améliorent (projet « Nigeria Agenda 2050 »).
\end{corrige}

\begin{corrige}[Exercice 2 -- Éléments attendus pour le croquis]
\begin{itemize}
    \item \textbf{Base cartographique :} Contours continents simplifiés (polygones grossiers), méridien de Greenwich, Équateur, Tropiques, Cercles polaires. Flèche Nord (haut de carte).
    \item \textbf{Zones (hachures/couleurs) :}
        \begin{itemize}
            \item Triade (USA, Canada, UE, Japon) = \textcolor{popblue}{\rule{5pt}{5pt}} (hachures verticales bleues).
            \item Chine = \textcolor{popred}{\rule{5pt}{5pt}} (hachures horizontales rouges).
            \item Inde, ASEAN, Brésil, Golfe, Afrique du Sud = \textcolor{poporange}{\rule{5pt}{5pt}} (points ou hachures fines orange) = « Relais de croissance ».
        \end{itemize}
    \item \textbf{Points (Métropoles mondiales) :} \textcolor{popdark}{\textbullet} New York, Los Angeles, Chicago ; Londres, Paris, Francfort, Amsterdam ; Tokyo, Osaka ; Shanghai, Shenzhen, Pékin, Hong Kong ; Singapour ; Séoul ; Dubaï ; Mumbai ; São Paulo ; Johannesburg. Noms en police 8-9 pt, lisibles, non chevauchants.
    \item \textbf{Flèches (Flux commerciaux majeurs) :}
        \begin{itemize}
            \item \textcolor{popblue}{\textbf{---\textgreater}} Épaisse (3-4 pt) : Atlantique Nord (UE $\leftrightarrow$ USA, investissements, services, haute tech).
            \item \textcolor{popred}{\textbf{---\textgreater}} Épaisse (3-4 pt) : Pacifique Nord (Chine/Japon/Corée/Taïwan $\leftrightarrow$ USA, biens manufacturés, composants électroniques).
            \item \textcolor{poporange}{\textbf{---\textgreater}} Moyenne (2 pt) : Europe $\leftrightarrow$ Asie (via Suez / Cap) : machines, chimie, biens de conso, énergie.
            \item \textcolor{popmuted}{\textbf{---\textgreater}} Fine (1 pt) : Afrique $\leftrightarrow$ Chine/UE/Inde (matières premières / biens manufacturés) ; Amérique Latine $\leftrightarrow$ Chine/USA.
        \end{itemize}
    \item \textbf{Interfaces (Détroits/Canaux stratégiques) :} Nommer obligatoirement : Malacca, Hormuz, Bab el-Mandeb, Suez, Panama, Bosphore/Dardanelles, Détroit de Taïwan, Détroit de Béring (Arctique).
    \item \textbf{Légende structurée (obligatoire, encadrée, coin bas droit) :}
        \begin{enumerate}
            \item Foyers moteurs (Triade historique, Chine, Relais de croissance).
            \item Métropoles mondiales (Commande, Finance, Innovation, Logistique).
            \item Flux commerciaux majeurs (Épaisseur $\propto$ volume/valeur).
            \item Points de passage stratégiques (Détroits, Canaux, Caps).
        \end{enumerate}
\end{itemize}
\end{corrige}

\coursec{L'Union Européenne : le premier foyer intégré par le droit et le marché unique}

Dans la section précédente, nous avons vu que l'économie mondiale s'organise autour de la \cle{Triade} (Amérique du Nord, Europe occidentale, Asie de l'Est), désormais concurrencée par de nouveaux moteurs comme la Chine. Parmi ces pôles, l'\cle{Union européenne} (UE) occupe une place singulière : c'est le seul foyer moteur qui ne soit pas un État, mais une construction politique et juridique inédite, bâtie volontairement par des nations qui ont choisi de partager une partie de leur souveraineté. Comprendre l'UE, c'est comprendre comment le \emph{droit} et le \emph{marché} peuvent créer une puissance économique de premier plan.

\courssub{Une construction historique par étapes : de la CECA aux 27}

\textbf{Situation d'accroche.} En 1950, l'Europe sort ruinée de deux guerres mondiales nées pour une large part de la rivalité franco-allemande autour du charbon et de l'acier de la Ruhr et de la Sarre. L'idée de Jean Monnet et Robert Schuman est simple : plaçons ces industries de guerre sous une autorité commune, et la guerre deviendra « non seulement impensable, mais matériellement impossible » (déclaration Schuman du 9 mai 1950). C'est l'acte de naissance de l'Europe unie.

\begin{experience}[Reconstituer la chronologie de la construction européenne]
\textbf{Matériel :} une frise chronologique vierge, les dates et événements suivants en désordre : Traité de Rome, CECA, Acte Unique, Maastricht, élargissement de 2004, Brexit.\par
\textbf{Protocole :} replacer chaque événement sur la frise et associer à chacun sa « logique » (paix, marché, monnaie, élargissement).\par
\textbf{Observations et interprétation :}
\begin{itemize}
\item \textbf{1951 — CECA} (Communauté européenne du charbon et de l'acier, traité de Paris) : 6 pays (France, RFA, Italie, Belgique, Pays-Bas, Luxembourg) mettent en commun le charbon et l'acier. Logique : la \cle{paix} par l'interdépendance économique.
\item \textbf{1957 — Traité de Rome} : création de la Communauté économique européenne (CEE) : union douanière et marché commun. Logique : le \cle{marché}.
\item \textbf{1986 — Acte Unique européen} : achèvement du grand marché intérieur au 1\textsuperscript{er} janvier 1993 (les 4 libertés de circulation).
\item \textbf{1992 — Traité de Maastricht} : naissance de l'Union européenne et feuille de route vers la monnaie unique, l'\cle{euro} (pièces et billets mis en circulation en 2002).
\item \textbf{2004, 2007, 2013 — Élargissements} : entrée de 10 pays, pour l'essentiel d'Europe centrale et orientale (2004, dont la Pologne), puis de la Roumanie et de la Bulgarie (2007), enfin de la Croatie (2013).
\item \textbf{2020 — Brexit} : le Royaume-Uni quitte l'UE (référendum de 2016, sortie effective le 31 janvier 2020). L'Union compte désormais \textbf{27 États membres}.
\end{itemize}
\textbf{Conclusion :} la construction européenne avance par crises et par traités, toujours dans le sens d'une intégration économique plus profonde.
\end{experience}

\begin{aretenir}[Les grandes étapes à connaître]
CECA (1951) $\rightarrow$ Traité de Rome (1957) $\rightarrow$ Acte Unique (1986) $\rightarrow$ Maastricht (1992) $\rightarrow$ élargissements à l'Est (2004, 2007, 2013) $\rightarrow$ Brexit (2020) : \textbf{27 membres}. Retiens la logique cumulative : \emph{paix} $\rightarrow$ \emph{marché} $\rightarrow$ \emph{monnaie} $\rightarrow$ \emph{élargissement}.
\end{aretenir}

\courssub{Le marché unique : le cœur de la puissance européenne}

\begin{definition}[Marché unique]
Le \cle{marché unique} (ou grand marché intérieur) est un espace sans frontières intérieures assurant la libre circulation des \textbf{biens}, des \textbf{services}, des \textbf{capitaux} et des \textbf{personnes} : ce sont les \textbf{quatre libertés} de circulation, achevées par l'Acte Unique de 1986 (effectif en 1993).
\end{definition}

Concrètement : un camion de bananes du Cameroun déchargé à Rotterdam peut être livré à Varsovie sans aucun contrôle douanier interne ; un ingénieur polonais peut travailler à Munich ; une banque française peut prêter à une entreprise espagnole. Ce marché de \textbf{environ 450 millions de consommateurs} aux revenus élevés fait de l'UE le \textbf{premier marché mondial} par le pouvoir d'achat.

L'euro renforce cette intégration : adopté par \textbf{20 pays} (la \emph{zone euro}, depuis l'entrée de la Croatie en 2023), il élimine les risques de change entre eux et s'impose comme la deuxième monnaie de réserve et de règlement mondiale, loin derrière le dollar mais devant toutes les autres.

\begin{exemplebox}[Quelques chiffres de la puissance européenne]
\begin{itemize}
\item Environ \textbf{450 millions d'habitants} : premier marché de consommation du monde par le pouvoir d'achat.
\item \textbf{Premier exportateur mondial de services} et l'un des tout premiers exportateurs de marchandises (avec la Chine et les États-Unis).
\item \textbf{L'une des toutes premières sources et destinations des IDE} : les flux d'investissements directs entre l'UE et le reste du monde se comptent en centaines de milliards d'euros par an, pour un stock de plusieurs milliers de milliards (ordres de grandeur, variables selon les années).
\item \textbf{Budget de l'UE} : environ \textbf{170 milliards d'euros par an}, soit seulement \textbf{1 \% environ du revenu national brut de l'Union} ; la Politique agricole commune (PAC) et la politique de cohésion (aides aux régions en retard) absorbent ensemble environ \textbf{deux tiers} de ce budget.
\end{itemize}
\end{exemplebox}

\courssub{Un modèle : l'économie sociale de marché et l'industrie de pointe}

L'UE s'est dotée d'un modèle original : l'\cle{économie sociale de marché}, qui combine libre concurrence et protection sociale (assurance maladie, retraites, droit du travail). Sa structure productive est très tertiarisée.

\begin{popfigure}
\begin{center}
\begin{tikzpicture}[scale=1]
% Camembert : Services 76%, Industrie 22%, Agriculture 1.6%
\fill[popblue] (0,0) -- (90:2.2) arc (90:90-273.6:2.2) -- cycle;
\fill[poporange] (0,0) -- (90-273.6:2.2) arc (90-273.6:90-352.8:2.2) -- cycle;
\fill[popgreen] (0,0) -- (90-352.8:2.2) arc (90-352.8:90:2.2) -- cycle;
\node[white, font=\small\bfseries] at (-45:1.2) {Services};
\node[white, font=\small\bfseries] at (-10:1.4) {76 \%};
\node[white, font=\footnotesize\bfseries] at (140:1.1) {Industrie 22 \%};
\node[font=\footnotesize] at (3.2,1.6) {Agriculture 1,6 \%};
\draw[->, thin] (2.9,1.5) -- (1.6,1.1);
\end{tikzpicture}
\end{center}
\legende{Structure du PIB de l'Union européenne par secteur (ordres de grandeur) : agriculture 1,6 \%, industrie 22 \%, services 76 \%. La tertiarisation atteint environ 70 à 76 \% du PIB selon les années : c'est la marque d'une économie post-industrielle.}
\end{popfigure}

L'industrie européenne, minoritaire en valeur ajoutée, reste \emph{de pointe} : aéronautique (\textbf{Airbus}, rival de Boeing), pharmacie (Sanofi, Novartis, Bayer), automobile haut de gamme (Volkswagen, Mercedes-Benz, Stellantis, Renault), luxe (LVMH, premier groupe mondial du secteur), agroalimentaire (Nestlé, Danone).

\begin{popfigure}
\begin{center}
\begin{tikzpicture}
\begin{axis}[
  ybar, bar width=0.9cm,
  width=0.8\linewidth, height=5cm,
  symbolic x coords={Intra-UE, Extra-UE},
  xtick=data,
  ymin=0, ymax=75,
  ylabel={Part des échanges (\%)},
  nodes near coords, every node near coord/.append style={font=\small\bfseries},
  ymajorgrids, grid style={popline!40},
]
\addplot[fill=popblue, draw=popdark] coordinates {(Intra-UE,60) (Extra-UE,40)};
\end{axis}
\end{tikzpicture}
\end{center}
\legende{Part des échanges intra-UE (environ 60 \%) et extra-UE (environ 40 \%) dans le commerce total des États membres (ordre de grandeur). Les Européens commercent d'abord entre eux : c'est la preuve statistique de l'intégration.}
\end{popfigure}

\begin{attention}[Erreur fréquente : « l'UE est un État »]
Non ! L'UE est une \textbf{organisation supranationale unique au monde}, mais ce n'est \emph{pas} un État fédéral comme les États-Unis. Les États membres conservent l'essentiel de leur souveraineté \textbf{fiscale} (les impôts), \textbf{sociale} (retraites, santé) et \textbf{diplomatique et militaire} (chacun son armée et sa politique étrangère). Dire à l'examen que « l'Europe est un pays » est une erreur éliminatoire.
\end{attention}

\courssub{L'« effet Bruxelles » : la puissance par les normes}

\begin{propriete}[L'effet Bruxelles (Anu Bradford, 2012)]
L'UE exporte ses \cle{normes} juridiques et techniques au monde entier, sans traité ni contrainte militaire, par la simple \textbf{taille de son marché} : aucune grande entreprise mondiale ne peut renoncer à vendre à 450 millions de consommateurs riches, donc toutes adoptent les règles européennes, souvent jusque dans leurs autres marchés. Exemples : le \textbf{RGPD} (protection des données, applicable depuis 2018) copié sur tous les continents, le \textbf{chargeur universel USB-C} imposé aux fabricants de smartphones, les \textbf{normes CO\textsubscript{2}} automobiles, le \textbf{Pacte vert} (Green Deal, 2019) visant la neutralité carbone en 2050. Cette propriété n'est pas à démontrer au Bac, mais tu dois savoir l'illustrer par un exemple.
\end{propriete}

\courssub{Un espace inégal : la mégapole européenne et les périphéries}

La richesse européenne se concentre dans une dorsale appelée \cle{mégapole européenne} (ou « banane bleue »), qui s'étire de Londres à Milan en passant par Paris, la vallée du Rhin et la Ruhr. Autour, les périphéries (Est, Sud méditerranéen, régions ultrapériphériques) restent en retard, malgré les fonds de cohésion.

\begin{popfigure}
\includegraphics[width=\linewidth,height=9.5cm,keepaspectratio]{N13/blue_banana.png}
\legende{La mégapole européenne ou « banane bleue » (aplat bleu) : dorsale qui s'étire de l'Angleterre du Sud-Est à l'Italie du Nord en passant par le Benelux, la vallée du Rhin et la Suisse. Autour de cet aplat, la France de l'Ouest et du Sud, l'Espagne, l'Italie du Sud, l'Europe de l'Est et du Nord figurent en gris : ce sont les périphéries, moins dynamiques malgré les fonds de cohésion. La carte ne montre pas les ports ; Rotterdam et Hambourg se situent sur la façade Nord de la dorsale, d'après le cours. \\ Source : Wikimedia Commons, « Blue Banana », ArnoldPlaton, CC BY-SA 3.0.}
\end{popfigure}

\begin{methode}[Analyser une carte des flux intra-européens]
\begin{enumerate}
\item \textbf{Identifier les nœuds} : repère les grandes plates-formes de concentration des flux — ports (\textbf{Rotterdam}, \textbf{Anvers}, \textbf{Hambourg}) et aéroports (\textbf{Paris-CDG}, \textbf{Francfort}).
\item \textbf{Repérer les corridors} : nomme les grands axes de transport — corridor \textbf{rhénan} (Rotterdam–Bâle–Gênes), corridor \textbf{méditerranéen}, corridor \textbf{Baltique–Adriatique}.
\item \textbf{Noter les disparités} : compare la densité des flux à l'Ouest (très dense) et à l'Est (mailles lâches), et cite un projet de rattrapage, par exemple le corridor ferroviaire \textbf{Rail Baltica} (Varsovie–Tallinn–Helsinki).
\item \textbf{Conclure} : relie la carte à la notion de mégapole européenne — les flux dessinent le cœur et marginalisent les périphéries.
\end{enumerate}
\end{methode}

\courssub{Les faiblesses du géant}

\begin{itemize}
\item \textbf{Vieillissement démographique} : l'Europe est le continent le plus âgé du monde (âge médian supérieur à 44 ans), ce qui pèse sur les retraites, la santé et le dynamisme économique.
\item \textbf{Dépendance énergétique} : l'UE importe l'essentiel de son pétrole et de son gaz ; la guerre en Ukraine (2022) a brutalement révélé la dépendance au gaz russe, depuis fortement réduite mais remplacée par d'autres fournisseurs (Norvège, États-Unis, Qatar).
\item \textbf{Fragmentation politique} : pas d'armée commune, des diplomaties nationales divergentes, des décisions souvent prises à l'unanimité, donc lentes.
\item \textbf{Croissance faible} : de l'ordre de \textbf{0,5 à 1,5 \% par an}, contre 2 à 3 \% aux États-Unis et davantage en Chine.
\item \textbf{Écarts internes Nord/Sud et Ouest/Est} : le PIB par habitant de l'Allemagne ou du Luxembourg est plusieurs fois supérieur à celui de la Grèce, de la Bulgarie ou de la Roumanie.
\end{itemize}

\begin{savaistu}[Une puissance incomplète]
L'UE est le \textbf{seul foyer moteur sans armée unifiée ni siège permanent propre au Conseil de sécurité de l'ONU} : seule la France y siège en membre permanent (avec droit de veto). Cette absence de « puissance dure » (\emph{hard power}) limite le poids politique de l'Europe face aux États-Unis, à la Chine et à la Russie, alors même qu'elle est un géant économique et normatif : on parle d'une « puissance civile ».
\end{savaistu}

\courssub{La stratégie d'« autonomie stratégique ouverte » et le lien avec le Cameroun}

Face à la rivalité sino-américaine, l'UE a adopté la doctrine d'\cle{autonomie stratégique ouverte} : rester ouverte au commerce mondial tout en réduisant ses dépendances critiques. Cela se traduit par :
\begin{itemize}
\item la \textbf{souveraineté technologique} : production européenne de puces électroniques (\emph{Chips Act}, adopté en 2023), cloud européen, hydrogène vert ;
\item des \textbf{accords de libre-échange} : CETA avec le Canada (appliqué provisoirement depuis 2017), accord avec le Mercosur (négocié), et accords avec l'Afrique via les APE et l'accord post-Cotonou dit \textbf{accord de Samoa} (2023) avec les pays ACP (Afrique, Caraïbes, Pacifique).
\end{itemize}

\textbf{Lien avec le Cameroun.} Le Cameroun a signé avec l'UE un \cle{Accord de partenariat économique} (APE) : accord intérimaire signé en 2009, entrée en vigueur en 2014. Il garantit aux exportations camerounaises — \textbf{banane} (PHP/Compagnie fruitière), \textbf{cacao}, \textbf{bois}, \textbf{aluminium} (Alucam) — un accès au marché européen \textbf{sans droits de douane ni contingents}. En échange, le Cameroun ouvre progressivement son marché aux produits européens. L'UE est ainsi le \textbf{premier partenaire commercial du Cameroun} et son \textbf{premier bailleur de fonds} (via le Fonds européen de développement, FED, qui finance routes, énergie et éducation).

\begin{exoresolu}[Commenter un tableau statistique : le commerce Cameroun–UE]
\textbf{Énoncé.} Document : « En ordre de grandeur, l'UE absorbe près de la moitié des exportations du Cameroun (banane, cacao, bois, aluminium, pétrole brut) et fournit environ un tiers de ses importations (machines, véhicules, produits pharmaceutiques). L'UE est le premier partenaire commercial et le premier bailleur de fonds du Cameroun. » À partir du document et de tes connaissances : 1) montre que le Cameroun est inséré dans la mondialisation via l'UE ; 2) relève une limite de cette relation.
\tcblower
\textbf{Solution détaillée.}\par
\textbf{1) Une insertion forte via l'UE.} Le document indique que près de la moitié des exportations camerounaises partent vers l'UE : le marché unique européen, premier marché mondial de 450 millions de consommateurs, est donc le débouché principal des produits camerounais. Cette insertion s'explique par l'APE (2009/2014), qui garantit un accès sans droits de douane ni contingents à la banane, au cacao, au bois et à l'aluminium. En sens inverse, l'UE fournit environ un tiers des importations du Cameroun : biens d'équipement, véhicules et médicaments, c'est-à-dire des produits à forte valeur ajoutée issus de l'industrie de pointe européenne. Le Cameroun est ainsi pleinement intégré aux flux de la mondialisation, avec l'UE comme porte d'entrée principale, complétée par l'aide au développement du FED.\par
\textbf{2) Une limite : la dépendance et l'asymétrie.} Le Cameroun exporte surtout des matières premières peu transformées (cacao en fèves, bois en grumes) et importe des produits manufacturés : il occupe une position défavorable dans la division internationale du travail. De plus, la concentration des exportations sur un seul partenaire rend l'économie camerounaise vulnérable à la conjoncture européenne (croissance faible de 0,5 à 1,5 \%). D'où la nécessité de la transformation locale et de la diversification des partenaires (Chine, marché africain de la ZLECAf).
\end{exoresolu}

\begin{exercice}[Construire un croquis de synthèse]
Sur une ébauche de carte de l'Europe, construis un croquis intitulé « L'Union européenne : un foyer moteur intégré mais inégal ». Tu placeras : la mégapole européenne, deux ports majeurs, deux périphéries, et un flux d'importation d'hydrocarbures venant de l'Est. N'oublie ni le titre, ni l'orientation, ni la légende organisée.
\end{exercice}

\begin{corrige}[Exercice : croquis de synthèse]
Éléments attendus : titre complet et orienté ; flèche du nord ; la mégapole figurée par un trait épais de Londres à Milan via Paris et la Rhine-Ruhr ; Rotterdam et Hambourg (ou Anvers) par des points nommés sur la façade Nord ; les périphéries Est (Pologne, Roumanie) et Sud (Grèce, sud de l'Italie) par des aplats ou des flèches ; une flèche large venue de l'Est légendée « importations de gaz et de pétrole (dépendance énergétique) » ; une légende numérotée et hiérarchisée. Le croquis doit rester simple : un contour simplifié cohérent suffit, la rigueur de la légende prime sur la précision du trait.
\end{corrige}

\begin{aretenir}[L'essentiel de la section]
L'UE, construite par étapes de la CECA (1951) à Maastricht (1992) et aux élargissements (2004--2013), réunit \textbf{27 États} autour d'un \textbf{marché unique} de 450 millions de consommateurs et, pour 20 d'entre eux, de l'\textbf{euro}. Premier marché mondial et acteur majeur des IDE, elle impose ses normes par l'\textbf{effet Bruxelles}. Mais sa richesse se concentre dans la \textbf{mégapole européenne}, sa croissance est faible, sa population vieillit, elle dépend de l'étranger pour son énergie et reste une \textbf{puissance politique incomplète}. Pour le Cameroun, elle est le premier partenaire commercial (APE) et le premier bailleur (FED). La section suivante examinera un foyer d'un type très différent : l'Amérique du Nord, puissance hégémonique construite autour d'un seul État.
\end{aretenir}

\coursec{L'Amérique du Nord : la puissance hégémonique (USA) et ses partenaires intégrés (USMCA)}

Après avoir analysé l'Union européenne, foyer moteur original par son intégration \emph{juridique et politique} (marché unique, euro, acquis communautaire), nous basculons vers le deuxième pilier historique de la Triade : l'Amérique du Nord. Si l'UE construit sa puissance par le droit et la régulation, le foyer nord-américain s'impose par la \textbf{puissance dure} (militaire, financière, technologique) de son leader, les États-Unis, relayée par une intégration économique \emph{de fait} (chaînes de valeur, flux financiers, main-d'œuvre) formalisée par l'USMCA. Observons comment ce « géant aux pieds d'argile » structure la mondialisation.

\courssub{Les États-Unis : le cœur hégémonique de la Triade}

\textbf{Une concentration de puissance sans équivalent.}\par
Les États-Unis demeurent, en 2024, l'unique \emph{superpuissance} planétaire. Ils concentrent les attributs classiques de la puissance : le 1er PIB mondial nominal (\SI{~27 400}{Md\$} en 2023, soit \SI{~25}{\%} du PIB mondial pour \SI{4,2}{\%} de la population), le 1er budget militaire (\SI{~916}{Md\$} budget 2023, plus que les 10 suivants réunis), la 1re place financière (Wall Street, dollar roi), et le leadership technologique (Silicon Valley, GAFAM/NATU). À cela s'ajoute depuis 2018 le statut de 1er producteur mondial d'hydrocarbures (pétrole + gaz de schiste), garantissant une indépendance énergétique stratégique.

\begin{popfigure}
\begin{tikzpicture}
\begin{axis}[
width=\linewidth,
height=7cm,
ybar,
legend pos=north west,
ylabel={PIB nominal (Md \$ 2023)},
symbolic x coords={USA,Chine,UE,Japon,Allemagne,Inde,UK,France,Russie,Canada},
xtick=data,
xticklabel style={rotate=45,anchor=east, font=\tiny},
ymax=29000,
bar width=0.6cm,
fill=popblue!70
]
\addplot coordinates {
(USA,27400)
(Chine,17700)
(UE,18300)
(Japon,4210)
(Allemagne,4460)
(Inde,3570)
(UK,3340)
(France,3030)
(Russie,1900)
(Canada,2140)
};
\end{axis}
\end{tikzpicture}
\legende{Figure 1 : Hiérarchie des principales économies mondiales (PIB nominal, FMI 2023). L'écart abyssal entre les USA (1er) et la Chine/UE (2e ex-æquo) illustre l'hégémonie économique persistante. Source : FMI, World Economic Outlook, avril 2024.}
\end{popfigure}

\textbf{Le « modèle américain » : libéralisme dérégulé, flexibilité et attractivité.}\par
Ce leadership repose sur un modèle socio-économique spécifique :
\begin{itemize}
\item \textbf{Libéralisme dérégulé :} Faible poids de l'État dans l'économie (dépenses publiques \SI{~36}{\%} PIB vs \SI{>50}{\%} en UE), fiscalité des entreprises compétitive (taux fédéral 21\%), droit des affaires favorable (Delaware).
\item \textbf{Flexibilité du travail :} « \emph{At-will employment} » (licenciement libre), mobilité géographique forte, faible syndicalisation (\SI{~10}{\%}).
\item \textbf{Brain gain (gain de cerveaux) :} 1re destination mondiale des talents (visas H-1B, cartes vertes). \SI{45}{\%} des fondateurs de licornes US sont nés à l'étranger (Musk, Nadella, Pichai, Huang).
\item \textbf{Marché intérieur massif :} \SI{335}{millions} d'habitants, revenu médian élevé, consommation moteur (\SI{68}{\%} du PIB).
\item \textbf{Privilège exorbitant du dollar :} Monnaie de réserve mondiale (\SI{~58}{\%} des réserves de change identifiées FMI T4 2023, \SI{~88}{\%} du Forex BIS 2022), permettant de financer déficits jumeaux (budgétaire + courant) à bas taux.
\end{itemize}

\begin{definition}[Privilège exorbitant (Valéry Giscard d'Estaing, 1965)]
Avantage structurel dont bénéficient les États-Unis du fait que le dollar est la principale monnaie de réserve et de transaction internationale. Ils peuvent emprunter massivement \emph{dans leur propre monnaie} sans risque de change, refinancer leur dette à des taux inférieurs à la croissance nominale, et sanctionner financièrement des États adverses (gel de réserves, exclusion SWIFT).
\end{definition}

\begin{exemplebox}[Chiffres clés de l'innovation et du commerce US (2023)]
\begin{itemize}
\item \textbf{Dépenses R\&D :} \SI{~780}{Md\$} (\SI{3,5}{\%} PIB), 1er mondial en volume absolu (NCSES).
\item \textbf{Brevets PCT (OMPI) :} USA \num{~55 000} (2e derrière Chine \num{~70 000}), mais 1er pour les citations/brevet (qualité).
\item \textbf{Déficit commercial biens :} \SI{~1 060}{Md\$} (record), excédent services \SI{~280}{Md\$}.
\item \textbf{Dette publique fédérale :} \SI{> 34 000}{Md\$} (\SI{> 120}{\%} PIB, T1 2024).
\item \textbf{Coefficient de Gini (revenus disponibles, OCDE) :} \num{0,39} (le plus élevé de l'OCDE) ; \num{0,49} pour les revenus primaires (avant redistribution).
\end{itemize}
\end{exemplebox}

\begin{attention}[Piège fréquent : « Amérique du Nord = États-Unis »]
\emph{Erreur majeure en croquis ou dissertation.} Le Canada (\SI{40}{M} hab, PIB \SI{2 140}{Md\$}, 1er exportateur énergie US, leader IA/quantique, stabilité politique) et le Mexique (\SI{128}{M} hab, PIB \SI{1 810}{Md\$}, 1er partenaire commercial US, « usine » intégrée, démographie jeune) ne sont pas de simples satellites. Ils sont \textbf{co-producteurs} de la compétitivité nord-américaine. Tout croquis de l'Amérique du Nord doit figurer les trois pays, leurs complémentarités et les flux transfrontaliers.
\end{attention}

\courssub{L'intégration nord-américaine : l'USMCA, un bloc inégal mais complémentaire}

\textbf{De l'ALENA (1994) à l'USMCA (2020) : une intégration « par le bas » (marché) plus que « par le haut » (politique).}\par
L'Accord États-Unis–Mexique–Canada (USMCA / ACEUM / T-MEC), entré en vigueur le 1er juillet 2020, renégocie l'ALENA sous pression protectionniste US (Trump). Il renforce : règles d'origine automobile (\SI{75}{\%} contenu régional, \SI{40-45}{\%} main-d'œuvre \num{> 16\$}/h), ouverture marché laitier canadien, protection propriété intellectuelle, clause de révision tous les 6 ans (\emph{sunset clause}). Contrairement à l'UE, \textbf{pas d'institutions supranationales}, pas de monnaie unique, pas de libre circulation des personnes (sauf visas TN professionnels).

\textbf{Une division internationale du travail (DIT) intra-régionale très hiérarchisée.}\par
\begin{center}
\begin{tabularx}{\linewidth}{|l|X|X|X|}
\hline
\rowcolor{popblueL}
\textbf{Pays} & \textbf{Spécialisation principale} & \textbf{Atouts stratégiques} & \textbf{Vulnérabilités} \\
\hline
\textbf{USA} & \textbf{Commandement :} R\&D, design, finance, marque, logiciels, services haute valeur. & Capital risque, universités (Stanford, MIT), marché intérieur, dollar. & Coûts salariaux élevés, déficit commercial, polarisation. \\
\hline
\textbf{Canada} & \textbf{Ressources \& Haute tech :} Énergie (pétrole, gaz, hydroélectricité, uranium), minerais critiques (lithium, cobalt, nickel), bois, céréales, IA (Montréal, Toronto), aérospatiale. & Dépendance export US (\SI{75}{\%}), climat, distance marchés asiatiques, infrastructure Nord faible. \\
\hline
\textbf{Mexique} & \textbf{Usine \& Assemblage :} Automobile, électronique, électroménager, aéronautique (maquiladoras), agro-export (avocats, tomates). & Main-d'œuvre qualifiée bon marché (\num{~4-6\$}/h), proximité US, accords de libre-échange (50+ pays). & Insécurité (cartels), corruption, infrastructure, dépendance US (\SI{80}{\%} export), eau rare. \\
\hline
\end{tabularx}
\end{center}

\begin{popfigure}
\includegraphics[width=\linewidth,height=9.5cm,keepaspectratio]{N13/usmca.png}
\legende{L'espace USMCA (en rouge) : États-Unis, Canada et Mexique forment un ensemble continu, de l'Arctique au Mexique, sur la moitié nord du continent américain. La carte montre les trois pays ensemble et leur frontière commune ; elle ne montre pas les mégalopoles ni les flux. D'après le cours, les mégalopoles états-uniennes (BosWash, ChiPitts, SanSan) pilotent les chaînes de valeur ; la frontière États-Unis--Mexique est à la fois un point de friction (migrations) et d'intégration (maquiladoras) ; le Nord canadien ouvre sur les ressources et les nouvelles routes maritimes. \\ Source : Wikimedia Commons, « USMCA on the globe (North America centered) », Addicted04, CC0.}
\end{popfigure}

\begin{savaistu}[Le corridor Texas – Nord-Est mexicain : 1er passage frontalier commercial mondial]
En 2023, le pont \textbf{Laredo (TX) – Nuevo Laredo (Tamaulipas)} a traité \SI{~300}{Md\$} de marchandises (camions + rail), dépassant Detroit–Windsor (USA-Canada). Ce corridor « \emph{NAFTA Highway} » (I-35 / I-69) concentre l'industrie automobile (usines GM, Ford, Stellantis, Tesla au Texas ; clusters autopartistes à Monterrey, Saltillo, Ramos Arizpe). Il illustre la \textbf{proximité géographique} comme avantage comparatif décisif face à l'Asie (délais, coûts logistiques, résilience post-Covid).
\end{savaistu}

\courssub{Faiblesses structurelles et réponses stratégiques : « Friend-shoring » et réindustrialisation sélective}

\textbf{Les talons d'Achille de l'hégémon.}\par
Malgré sa puissance, le modèle US accumule des contradictions majeures :
\begin{enumerate}
\item \textbf{Inégalités record :} Gini \num{0,39} (disponibles) / \num{0,49} (primaires) ; 1\% des ménages détiennent \SI{30}{\%} de la richesse ; espérance de vie en baisse (opioïdes, armes, obésité).
\item \textbf{Dette publique explosive :} \SI{> 120}{\%} PIB (\SI{> 34 000}{Md\$}), service de la dette \SI{> 1 000}{Md\$}/an (dépasse budget Défense). Risque de « \emph{fiscal cliff} » (plafond dette, shutdowns).
\item \textbf{Polarisation politique / Dysfonctionnement institutionnel :} Blocage Congrès, remise en cause État de droit (Capitol 2021), incertitude réglementaire pour investisseurs.
\item \textbf{Infrastructures vieillissantes :} Note ASCE « C- » (ponts, routes, réseau électrique, ports, broadband rural). Plan Biden \SI{1 200}{Md\$} (IIJA, 2021) insuffisant.
\item \textbf{Déficit commercial structurel de biens :} \SI{~1 000}{Md\$} (importations biens de consommation, électronique, habillement, machinerie). Dépendance chaînes de valeur asiatiques (Chine, Taïwan pour puces).
\end{enumerate}

\textbf{La rupture stratégique : « Friend-shoring » / « Near-shoring » (Biden / Yellen).}\par
Face aux chocs (Covid, guerre Ukraine, tensions Taïwan), Washington opère un virage \textbf{industriel et géoéconomique} : relocaliser les chaînes critiques vers des pays \textbf{alliés, démocratiques, proches géographiquement}.
\begin{itemize}
\item \textbf{CHIPS and Science Act (2022) :} \SI{52}{Md\$} subventions + \SI{25}{\%} crédit d'impôt (Section 48D) pour usines semi-conducteurs US (TSMC Arizona, Intel Ohio, Samsung Texas). Objectif : \SI{20}{\%} production mondiale puces avancées < 10nm en 2030 (vs 0\% aujourd'hui).
\item \textbf{Inflation Reduction Act - IRA (2022) :} \SI{370}{Md\$} climat/énergie. Crédits d'impôt véhicules électriques \textbf{conditionnés} à assemblage en Amérique du Nord + batteries/minerais critiques extraits/transformés en US ou pays FTA (Accord de libre-échange). \rightarrow \textbf{Effet aimant pour Mexique (lithium, assemblage EV) et Canada (minerais critiques, hydroélectricité verte).}
\item \textbf{Alliances minérales :} \emph{Minerals Security Partnership} (MSP, 14 pays + UE) pour diversifier approvisionnement hors Chine (Graphite Mozambique, Lithium Australie/Chili, Cobalt RDC/Indonésie).
\end{itemize}

\begin{definition}[Friend-shoring (alliage) / Near-shoring (proximité)]
Stratégie de diversification des chaînes d'approvisionnement (\emph{supply chains}) consistant à privilégier des partenaires \textbf{politiquement alignés} (valeurs démocratiques, sécurité partagée) et/ou \textbf{géographiquement proches} (réduction délais, risques logistiques, empreinte carbone), au détriment de l'efficacité pure coûts (\emph{just-in-time} mondialisé). C'est la géoéconomie primant sur l'économie pure.
\end{definition}

\begin{methode}[Commenter un graphique d'évolution du PIB par tête (USA vs UE vs Chine)]
\textbf{Étape 1 : Lecture technique.} Identifier axes (années, \$ PPP ou nominal), échelle (log/linéaire), base 100. Repérer les chocs (2008, 2020).
\textbf{Étape 2 : Tendances longues.} USA : croissance trend \SI{~1,5-2}{\%}/an, résilience post-chocs. UE : convergence ralentie, écart se creuse (\SI{~30}{\%} PIB/tête inf. USA). Chine : rattrapage spectaculaire (x10 depuis 2000), mais ralentissement récent (vieillissement, dette immobilière).
\textbf{Étape 3 : Facteurs explicatifs.} USA : productivité horaire + heures travaillées/an (\num{1 800}h vs \num{1 500}h UE) + croissance démographique (immigration). UE : productivité proche, mais choix loisirs/protection sociale + démographie stagnante. Chine : investissement massif (\SI{>40}{\%} PIB), transfert rural-urbain, gains productivité manufacturière.
\textbf{Étape 4 : Conclusion synthétique.} L'écart USA/UE est structurel (modèle social, énergie, tech, défense). Le rattrapage Chine ralentit (piège revenu intermédiaire, démographie).
\end{methode}

\begin{exoresolu}[Analyse de la stratégie « Friend-shoring » à travers l'exemple de la filière batterie EV]
\textbf{Document :} Extrait de l'\emph{Inflation Reduction Act} (IRA) § 30D : Crédit d'impôt \SI{7 500}{\$} pour achat VE neuf. Condition A : Assemblage final en Amérique du Nord. Condition B (Batterie) : \SI{50}{\%} composants batterie fabriqués/assemblés en US/ALÉNA (montée à \SI{100}{\%} en 2029). Condition C (Minéraux critiques) : \SI{40}{\%} extraits/transformés en US ou pays FTA (montée à \SI{80}{\%} en 2027).
\textbf{Question :} En quoi l'IRA illustre-t-il le « friend-shoring » et redessine-t-il l'espace USMCA ?
\textbf{Solution détaillée :}
\begin{enumerate}
\item \textbf{Identification de la contrainte géoéconomique :} Les USA ne produisent presque pas de lithium, cobalt, graphite, nickel raffiné. La Chine contrôle \SI{60-80}{\%} du raffinage mondial. L'IRA force la création d'une chaîne de valeur \textbf{occidentale intégrée}.
\item \textbf{Rôle du Mexique (Near-shoring) :} Gisements lithium (Sonora, Zacatecas) + usines cathodes (ex: Ganfeng Lithium, sous pression pour JV locales) + assemblage packs (Tesla Gigafactory Texas, GM Ramos Arizpe). Main-d'œuvre qualifiée (\SI{~5}{\$}/h) + proximité logistique (rail/camion). \rightarrow \textbf{Mexique devient « Battery Belt » Sud.}
\item \textbf{Rôle du Canada (Friend-shoring ressources) :} Mines lithium (Québec, Manitoba), cobalt (Ontario), nickel (Voisey's Bay), graphite. Raffinage hydroélectrique bas carbone (Québec). Usines cathodes/anodes (GM-POSCO Bécancour, Northvolt Montréal, Volkswagen St. Thomas). Accord US-Canada « Critical Minerals » (2024). \rightarrow \textbf{Canada = « Mine to Battery » vert.}
\item \textbf{Rôle des USA (Commandement/Intégration) :} Cellules (Panasonic/Toyota Kansas, LG Energy Solution Arizona, Ford/SK Kentucky/Tennessee), assemblage final, R\&D (DOE labs), finance (IRA crédits d'impôt \SI{45X}{/kWh} production cellules/modules).
\item \textbf{Synthèse :} L'IRA transforme l'USMCA en \textbf{bloc autonome « Mine-to-EV »}, réduisant dépendance Chine, créant emplois industriels US (Swing States), valorisant ressources alliés. Risque : contentieux OMC (subventions discriminatoires), coûts élevés transférés au consommateur, tensions eau/environnement locales (Mexique, Nevada).
\end{enumerate}
\end{exoresolu}

\courssub{Le Cameroun face à la puissance américaine : AGOA, investissements et diaspora}

\textbf{Un partenaire commercial asymétrique mais structurant.}\par
Les USA sont la \textbf{2e destination des exportations camerounaises} (\SI{~10-12}{\%} du total, après UE/Chine), principalement via le régime préférentiel \textbf{AGOA} (\emph{African Growth and Opportunity Act}, renouvelé jusqu'en 2025).
\begin{itemize}
\item \textbf{Exportations Cameroun $\to$ USA :} Bois sciés/transformés (\SI{~40}{\%} export AGOA), Pétrole brut (\SI{~30}{\%}, volatile), Cacao/café, Coton, Aluminium (Alucam). \textbf{Valeur :} \SI{~200-300}{M\$}/an (hors pétrole : \SI{~100}{M\$}).
\item \textbf{Importations USA $\to$ Cameroun :} Machines, véhicules, produits pharmaceutiques, céréales (blé, riz), équipements pétroliers. \textbf{Valeur :} \SI{~300-400}{M\$}/an. \textbf{Déficit commercial Cameroun.}
\end{itemize}

\textbf{Investissements directs (IDE) et coopération technique.}\par
\begin{itemize}
\item \textbf{Énergie (Historique) :} \textbf{Perenco} (Franco-Britannique, opérateur majeur SNH), \textbf{Noble Energy} (devenu \textbf{Chevron}, US) : exploitation champs offshore (Bolongo, Sanaga Sud). Nouveaux acteurs : \textbf{Bowleven}, \textbf{New Age} (région Rio del Rey).
\item \textbf{Technologie / Numérique :} Programmes \textbf{USAID} / \textbf{Power Africa} (électricité rurale, mini-réseaux solaires), \textbf{Google} / \textbf{Microsoft} (hubs développeurs, câbles sous-marins 2Africa / Equiano atterrissant à Kribi/Limbe).
\item \textbf{Santé / Capital humain :} \textbf{PEPFAR} (\SI{> 1}{Md\$} investis depuis 2004, lutte VIH/Sida), \textbf{CDC Atlanta} (laboratoire référence Yaoundé, ripostes Ebola/Choléra/Covid), \textbf{USAID} (éducation, gouvernance, humanitaire Nord/Extrême-Nord).
\item \textbf{Diaspora :} \num{~60-80 000} Camerounais aux USA (MD, TX, VA, GA). Transferts \textbf{formels} estimés à une part significative des \SI{~300-400}{M\$} annuels totaux (sous-estimés), investissements immobiliers/PME au pays, transfert compétences (médecins, ingénieurs, entrepreneurs).
\end{itemize}

\begin{aretenir}[Bilan : L'Amérique du Nord, foyer moteur à géométrie variable]
\begin{itemize}
\item \textbf{USA :} Moteur principal (Innovation, Finance, Défense, Dollar, Énergie). Modèle libéral flexible mais fragilisé (inégalités, dette, polarisation, infrastructure).
\item \textbf{USMCA :} Bloc économique intégré \textbf{par les chaînes de valeur} (Auto, Énergie, Agro, Tech). Complémentarité : Cerveau/Capital (US) + Ressources/Énergie verte (CA) + Main-d'œuvre/Assemblage/Proximité (MX).
\item \textbf{Stratégie 2020s :} \textbf{Friend-shoring / Near-shoring} (CHIPS Act, IRA) $\rightarrow$ Réindustrialisation sélective critique (Puces, Batteries, Pharma, Minéraux) \textbf{au sein même de l'USMCA} et avec alliés (Japon, Corée, UE, Australie).
\item \textbf{Cameroun :} Partenaire secondaire pour US, mais vital pour diversification (AGOA bois/agro), énergie (offshore), santé (PEPFAR), diaspora. Enjeu : passer de rente (pétrole/bois brut) à transformation locale (AGOA exige \emph{value addition}) et attirer IDE tech/vert via diaspora.
\end{itemize}
\end{aretenir}

\coursec{L'Asie de l'Est : l'atelier du monde et l'innovation (Chine, Japon, Corée du Sud)}

Après avoir analysé la puissance hégémonique nord-américaine et son espace intégré \emph{USMCA}, tournons notre regard vers l'ouest du Pacifique. C'est là que se joue aujourd'hui la plus forte dynamique de croissance manufacturière et technologique. Si les États-Unis restent le centre de commande financier et numérique, l'Asie de l'Est est devenue l'\textbf{atelier du monde} et un foyer d'innovation majeur, rivalisant désormais avec la \cle{Triade} historique sur le terrain de la haute technologie.

\courssub{La Chine : l'émergence du premier manufacturier mondial par le capitalisme d'État}

Commençons par une observation concrète. Regardez l'étiquette de votre smartphone, de vos chaussures ou des machines-outils du lycée technique de Douala : la mention \og Made in China \fg{} est omniprésente. Ce n'est pas un hasard.

\begin{definition}[Capitalisme d'État chinois]
Modèle économique où l'État, via le Parti communiste chinois (PCC), pilote l'allocation des ressources (crédits directs des banques publiques, entreprises d'État \emph{SOE}, plans quinquennaux) tout en intégrant les mécanismes de marché (concurrence, investissements privés, IDE, exportations) pour assurer la croissance, la stabilité sociale et la montée en gamme technologique.
\end{definition}

Depuis les réformes de \textbf{Deng Xiaoping (1978)} et l'adhésion à l'\textbf{OMC (2001)}, la Chine a construit sa puissance sur :
\begin{itemize}
    \item Une \textbf{main-d'œuvre abondante et bon marché} (dividende démographique) attirant les IDE pour l'assemblage (phase \og usine du monde \fg).
    \item Un \textbf{investissement massif en infrastructures} (routes, ports, TGV, énergie) financé par l'épargne nationale (\textasciitilde 44\% du PIB) et la dette publique locale.
    \item Une \textbf{stratégie mercantiliste} : excédents commerciaux réinvestis en réserves de change (\textasciitilde \num{3200} Md \$ en 2023) et en obligations du Trésor US.
\end{itemize}

\paragraph{La montée en gamme : du « Made in China » au « Created in China ».}
Le \textbf{XIV\textsuperscript{e} Plan quinquennal (2021-2025)} et la stratégie \textbf{\og Industries du futur \fg} (ex-\og Made in China 2025 \fg) visent l'autonomie dans 10 secteurs critiques : semi-conducteurs, IA, biotech, aéronautique (C919), véhicules électriques (BYD, NIO), énergies nouvelles.
\begin{exemplebox}[Chiffres clés Chine 2023]
\begin{itemize}
    \item \textbf{PIB :} \textasciitilde 17\,800 Md \$ (2\textsuperscript{e} nominal, 1\textsuperscript{er} PPA \textasciitilde 33\,000 Md \$).
    \item \textbf{Exportations marchandises :} 1\textsuperscript{er} mondial (\textasciitilde 3\,380 Md \$), excédent commercial \textasciitilde 820 Md \$.
    \item \textbf{Valeur ajoutée manufacturière :} 28\% du total mondial (1\textsuperscript{er} depuis 2010).
    \item \textbf{R\&D :} 2,6\% du PIB (\textasciitilde 660 Md \$), 1\textsuperscript{er} dépôt de brevets PCT (OMPI) avec \textasciitilde 1,6 million (2023).
\end{itemize}
\end{exemplebox}

\begin{definition}[BRI — Belt and Road Initiative / Nouvelles Routes de la Soie]
Lancée par Xi Jinping en 2013. Stratégie d'infrastructures (ports, rails, autoroutes, pipelines, fibre optique, zones économiques spéciales) reliant la Chine à l'Eurasie, l'Afrique et l'Amérique latine. Objectif officiel : \og connectivité gagnant-gagnant \fg. Critique occidentale : \og \emph{debt-trap diplomacy} \fg{} (piège de la dette) pour acquérir des actifs stratégiques (ex. port du Pirée en Grèce, port de Hambantota au Sri Lanka).
\end{definition}

\begin{popfigure}
\begin{tikzpicture}
\begin{axis}[
    width=12cm, height=7cm,
    stack plots=y,
    area style,
    legend style={at={(0.02,0.98)}, anchor=north west, fill=white, draw=popline},
    ylabel={Exportations high-tech (Md \$)}, 
    symbolic x coords={2000,2005,2010,2015,2020,2023},
    xtick=data,
    xlabel={Année},
    ymin=0, ymax=1400,
    tick label style={font=\small},
    label style={font=\small},
    legend style={font=\small}
]
\addplot[fill=poppink!30, draw=poppink] coordinates {(2000,50) (2005,150) (2010,400) (2015,600) (2020,800) (2023,900)}; \addlegendentry{Chine}
\addplot[fill=popblue!30, draw=popblue] coordinates {(2000,120) (2005,130) (2010,150) (2015,160) (2020,170) (2023,180)}; \addlegendentry{Japon}
\addplot[fill=popgreen!30, draw=popgreen] coordinates {(2000,30) (2005,60) (2010,120) (2015,180) (2020,220) (2023,250)}; \addlegendentry{Corée du Sud}
\end{axis}
\end{tikzpicture}
\legende{Évolution des exportations de produits de haute technologie (en milliards de dollars courants). Source : OCDE / Banque mondiale, données retraitées. La Chine rattrape et dépasse largement le Japon et la Corée du Sud depuis 2010.}
\end{popfigure}

\courssub{Le Japon : la puissance mature, créancière et robotisée}

Le Japon, 4\textsuperscript{e} PIB mondial (\textasciitilde 4\,200 Md \$ en 2023), présente un profil radicalement différent : pays vieilli (médiane 49 ans), dette publique \textgreater 260\% du PIB, mais \textbf{1\textsuperscript{er} créancier net mondial} (\textasciitilde 3\,400 Md \$ d'actifs nets extérieurs).

\begin{definition}[Keiretsu]
Groupes industriels et financiers japonais (ex. Mitsubishi, Mitsui, Sumitomo) caractérisés par des participations croisées, une banque centrale du groupe et des relations stables \emph{keiretsu} (horizontal) ou \emph{vertical} (sous-traitance). Ils assurent la stabilité de l'emploi et le financement patient de l'innovation (long terme), au détriment parfois de la rentabilité à court terme.
\end{definition}

Le modèle japonais repose sur :
\begin{itemize}
    \item \textbf{La qualité et la productivité} (\emph{Kaizen}, \emph{Just-in-Time} Toyota).
    \item \textbf{L'innovation incrémentale} : robotique industrielle (Fanuc, Yaskawa \textasciitilde 50\% part mondiale), composants électroniques passifs/actifs (Murata, TDK), matériaux avancés (carbone, céramiques), automobile (Toyota \#1 mondial).
    \item \textbf{Une spécialisation en amont des chaînes de valeur} : le Japon vend les machines-outils, les précurseurs chimiques et les équipements de lithographie (Canon, Nikon) aux usines d'assemblage chinoises ou coréennes.
\end{itemize}

\begin{attention}[Piège : « Le Japon est en déclin »]
Ne confondez pas \emph{stagnation démographique/PIB par tête} et \emph{perte de puissance technologique}. Le Japon reste \textbf{indispensable} aux chaînes de valeur mondiales (semi-conducteurs, automobile, chimie fine). Son excédent courant (\textasciitilde 150 Md \$) finance son vieillissement. Au Bac, citez : \og Puissance mature, relais technologique de la Triade, premier créancier net \fg.
\end{attention}

\courssub{La Corée du Sud : le « Miracle du Han » et l'empire des Chaebols}

De la guerre (1950-53) au 12\textsuperscript{e} PIB mondial (\textasciitilde 1\,700 Md \$), la Corée du Sud a compressé en 50 ans ce que l'Europe a mis 150 ans à faire.

\begin{definition}[Chaebol]
Conglomérats familiaux (Samsung, Hyundai Motor, LG, SK) diversifiés (électronique, auto, navale, chimie, finance, assurance, construction), pilotés par le fondateur/héritier, soutenus historiquement par l'État (crédits dirigés, protection marché intérieur). Représentent \textasciitilde 60\% du PIB et 70\% des exportations.
\end{definition}

La transition réussie : \textbf{Bas salaires (années 70-80) $\rightarrow$ R\&D intensive (aujourd'hui)}.
\begin{itemize}
    \item \textbf{Effort R\&D :} \textbf{4,9\% du PIB (2022)}, \textbf{record mondial} (Israël \textasciitilde 5,6\% mais PIB plus petit). Samsung Electronics seul investit \textasciitilde 20 Md \$/an.
    \item \textbf{Spécialisations dominantes :} 1\textsuperscript{er} exportateur mondial puces mémoires (Samsung, SK Hynix \textasciitilde 60-70\% part), navires (HD Hyundai, Hanwha Ocean), batteries LG/Samsung/SK, écrans OLED.
    \item \textbf{Modèle \og Fast Follower \fg $\rightarrow$ \og First Mover \fg} : passage de la copie sous licence à la définition des standards (5G, 6G, IA embarquée).
\end{itemize}

\begin{exoresolu}[Analyse de la structure d'exportation coréenne (Données 2023 approx.)]
\textbf{Énoncé :} À partir du tableau ci-dessous, calculez la part des 3 premiers postes et commentez la spécialisation du pays.

\begin{center}
\begin{tabularx}{\linewidth}{|l|X|X|}\hline
\textbf{Produit} & \textbf{Valeur (Md \$)} & \textbf{Part (\%)} \\ \hline
Semi-conducteurs (puces mémoires) & 110 & 19,5\% \\ \hline
Pétrole raffiné & 55 & 9,8\% \\ \hline
Automobiles & 52 & 9,2\% \\ \hline
Navires / Plateformes offshore & 28 & 5,0\% \\ \hline
Écrans plats / OLED & 22 & 3,9\% \\ \hline
Autres & 300 & 52,6\% \\ \hline
\textbf{Total} & \textbf{567} & \textbf{100\%} \\ \hline
\end{tabularx}
\end{center}

\textbf{Solution détaillée :}
\begin{enumerate}
    \item \textbf{Calcul part Top 3 :} Puces (19,5\%) + Pétrole raffiné (9,8\%) + Auto (9,2\%) = \textbf{38,5\%}. Presque 40\% des recettes d'exportation reposent sur 3 secteurs.
    \item \textbf{Interprétation :} Spécialisation \textbf{très concentrée} (risque de \emph{Dutch disease} sectoriel). Dominance \textbf{haute technologie} (puces, auto, écrans = 32,6\%). Le pétrole raffiné montre l'intégration aval (pétrochimie SK, S-Oil). La construction navale (5\%) reste un atout stratégique (porte-conteneurs, GNL).
    \item \textbf{Vocabulaire Bac :} \og Spécialisation comparative révélée \fg (Balassa) forte dans l'électronique et le transport maritime.
\end{enumerate}
\end{exoresolu}

\courssub{Intégration régionale : le RCEP et les chaînes de valeur « Usine Asie »}

L'Asie de l'Est n'est pas une juxtaposition d'États rivaux, mais un \textbf{système productif intégré} par les IDE japonais/coréens/taïwanais et le marché chinois.

\begin{propriete}[RCEP — Regional Comprehensive Economic Partnership]
Entré en vigueur \textbf{janvier 2022}. Plus grand ALE (Accord de Libre-Échange) mondial : \textbf{15 pays} (ASEAN 10 + Chine + Japon + Corée du Sud + Australie + Nouvelle-Zélande). \textasciitilde 30\% PIB mondial, 30\% population, 28\% commerce mondial. \textbf{Innovation :} Règles d'origine \emph{cumulatives} (un produit assemblé au Vietnam avec composants japonais/chinois/koréens est \og originaire RCEP \fg). Renforce l'usine Asie.
\end{propriete}

\begin{methode}[Réaliser un croquis des chaînes de valeur électroniques (Spécialité Bac)]
\textbf{Consigne :} Représenter la fragmentation internationale de la production d'un smartphone.
\begin{enumerate}
    \item \textbf{Identifier les étapes (amont $\rightarrow$ aval) :}
        \begin{itemize}
            \item \textbf{R\&D / Design :} USA (Apple, Qualcomm), Japon, Corée, Taïwan.
            \item \textbf{Wafers / Matériaux :} Japon (Shin-Etsu, SUMCO \textasciitilde 50\% wafers Si), USA, Allemagne (Siltronic), Corée.
            \item \textbf{Équipements lithographie :} \textbf{ASML (Pays-Bas)} \emph{unique fournisseur EUV}, tech US/Japon/EU.
            \item \textbf{Fonderies (Front-end) :} \textbf{TSMC (Taïwan) \textasciitilde 60\% foundry, 90\% <10nm}, Samsung (Corée), Intel (USA), SMIC (Chine).
            \item \textbf{Assemblage / Test (Back-end) :} Chine (Foxconn, Luxshare), Vietnam, Malaisie, Inde (montée).
            \item \textbf{Marque / Distribution :} USA, Corée, Chine, Japon.
        \end{itemize}
    \item \textbf{Flèches proportionnelles :} Épaisseur = valeur ajoutée / volume. Flèche massive \og Composants $\rightarrow$ Assemblage Chine/Vietnam \fg. Flèche fine \og Assemblage $\rightarrow$ Marque \fg (faible VA en assemblage \textasciitilde 3-5\% prix iPhone).
    \item \textbf{Légende structurée :} 3 colonnes : \og Pays / Région \fg, \og Spécialisation (Étape) \fg, \og Part de marché / Acteurs clés \fg.
    \item \textbf{Titre + Échelle + Nord + Sources.}
\end{enumerate}
\end{methode}

\begin{popfigure}
\begin{tikzpicture}[scale=0.85, every node/.style={font=\small}]
% Nodes: Countries/Regions
\node[draw, rounded corners, fill=popblue!10, minimum width=2.5cm, minimum height=1.2cm, align=center] (US) at (0,4){USA\\Design, IP, EDA};
\node[draw, rounded corners, fill=poporange!10, minimum width=2.5cm, minimum height=1.2cm, align=center] (JP) at (0,2){JAPON\\Matériaux, Équipements, Composants passifs};
\node[draw, rounded corners, fill=poppurple!10, minimum width=2.5cm, minimum height=1.2cm, align=center] (KR) at (0,0){CORÉE DU SUD\\Puces mémoire, Écrans, Batteries};
\node[draw, rounded corners, fill=poppink!10, minimum width=2.5cm, minimum height=1.2cm, align=center] (TW) at (0,-2){TAÏWAN\\Fonderies (TSMC)};
\node[draw, rounded corners, fill=popgreen!10, minimum width=3cm, minimum height=1.5cm, align=center] (CN) at (5,1){CHINE / VIETNAM / INDE\\Assemblage Final, Test, Emballage};
\node[draw, rounded corners, fill=popgold!10, minimum width=2.5cm, minimum height=1.2cm, align=center] (MKT) at (9,1){MARCHÉS MONDIAUX\\Marques (Apple, Samsung, Xiaomi...)};

% Arrows
\draw[-{Stealth[length=3mm]}, thick, popblue] (US.east) -- node[above, sloped, pos=0.5] {Design, IP} (CN.west);
\draw[-{Stealth[length=3mm]}, thick, poporange] (JP.east) -- node[above, sloped, pos=0.5] {Wafers, Chimie, Machines} (CN.west);
\draw[-{Stealth[length=3mm]}, thick, poppurple] (KR.east) -- node[above, sloped, pos=0.5] {DRAM/NAND, OLED, Batteries} (CN.west);
\draw[-{Stealth[length=3mm]}, thick, poppink, line width=1.5pt] (TW.east) -- node[above, sloped, pos=0.5] {\textbf{Logique <10nm (90\%)}} (CN.west);
\draw[-{Stealth[length=3mm]}, thick, popgreen] (CN.east) -- node[above] {Produits finis} (MKT.west);

% Ports
\node[anchor=west, font=\tiny, text=popmuted] at (5.5, -1.5) {Ports: Shanghai, Shenzhen, Busan, Yokohama, Singapour (Hub)};
\node[anchor=west, font=\tiny, text=popmuted] at (5.5, -1.8) {BRI: Routes ferroviaires (Chine-Europe), Maritimes};
\end{tikzpicture}
\legende{Schéma simplifié de la chaîne de valeur électronique en Asie de l'Est (Flux amont $\rightarrow$ assemblage $\rightarrow$ marchés). Épaisseur des flèches $\approx$ criticité / valeur technologique. Taïwan (TSMC) et ASML (Pays-Bas, non montré ici mais fournisseur du Japon/EU) sont les goulets d'étranglement (choke points).}
\end{popfigure}

\courssub{Tensions géopolitiques : Mer de Chine, Taïwan et « dérisquage »}

Cette interdépendance est le théâtre de la rivalité sino-américaine.

\begin{itemize}
    \item \textbf{Mer de Chine méridionale :} Revendications chinoises (ligne en neuf traits) vs Philippines, Vietnam, Malaisie, Brunei (décision CPI 2016 ignorée par Pékin). Militarisation d'îlots artificiels. Enjeu : \textbf{30\% commerce maritime mondial}, ressources halieutiques, hydrocarbures.
    \item \textbf{Taïwan (République de Chine) :} \textbf{TSMC = 90\% puces <10 nm mondiales}. « Bouclier de silicium » : une invasion ou un blocus provoquerait un effondrement économique mondial (automobile, défense, IA, smartphones). Stratégie US \emph{CHIPS Act} (2022, 52 Md \$) + \emph{CHIPS Act} UE/Japon/Corée pour relocaliser la fonderie (usines TSMC Arizona, Kumamoto Japon, Dresde Allemagne).
    \item \textbf{Dérisquage (\emph{De-risking}) occidental :} Ne plus dépendre \emph{uniquement} de la Chine. Diversification \textbf{\og China+1 \fg} : Vietnam (Samsung, Foxconn), Inde (Apple, Tata), Mexique (near-shoring US). \textbf{Résultat :} IDE en Chine en baisse (2023: premier flux net négatif trimestriel depuis 1998), mais Chine reste \textbf{centre d'assemblage final dominant} (écosystème complet, logistique, main-d'œuvre qualifiée).
\end{itemize}

\begin{attention}[Erreur fréquente : « La Chine fabrique l'iPhone »]
\textbf{Faux.} La Chine (et maintenant l'Inde/Vietnam) fait l'\textbf{assemblage final} (Foxconn, Luxshare). La \textbf{valeur ajoutée} (R\&D, design, puces, OS, écrans, brevets) reste \textbf{majoritairement américaine, taïwanaise, japonaise, coréenne, européenne}. Au Bac, précisez : \og La Chine capture la faible VA de l'assemblage (\textasciitilde 3-5\%), mais monte en gamme (conception Huawei, Xiaomi, BYD) \fg.
\end{attention}

\courssub{L'Asie de l'Est et le Cameroun : partenariats asymétriques mais structurants}

Le Cameroun tourne résolument vers l'Est depuis les années 2000.

\begin{exemplebox}[Coopération Cameroun — Asie de l'Est (Données récentes)]
\begin{center}
\begin{tabularx}{\linewidth}{|l|X|X|}\hline
\textbf{Pays} & \textbf{Principaux flux / Projets} & \textbf{Nature du partenariat} \\ \hline
\textbf{CHINE (1\textsuperscript{er} partenaire commercial)} & \textbf{Import :} Machines, textiles, électronique, véhicules, BTP. \textbf{Export :} Bois (grumes, sciés), Coton, Minerais (Fer Mbalam/Nabeba - en projet, Bauxite), Pétrole brut. \textbf{Projets BRI :} Port en eau profonde \textbf{Kribi} (CHEC), Barrages \textbf{Mékin, Lom Pangar, Memve'ele}, Routes (Yaoundé-Nsimalen, Douala-Yaoundé), Stade Olembé, Immeubles ministériels. & \textbf{Infrastructures contre ressources / Marché.} Financement \emph{Exim Bank / China Dev Bank}. Dette \textasciitilde 40\% de la dette extérieure totale (part prépondérante de la dette bilatérale). \\ \hline
\textbf{CORÉE DU SUD} & \textbf{Agricole :} Projet rizicole \textbf{SEMRY II} (vallée Logone), Centre de formation agricole. \textbf{Santé :} Hôpital de référence de Douala (KORCAM), équipement. \textbf{Éducation :} Bourses KOICA, Centre Corée-Cameroun. & \textbf{Coopération technique / ODA (Aide publique au développement).} Modèle \og Saemaul Undong \fg (nouveau village). \\ \hline
\textbf{JAPON} & \textbf{Pêche :} Accords d'accès thoniers (navires japonais en ZEE). \textbf{Infrastructures :} Pont Wouri (2\textsuperscript{e} pont), Route Bamenda-Mamfé (JICA). \textbf{Formation :} Centres de formation professionnelle (CFP), Bourses MEXT. \textbf{TICAD :} Cadre multilatéral (Tokyo International Conference on African Development). & \textbf{Aide liée / Investissement qualité / Transfert de compétences.} Moins de prêts massifs, plus de dons/expertise. \\ \hline
\end{tabularx}
\end{center}
\end{exemplebox}

\begin{savaistu}[Le « Goulet d'étranglement » des semi-conducteurs : TSMC et ASML]
Le monde numérique repose sur deux entreprises \textbf{quasi-monopolistiques} :
\begin{enumerate}
    \item \textbf{TSMC (Taiwan Semiconductor Manufacturing Co., Taïwan) :} \textasciitilde 60\% fonderie mondiale, \textbf{90\% puces <10 nm} (IA, smartphones haut de gamme, défense). Aucune alternative crédible à court terme.
    \item \textbf{ASML (Pays-Bas) :} \textbf{Seul fabricant mondial} de machines de lithographie \textbf{EUV (Extreme Ultraviolet)} indispensables pour graver <7 nm. Technologie issue de 40 ans de R\&D US (Intel), EU (IMEC), Japon (optique Zeiss/Carl Zeiss SMT), composants critiques japonais.
\end{enumerate}
\textbf{Conséquence géopolitique :} Taïwan devient la \og montagne la plus précieuse du monde \fg. Les US interdisent à ASML d'exporter l'EUV vers la Chine (régime Wassenaar). La Chine (SMIC) a réussi du 7 nm (DUV multi-patterning) mais avec rendements faibles et coûts élevés. La course au \textbf{3 nm / 2 nm (Gate-All-Around)} se joue \textbf{aujourd'hui} entre TSMC, Samsung, Intel.
\end{savaistu}

\begin{exercice}[Entraînement Bac — Sujet type : « L'Asie de l'Est, foyer moteur de la mondialisation »]
\textbf{Document 1 :} Carte des flux de composants électroniques en Asie de l'Est (type croquis ci-dessus).
\textbf{Document 2 :} Tableau : Parts de marché mondiales 2023 (Puces logique TSMC/Samsung/Intel ; Puces mémoire Samsung/SK Hynix/Micron ; Navires Corée/Chine/Japon ; VE BYD/Tesla/VW).
\textbf{Document 3 :} Extrait discours Xi Jinping (2023) sur « Nouvelle qualité de forces productives » + Extrait \emph{CHIPS Act} US (2022).

\textbf{Questions :}
\begin{enumerate}
    \item À l'aide du Document 1, réalisez un croquis légendé des chaînes de valeur électroniques en Asie de l'Est. (6 pts)
    \item En vous appuyant sur les Documents 2 et 3, montrez que l'Asie de l'Est est à la fois un espace \textbf{intégré par les chaînes de valeur} et un espace \textbf{fragmenté par les rivalités géopolitiques}. (14 pts)
\end{enumerate}
\end{exercice}

\begin{corrige}[Exercice n°13-4]
\textbf{Q1. Croquis (Attendus) :}
\begin{itemize}
    \item \textbf{Fond de carte :} Contours Chine, Japon, Corées, Taïwan, Mer de Chine, Océan Pacifique.
    \item \textbf{Symboles :} Carrés (R\&D/Design : USA/Japon/Corée/Taïwan), Losanges (Fonderies : TSMC Taïwan, Samsung Corée), Cercles (Assemblage : Chine côtière, Vietnam, Inde), Flèches proportionnelles (Composants $\rightarrow$ Assemblage ; Produits finis $\rightarrow$ Monde).
    \item \textbf{Légende organisée :} 1. Amont (R\&D, Matériaux, Équipements) — 2. Front-end (Fonderies) — 3. Back-end (Assemblage/Test) — 4. Ports/Hubs (Shanghai, Shenzhen, Busan, Singapour) — 5. Tensions (Mer de Chine, Détroit Taïwan).
    \item \textbf{Titre :} \og Asie de l'Est : fragmentation spatiale de la chaîne de valeur des semi-conducteurs et de l'électronique (2023) \fg.
\end{itemize}
\textbf{Q2. Argumentation structurée (Intégration vs Fragmentation) :}
\begin{itemize}
    \item \textbf{Intégration (RCEP, IDE, CVM) :} Flux massifs composants JP/KR/TW $\rightarrow$ Assemblage CN/VN (Doc 1). Spécialisations complémentaires (Doc 2 : JP matériaux/équipements, KR mémoire/écrans, TW fonderie, CN assemblage/marché). RCEP règles origine cumulatives facilitent.
    \item \textbf{Fragmentation (Géopolitique, Sécurité) :} Rivalité US/Chine $\rightarrow$ Sanctions US (Entity List, contrôles export EUV/IA chips vers CN). \emph{CHIPS Act} US/Japon/UE/KR = relocalisation/subvention fonderies hors Chine (Arizona, Kumamoto, Dresde, Yongin). Taïwan = enjeu souveraineté (Doc 3). \og Friend-shoring \fg / \og De-risking \fg cassent l'optimalité économique pure.
    \item \textbf{Synthèse :} \og Interdépendance contrainte \fg : on ne peut pas se passer de l'usine Asie (coûts, écosystème), mais on veut en réduire la vulnérabilité stratégique. L'Asie de l'Est reste le \textbf{foyer moteur n°1 de la fabrication mondiale}, mais sa gouvernance est contestée.
\end{itemize}
\end{corrige}

\begin{aretenir}[Synthèse : L'Asie de l'Est, foyer moteur polycentrique]
\begin{itemize}
    \item \textbf{Chine :} Moteur principal (taille, croissance, BRI), \textbf{atelier final}, montée en gamme forcée (autonomie puces, VE, IA). Modèle : Capitalisme d'État planifié.
    \item \textbf{Japon :} Moteur technologique amont (matériaux, équipements, robotique), créancier net, modèle \emph{Keiretsu} / qualité. Défi démographique.
    \item \textbf{Corée du Sud :} Moteur d'innovation privée (\emph{Chaebols}), record R\&D/PIB, spécialisation puces/écrans/batteries/navires. Vulnérabilité géopolitique (Corée du Nord, Chine, US).
    \item \textbf{Intégration :} \textbf{RCEP} + \textbf{Chaînes de valeur verticales} (Amont JP/KR/TW $\rightarrow$ Aval CN/ASEAN) = \og Usine Asie \fg.
    \item \textbf{Tensions :} Taïwan (puces), Mer de Chine (routes), Dérisquage occidental (China+1).
    \item \textbf{Cameroun :} Chine = partenaire n°1 (infra/ressources), Corée/Japon = partenaires ciblés (agri/santé/formation/pêche).
\end{itemize}
\end{aretenir}

\coursec{La Russie : foyer énergétique, militaire et pivot eurasiatique}

Après avoir analysé la dynamique de l'Asie de l'Est, atelier du monde et foyer d'innovation technologique, nous tournons notre regard vers le nord de l'Eurasie. La Russie occupe une place à part dans la géographie de la mondialisation : elle n'est pas une puissance économique de premier rang par la taille de son PIB (de l'ordre de 2 000 milliards de dollars en 2023 selon le FMI, soit environ 2 \% du PIB mondial, comparable à celui de l'Italie ou du Canada), mais elle demeure une \cle{superpuissance stratégique} par son arsenal nucléaire, son siège de membre permanent au Conseil de sécurité de l'ONU, sa maîtrise de l'espace et du cyberespace, et surtout son statut de \cle{foyer énergétique} planétaire. Comprendre la Russie, c'est comprendre comment un pays peut transformer une rente naturelle en levier de puissance géopolitique, tout en gérant des fragilités structurelles profondes.

\courssub{Une spécificité unique : puissance économique moyenne, superpuissance stratégique}

\emph{Observation introductive.} Regardons la facture de gaz naturel payée par une usine chimique allemande en 2021, puis celle de la même usine en 2023 : le fournisseur a changé, le prix a fortement augmenté, mais le gaz continue de chauffer l'Europe par d'autres voies (GNL américain, gaz norvégien, gaz azerbaïdjanais). Derrière cette facture, il y a la Russie : un pays dont le PIB nominal (environ 2 000 Md \$ en 2023, FMI) ne représente qu'environ 2 \% du PIB mondial, mais qui détient les \textbf{premières réserves mondiales de gaz naturel}, figure parmi les tout premiers producteurs mondiaux de pétrole, et domine les exportations mondiales de blé, tout en étant un acteur majeur du palladium, du nickel et de l'uranium enrichi.

\begin{definition}[Pivot vers l'Asie (\emph{Turn to the East})]
Réorientation stratégique des flux commerciaux, financiers et diplomatiques de la Russie vers la Chine, l'Inde, l'Iran, la Turquie et les BRICS, à la suite des sanctions occidentales adoptées après 2022. Ce pivot vise à réduire la dépendance aux marchés et aux technologies de l'UE et des États-Unis, en réorientant les exportations d'hydrocarbures (gazoduc \emph{Force de Sibérie}), en utilisant les monnaies nationales (rouble, yuan, roupie) dans le commerce bilatéral, et en approfondissant la coopération militaire et technologique avec Pékin.
\end{definition}

La Russie est une \cle{économie de rente} : les hydrocarbures représentent environ 30 à 40 \% des recettes du budget fédéral, près de 60 \% des exportations totales et 15 à 20 \% du PIB (ordres de grandeur, sources FMI et Banque mondiale). Cette structure crée une vulnérabilité aux chocs de prix, mais offre aussi une \cle{rente de situation} géopolitique : qui a besoin de l'énergie russe ? L'Europe (historiquement), la Chine (désormais), l'Inde (de plus en plus). L'État fort (pouvoir présidentiel centralisé) capte cette rente via des champions nationaux (Gazprom, Rosneft, Rosatom, Rostec) souvent dirigés par des \cle{oligarques} loyaux au Kremlin. Depuis 2014 (annexion de la Crimée) et surtout 2022 (offensive en Ukraine), des sanctions occidentales sans précédent (gel d'une partie des réserves de change, exclusion de nombreuses banques du réseau SWIFT, embargo technologique, plafonnement du prix du pétrole russe) ont tenté d'étrangler ce modèle. La réponse russe a été le \emph{Pivot vers l'Asie} et le renforcement de la coopération au sein de l'\cle{OPEP+} avec l'Arabie saoudite pour gérer l'offre mondiale de pétrole.

\begin{definition}[OPEP+]
Alliance élargie réunissant les pays membres de l'OPEP (Arabie saoudite, Irak, Émirats arabes unis, etc.) et plusieurs producteurs non membres de l'OPEP, dont la Russie (principal partenaire), le Kazakhstan ou encore le Mexique. Objectif : coordonner les niveaux de production pour stabiliser ou soutenir les cours du pétrole. La Russie y joue un rôle clé de producteur d'ajustement non-OPEP, acceptant des réductions de production en échange de prix élevés qui financent son budget.
\end{definition}

\courssub{Structure des exportations : la dépendance aux matières premières}

Pour visualiser cette économie de rente, analysons la structure des exportations russes en 2023 (ordres de grandeur établis à partir des données du Service fédéral des douanes de Russie et de l'UN Comtrade, arrondis).

\begin{popfigure}
\begin{tikzpicture}
\begin{axis}[
    width=\linewidth,
    height=7cm,
    ybar,
    ylabel={Exportations 2023 (Md \$, ordres de grandeur)},
    symbolic x coords={Pétrole/Gaz, Métaux/Minéraux, Céréales, Armement, Nucléaire civil, Autres},
    xtick=data,
    xticklabel style={rotate=30, anchor=east, font=\small},
    ymin=0,
    ymax=400,
    bar width=12pt,
    fill=poporange!70,
    draw=popdark,
    nodes near coords,
    nodes near coords align={vertical},
    every node near coord/.append style={font=\tiny, rotate=90, anchor=west, yshift=2pt}
]
\addplot coordinates {
    (Pétrole/Gaz, 280)
    (Métaux/Minéraux, 50)
    (Céréales, 25)
    (Armement, 15)
    (Nucléaire civil, 10)
    (Autres, 45)
};
\end{axis}
\end{tikzpicture}
\legende{Figure 1 : Structure estimée des exportations russes en 2023 (ordres de grandeur). Source : Douanes russes, UN Comtrade, estimations. La domination des hydrocarbures (pétrole brut, produits raffinés, gaz naturel, GNL) est écrasante : environ 60 \% du total. Les métaux (aluminium, nickel, palladium, or) et les céréales (la Russie est le premier exportateur mondial de blé) sont les deux autres piliers. L'armement et le nucléaire civil (Rosatom) sont des niches de haute technologie à forte valeur géopolitique.}
\end{popfigure}

\begin{methode}[Analyser la structure des exportations d'un pays rentier]
\begin{enumerate}
    \item \textbf{Classer les exportations par produit.} Regrouper les exportations par grandes catégories (hydrocarbures, métaux, produits agricoles, armement, etc.). Calculer la part des trois premiers postes : si elle dépasse 50 \%, on parle de forte spécialisation rentière.
    \item \textbf{Cartographier les destinations.} Identifier les cinq premiers clients et calculer leur part cumulée. Une concentration géographique élevée signifie une vulnérabilité politique (risque de sanctions, rupture de marché).
    \item \textbf{Mesurer la concentration.} On peut utiliser l'indice de Hirschman-Herfindahl : $HHI = \sum (s_i)^2$, où $s_i$ est la part (en \%) du produit ou du client $i$. Un HHI supérieur à 2 500 traduit une concentration forte.
    \item \textbf{Évaluer la vulnérabilité.} Croiser concentration par produit, concentration géographique et sensibilité aux sanctions (technologies, financement, assurance, transport maritime). Identifier les stratégies de contournement (flotte parallèle, paiements en monnaies nationales, réexportation via des pays tiers).
\end{enumerate}
\end{methode}

\begin{exoresolu}[Analyse de la résilience des exportations russes 2022-2023]
\textbf{Énoncé :} À partir de la figure 1 et de vos connaissances, expliquez pourquoi les recettes d'exportation de la Russie ne se sont pas effondrées en 2023 malgré l'embargo européen sur le pétrole brut russe (décembre 2022) et sur les produits raffinés (février 2023), ainsi que le plafonnement du prix du baril russe à 60 \$ décidé par le G7.
\tcblower
\textbf{Solution détaillée :}
\begin{enumerate}
    \item \textbf{Réorientation géographique massive :} l'UE absorbait environ la moitié des exportations pétrolières russes en 2021. En 2023, l'Inde et la Chine absorbent l'essentiel des flux maritimes russes (bruts Urals et ESPO), achetés avec une décote par rapport au Brent, mais souvent au-dessus du plafond de 60 \$.
    \item \textbf{Flotte parallèle (\emph{shadow fleet}) :} la Russie s'appuie sur une flotte de plusieurs centaines de pétroliers vieillissants, opérés par des sociétés-écrans et assurés hors des clubs occidentaux, pour contourner l'interdiction d'assurance et de transport imposée par le G7 au-delà du plafond de prix.
    \item \textbf{Des prix de l'énergie encore rémunérateurs :} le Brent a tourné en moyenne autour de 80-85 \$ le baril en 2023. Même avec décote, le baril russe reste vendu à un prix supérieur au seuil d'équilibre budgétaire russe (de l'ordre de 45-50 \$).
    \item \textbf{Diversification partielle :} hausse des exportations de produits raffinés vers la Turquie, l'Afrique et l'Amérique latine ; poursuite des exportations de GNL (Yamal LNG) vers l'Asie et même vers l'Europe, le GNL russe n'étant pas frappé d'embargo.
    \item \textbf{Réserves et changes :} l'excédent courant record de 2022 (plus de 200 Md \$) a permis d'accumuler des réserves (or, yuan) et de financer l'effort de guerre sans recours massif à l'endettement extérieur.
\end{enumerate}
\textbf{Conclusion :} l'économie russe a subi une \emph{transformation forcée} (réallocation des ressources vers la défense, inflation, pénurie de main-d'œuvre) mais pas un effondrement, grâce à la rente énergétique et à une adaptation logistique et financière rapide.
\end{exoresolu}

\begin{attention}[Erreur fréquente : sous-estimer la résilience russe]
Ne pas confondre \textbf{effondrement} et \textbf{transformation forcée}. Malgré des sanctions d'une ampleur historique, le PIB réel russe a crû d'environ 3,6 \% en 2023 (FMI, Rosstat). Moteurs : économie de guerre (dépenses militaires de l'ordre de 6 à 7 \% du PIB, stimulant la demande interne), réorientation des flux vers l'Asie, prix de l'énergie encore rémunérateurs, politique monétaire très stricte (taux directeur porté à 16 \% fin 2023) pour contenir l'inflation (environ 7 \%). Piège au Bac : écrire « la Russie est ruinée par les sanctions ». Ce qui est vrai : le potentiel de croissance à long terme est hypothéqué (investissement, technologie, démographie). Ce qui est faux : l'effondrement à court terme.
\end{attention}

\courssub{Atouts majeurs et faiblesses structurelles : le paradoxe russe}

\emph{Atouts (la carte du pouvoir).}
\begin{itemize}
    \item \textbf{Ressources immenses :} la Sibérie et l'Arctique russe concentrent des richesses considérables. Gazprom (premier détenteur mondial de réserves de gaz), Rosneft (major pétrolier public), Norilsk Nickel (palladium, nickel), Alrosa (diamants). \textbf{Agriculture :} la Russie est le premier exportateur mondial de blé (de l'ordre de 45 à 50 Mt par an), nourrissant l'Afrique du Nord, le Moyen-Orient et une partie de l'Asie.
    \item \textbf{Capital humain scientifique et technique :} héritage soviétique d'ingénieurs nucléaires, aérospatiaux et informaticiens (Yandex, Kaspersky). Mais ce capital s'érode : fuite des cerveaux massive depuis 2022 (plusieurs centaines de milliers de départs de personnes qualifiées).
    \item \textbf{Industrie de défense et nucléaire civil (Rosatom) :} Rosatom est le premier exportateur mondial de centrales nucléaires (projets en Turquie, en Égypte, en Hongrie, au Bangladesh, en Chine, en Inde) et maîtrise le cycle complet du combustible (environ 40 \% des capacités mondiales d'enrichissement). Le complexe militaro-industriel (Rostec) assure une large autosuffisance (missiles, blindés, aviation, guerre électronique).
    \item \textbf{Géographie stratégique :} la Route maritime du Nord (NSR) offre un raccourci Europe-Asie (gain de 15 à 20 jours par rapport à Suez). La Russie projette sa puissance dans l'Arctique (bases militaires, brise-glaces).
\end{itemize}

\emph{Faiblesses (le talon d'Achille).}
\begin{itemize}
    \item \textbf{Démographie dégradée :} population d'environ 146 millions d'habitants (2024), vieillissante, fécondité faible (environ 1,4 enfant par femme), espérance de vie des hommes autour de 67-68 ans. Le solde naturel est négatif et la pénurie de main-d'œuvre devient structurelle (taux de chômage historiquement bas, autour de 2,5-3 \%).
    \item \textbf{Dépendance technologique occidentale :} machines-outils, semi-conducteurs, logiciels industriels, équipements pétroliers et gaziers sophistiqués (forages complexes, GNL), composants aéronautiques. Le contournement via des pays tiers (Chine, Turquie, Kazakhstan, Émirats) est possible mais augmente les coûts, les délais et dégrade la qualité.
    \item \textbf{Infrastructures inadaptées au nouveau contexte :} les gazoducs vers l'UE (Yamal-Europe, Nord Stream 1 et 2, Fraternité via l'Ukraine) sont largement inutilisés ou détruits. Les nouveaux gazoducs vers la Chine (Force de Sibérie 1, et le projet Force de Sibérie 2) coûtent des dizaines de milliards de dollars et demandent de longs délais.
    \item \textbf{Échec relatif de la diversification :} le projet Skolkovo (la « Silicon Valley russe ») n'a pas créé d'écosystème d'innovation autonome. Hors défense et nucléaire, la valeur ajoutée du secteur non énergétique reste faible.
\end{itemize}

\begin{exemplebox}[Chiffres clés : le basculement gazier Russie-UE-Chine]
\begin{itemize}
    \item Exportations de gaz par gazoduc vers l'UE : environ \textbf{150-155 Gm\textsuperscript{3} en 2021} contre environ \textbf{25-30 Gm\textsuperscript{3} en 2023} (via l'Ukraine et TurkStream). Un effondrement historique.
    \item Exportations de gaz vers la Chine (Force de Sibérie 1) : \textbf{0 en 2019}, environ \textbf{23 Gm\textsuperscript{3} en 2023}, avec une cible de \textbf{38 Gm\textsuperscript{3} en 2025} (pleine capacité).
    \item Force de Sibérie 2 (via la Mongolie) : capacité envisagée de 50 Gm\textsuperscript{3}, mais négociations difficiles sur les prix et les délais, la Chine étant en position de force.
    \item Réserves d'or de la Banque de Russie : environ \textbf{2 300 tonnes} (parmi les cinq premières mondiales), stockées en Russie et donc non gelables.
    \item Budget fédéral 2024 : forte hausse des dépenses de défense (de l'ordre de 10 800 milliards de roubles, soit environ 110 Md \$), pour un déficit public maîtrisé (environ 1 \% du PIB).
\end{itemize}
\end{exemplebox}

\courssub{Géopolitique économique : nouveaux flux et intégration eurasiatique}

La réorientation des flux ne se fait pas au hasard : elle suit des infrastructures lourdes (gazoducs, ports, voies ferrées) et des cadres institutionnels (UEEA, BRICS, OCS).

\begin{popfigure}
\includegraphics[width=\linewidth,height=9.5cm,keepaspectratio]{N13/arctic_en.jpg}
\legende{L'Arctique vu du pôle Nord, avec le Passage du Nord-Ouest (rose, par le Canada), le Passage du Nord-Est (orange) et la Route maritime du Nord (NSR, pointillé jaune et blanc) le long de la façade arctique de la Russie, entre le détroit de Béring et la Norvège. La carte montre aussi la profondeur de l'océan (en bleu). Enjeux d'après le cours : la Russie possède la plus longue façade arctique du monde (environ \num{24000} km de côtes) et l'unique flotte de brise-glaces à propulsion nucléaire ; la NSR est un raccourci entre l'Europe et l'Asie (15 à 20 jours de moins que par Suez), avec les ressources offshore (gaz de Yamal, pétrole de la mer de Kara). Carte en anglais. \\ Source : Wikimedia Commons, « Map of the Arctic region showing the Northeast Passage, the Northern Sea Route and Northwest Passage, and bathymetry », Susie Harder, domaine public.}
\end{popfigure}

\begin{itemize}
    \item \textbf{Gazoducs :} \emph{Force de Sibérie 1} (capacité de 38 Gm\textsuperscript{3}/an, en service depuis 2019, vers la Chine). \emph{Force de Sibérie 2} (via la Mongolie, 50 Gm\textsuperscript{3} envisagés, en négociation). \emph{TurkStream} (31,5 Gm\textsuperscript{3}, vers la Turquie et les Balkans). \emph{Nord Stream 1 et 2} (110 Gm\textsuperscript{3} de capacité, sabotés en septembre 2022).
    \item \textbf{Union économique eurasiatique (UEEA) :} Russie, Biélorussie, Kazakhstan, Arménie, Kirghizistan. Marché commun d'environ 180 millions d'habitants (PIB cumulé de l'ordre de 2 500 Md \$). Libre circulation des biens, des services, des capitaux et de la main-d'œuvre. La Russie en est le moteur (environ 85-90 \% du PIB de l'Union). Le Kazakhstan est un partenaire clé (pétrole, uranium, transit Chine-Europe).
    \item \textbf{Route Maritime du Nord (NSR) :} elle devient une \cle{artère stratégique}. Trafic 2023 : environ 36 Mt (record). Objectif officiel : 150 à 200 Mt/an à l'horizon 2030. Elle nécessite des brise-glaces (flotte nucléaire unique au monde), des ports (Mourmansk, Arkhangelsk, Dikson, Pevek) et des satellites d'observation (Arktika-M).
\end{itemize}

\begin{savaistu}[La flotte de brise-glaces nucléaires : la clé de l'Arctique]
La Russie est le \textbf{seul pays au monde} à exploiter une flotte de brise-glaces civils à propulsion nucléaire (Rosatomflot). Les brise-glaces du projet 22220 (\emph{Arktika}, \emph{Sibir}, \emph{Ural}) peuvent briser une banquise de près de 3 mètres d'épaisseur. Pourquoi le nucléaire ? Parce qu'il offre une autonomie de plusieurs années sans ravitaillage en combustible et une puissance considérable pour ouvrir la voie aux convois de méthaniers (Yamal LNG) et de porte-conteneurs, y compris en hiver. L'enjeu est de transformer l'Arctique en « autoroute maritime » russe, génératrice de droits de passage, et d'y projeter la puissance navale.
\end{savaistu}

\courssub{Le Cameroun et la Russie : un partenariat minoritaire mais symbolique}

Si les échanges commerciaux restent modestes (de l'ordre de quelques dizaines de milliards de FCFA par an, largement déficitaires pour le Cameroun : importations d'engrais, de blé, d'hydrocarbures et de matériel militaire ; exportations de bois, de cacao, de coton et d'aluminium brut), la relation a une dimension \cle{stratégique et diplomatique} :

\begin{itemize}
    \item \textbf{Coopération militaire :} formation d'officiers camerounais dans des académies russes ; livraisons d'équipements (hélicoptères Mi-17, blindés, armes légères). La Russie constitue une alternative aux fournisseurs traditionnels (France, États-Unis, Chine, Israël).
    \item \textbf{Énergie :} discussions avec Rosatom sur le nucléaire civil (accords-cadres, études préliminaires) ; intérêt russe pour le secteur gazier camerounais.
    \item \textbf{Éducation et culture :} bourses gouvernementales russes pour des étudiants camerounais (médecine, ingénierie, économie) ; enseignement de la langue russe à Yaoundé ; petite diaspora camerounaise en Russie.
    \item \textbf{Géologie et mines :} expertise russe en cartographie géologique et en exploration minière, en alternative aux compagnies occidentales et chinoises.
    \item \textbf{Diplomatie (ONU) :} le Cameroun choisit souvent l'\emph{abstention} lors des résolutions de l'Assemblée générale condamnant la Russie. Cette position de \emph{non-alignement actif} permet de préserver les relations avec l'Occident (aide, FMI, UE) tout en cultivant l'alternative russe (armement, soutien diplomatique).
\end{itemize}

\begin{aretenir}[Synthèse : la Russie, foyer moteur \emph{sui generis}]
La Russie n'est pas un foyer moteur classique (comme l'UE, les États-Unis ou la Chine) tirant sa puissance de l'innovation diffuse, de la consommation de masse ou de l'intégration par le marché. C'est un \textbf{foyer de puissance par la rente et la contrainte} :
\begin{enumerate}
    \item \textbf{Rente énergétique et minière} monétisée pour financer l'État, l'armée, le nucléaire et l'espace.
    \item \textbf{Réponse à la contrainte géopolitique :} sanctions $\rightarrow$ Pivot vers l'Asie (la Chine comme partenaire dominant, l'Inde comme acheteur opportuniste) + OPEP+ (gestion des prix) + BRICS (cadre politique non occidental).
    \item \textbf{Vulnérabilités critiques :} démographie, technologie, dépendance croissante à la Chine (asymétrie : la Russie vend des ressources brutes, la Chine vend des technologies et des biens finis).
    \item \textbf{Projet arctique :} tentative de transformer la géographie (fonte de la banquise) en avantage durable (NSR, ressources offshore).
\end{enumerate}
Pour le Cameroun, la Russie est un \textbf{partenaire de diversification} (armement, diplomatie, formation), non un moteur de développement économique structurel.
\end{aretenir}

\begin{exercice}[Entraînement Bac : la Russie dans la mondialisation]
\textbf{Document 1 :} Extrait du discours de V. Poutine au Forum de Valdaï, octobre 2023 : « Le monde change... La majorité mondiale ne veut plus se plier aux règles dictées par l'Occident collectif... La Russie construit un partenariat stratégique avec la Chine et développe ses relations avec l'Inde, le monde islamique, l'Afrique et l'Amérique latine. »

\textbf{Document 2 :} Carte des flux gaziers russes en 2021 et en 2023 (flèches d'épaisseur proportionnelle : UE : 150 $\rightarrow$ 25-30 Gm\textsuperscript{3} ; Chine : 0 $\rightarrow$ 23 Gm\textsuperscript{3} ; Turquie : stable, environ 15-20 Gm\textsuperscript{3}).

\textbf{Document 3 :} Tableau du PIB nominal de la Russie (Md \$ courants, FMI, ordres de grandeur) : 2013 : 2 200 ; 2021 : 1 840 ; 2022 : 2 240 ; 2023 : 2 020.

\textbf{Questions :}
\begin{enumerate}
    \item Caractérisez la « majorité mondiale » évoquée par V. Poutine (composition, intérêts communs).
    \item Analysez la réorientation des flux gaziers (document 2) : causes, conséquences pour la Russie (revenus, infrastructures) et pour l'UE (sécurité énergétique, prix).
    \item Interprétez l'évolution du PIB nominal (document 3) : effet prix (pétrole/gaz), effet change (rouble), effet volume. La Russie est-elle en récession ?
    \item Croquis : sur un planisphère simplifié, représentez le \emph{Pivot vers l'Asie} (flux énergétiques, partenariats au sein des BRICS et de l'OCS, NSR). Légendez soigneusement.
\end{enumerate}
\end{exercice}

\begin{corrige}[Exercice : la Russie dans la mondialisation]
\begin{enumerate}
    \item \textbf{Majorité mondiale :} concept russe (proche de celui de \emph{Sud global}) désignant les pays qui ne se sont pas alignés sur les sanctions occidentales : Chine, Inde, Indonésie, Brésil, Afrique du Sud (BRICS), Iran, Turquie, pays arabes, majeure partie de l'Afrique et de l'Amérique latine. Intérêts communs : multipolarité, non-ingérence, accès aux ressources et aux marchés russes et chinois, réforme de la gouvernance mondiale (FMI, Banque mondiale, Conseil de sécurité).
    \item \textbf{Réorientation du gaz :} \emph{Causes} : embargo et sanctions de l'UE, sabotage de Nord Stream, décision russe de réduire les livraisons (arme énergétique). \emph{Conséquences pour la Russie} : perte d'un marché captif et rémunérateur, nécessité d'investissements massifs dans de nouveaux gazoducs (Force de Sibérie 2, coûteux et long à construire), dépendance accrue à une Chine en position de force pour négocier les prix. \emph{Conséquences pour l'UE} : choc des prix en 2022, diversification coûteuse (GNL américain, gaz norvégien et algérien), accélération de la transition énergétique, perte de compétitivité des industries énergivores (chimie, sidérurgie).
    \item \textbf{PIB nominal :} le pic de 2022 s'explique par des prix de l'énergie records et un rouble soutenu (contrôle des capitaux, excédent courant exceptionnel). La baisse de 2023 traduit la normalisation des prix et l'affaiblissement du rouble (environ 90-100 roubles pour 1 dollar). \textbf{Il n'y a pas de récession technique :} le PIB réel (en volume) a progressé d'environ 3,6 \% en 2023, tiré par la dépense publique de guerre, la substitution aux importations et la consommation des ménages (salaires réels en hausse du fait de la pénurie de main-d'œuvre).
    \item \textbf{Croquis attendu :} planisphère simplifié avec la Russie au centre. Flèches épaisses : pétrole et GNL vers l'Inde, la Chine, la Turquie. Flèches de gazoducs : Force de Sibérie 1 (et 2 en pointillés) vers la Chine ; TurkStream vers la Turquie et les Balkans. Flèche maritime : NSR de Mourmansk au détroit de Béring. Symboles : centrale nucléaire (exportations Rosatom), épi de blé (exportations vers l'Afrique et le Moyen-Orient), matériel militaire. Légende organisée en quatre rubriques : flux énergétiques, flux commerciaux non énergétiques, coopération politico-militaire (BRICS, OCS, UEEA), infrastructures stratégiques (NSR, gazoducs). Ne pas oublier le titre, l'orientation et l'échelle indicative.
\end{enumerate}
\end{corrige}

\coursec{Synthèse : Hiérarchie, interdépendances et croquis de synthèse mondiale}

Après avoir étudié la Russie, ce foyer singulier dont la puissance repose sur l'énergie et les armes plutôt que sur l'industrie manufacturière, nous disposons désormais de tous les éléments du puzzle. Il reste à les assembler. C'est toute l'ambition de cette synthèse : hiérarchiser les cinq foyers moteurs, comprendre les interdépendances qui les relient, identifier les relais émergents et la gouvernance mondiale, situer le Cameroun, puis construire pas à pas le croquis final qui pourrait t'être demandé au Baccalauréat.

\courssub{Hiérarchiser les foyers moteurs : le tableau comparatif}

Comparer des foyers aussi différents exige des indicateurs communs. Retiens la grille de lecture suivante : la \cle{masse} (PIB total), la \cle{richesse individuelle} (PIB par habitant, IDH), l'\cle{insertion commerciale} (part dans les exportations mondiales), l'\cle{attractivité et l'influence financière} (IDE entrants et sortants), la \cle{capacité d'innovation} (dépenses de R\&D), la \cle{démographie}, la \cle{spécialisation} et enfin la \cle{vulnérabilité}. Aucun foyer n'est fort sur tous les critères : c'est précisément ce qui crée les interdépendances.

\begin{center}
\small
\begin{tabularx}{\linewidth}{|l|X|X|X|X|X|}
\hline
\rowcolor{popblueL} \textbf{Indicateur} & \textbf{UE} & \textbf{USA} & \textbf{Chine} & \textbf{Japon / Corée du Sud} & \textbf{Russie} \\
\hline
PIB (ordre de grandeur, 2023) & $\approx$ 18 000 Md\$ & $\approx$ 27 000 Md\$ & $\approx$ 17 500 Md\$ & $\approx$ 4 200 / 1 700 Md\$ & $\approx$ 2 000 Md\$ \\
\hline
PIB/hab & $\approx$ 40 000 \$ & $\approx$ 80 000 \$ & $\approx$ 12 500 \$ & $\approx$ 34 000 / 33 000 \$ & $\approx$ 13 000 \$ \\
\hline
IDH & Très élevé & Très élevé & Élevé & Très élevé & Élevé \\
\hline
Part export. mondiales de biens & $\approx$ 14--15 \% (hors intra-UE) & $\approx$ 8--9 \% & $\approx$ 14 \% (1\textsuperscript{re}) & $\approx$ 3 / 3 \% & $\approx$ 2 \% \\
\hline
IDE (stock) & Très fort entrant et sortant & 1\textsuperscript{er} entrant et sortant & Fort entrant, sortant croissant & Sortant $>$ entrant & Faible, en retrait depuis 2022 \\
\hline
R\&D (\% du PIB) & $\approx$ 2,3 \% & $\approx$ 3,5 \% & $\approx$ 2,5 \% & $\approx$ 3,3 / 4,9 \% & $\approx$ 1 \% \\
\hline
Démographie & $\approx$ 450 M hab, vieillissante & $\approx$ 335 M, dynamique (immigration) & $\approx$ 1,4 Md, déclin amorcé & Vieillissement avancé & $\approx$ 145 M, en déclin \\
\hline
Spécialisation & Industrie haut de gamme, services, agroalimentaire & Technologies, finance, énergie, agriculture & Manufacture, BTP, électronique & Électronique, automobile, robotique & Hydrocarbures, armement, céréales \\
\hline
Vulnérabilité & Dépendance énergétique, fragmentation politique & Endettement, polarisation interne & Dépendance aux puces haut de gamme, tensions avec l'Occident & Dépendance énergétique et alimentaire, démographie & Dépendance aux hydrocarbures, sanctions, isolement \\
\hline
\end{tabularx}
\end{center}

\begin{aretenir}[Trois enseignements du tableau]
\begin{enumerate}
\item La \cle{Triade} historique (Amérique du Nord, Europe occidentale, Japon) domine encore par la richesse individuelle, la R\&D et la finance.
\item La \cle{Chine} a rejoint le club des moteurs par la masse (PIB, commerce, manufacture) mais reste un pays à revenu intermédiaire par habitant.
\item La \cle{Russie} est un foyer incomplet : puissance énergétique et militaire, mais poids économique comparable à celui de l'Italie ou de l'Espagne.
\end{enumerate}
\end{aretenir}

\courssub{Les interdépendances : la voiture, fil rouge de la mondialisation}

Pour comprendre concrètement ce qu'est une \cle{chaîne de valeur mondiale}, suis le parcours d'une voiture allemande vendue en 2024 :

\begin{itemize}
\item la \textbf{conception} et le marketing sont réalisés en Allemagne (haute valeur ajoutée, R\&D) ;
\item les \textbf{puces électroniques} viennent de Taïwan (TSMC) ou de Corée du Sud (Samsung), les écrans et batteries du Japon ou de Chine ;
\item l'\textbf{acier} peut être produit à partir de minerai brésilien et, jusqu'en 2022, d'énergie russe bon marché ;
\item l'\textbf{assemblage} a lieu en Slovaquie (main-d'œuvre qualifiée et moins chère, au sein du marché unique) ou en Chine pour le marché chinois ;
\item la \textbf{vente} se fait sur les grands marchés solvables : États-Unis, Chine, Europe.
\end{itemize}

Chaque étape est localisée là où le rapport coût/qualité est optimal. C'est la \cle{fragmentation internationale du processus productif}, rendue possible par la baisse des coûts de transport (conteneur) et de communication (internet).

\begin{attention}[La fragilité des chaînes de valeur]
Ces chaînes sont efficaces mais vulnérables. Trois chocs récents le prouvent : la pandémie de \textbf{Covid-19} (2020) a bloqué les ports et provoqué la pénurie de puces ; la \textbf{guerre en Ukraine} (2022) a coupé l'Europe du gaz russe ; une crise autour de \textbf{Taïwan} priverait le monde de l'essentiel des puces avancées. D'où les stratégies actuelles de \emph{réindustrialisation}, de \emph{nearshoring} (relocalisation proche) et de \emph{friendshoring} (relocalisation chez des alliés).
\end{attention}

\begin{propriete}[La loi de gravité appliquée à l'économie]
L'intensité des flux (commerciaux, financiers, migratoires) entre deux pôles est approximativement \textbf{proportionnelle au produit de leurs masses économiques} (PIB) et \textbf{inversement proportionnelle à la distance} qui les sépare (coûts de transport, barrières tarifaires, différences culturelles) :
\[ F_{ij} \propto \frac{PIB_i \times PIB_j}{D_{ij}} \]
Cette loi explique pourquoi les flux sont denses \emph{à l'intérieur} de la Triade et entre l'Asie de l'Est et l'Occident (masses énormes, distance réduite par le maritime), et pourquoi les flux \emph{Afrique-Afrique} restent faibles (masses faibles, distances internes coûteuses faute d'infrastructures). Elle n'est pas exigible en tant que formule au Bac, mais elle structure le raisonnement.
\end{propriete}

\courssub{Les nouveaux relais de la mondialisation}

La carte des foyers moteurs n'est pas figée. Quatre ensembles jouent un rôle croissant de \cle{relais} :

\begin{itemize}
\item \textbf{L'Inde} : premier pays démographique du monde depuis 2023 ($\approx$ 1,43 Md d'habitants), championne des services informatiques (Bangalore) et de la pharmacie générique, avec une croissance parmi les plus fortes du G20.
\item \textbf{L'ASEAN} : le Vietnam et l'Indonésie notamment deviennent des \emph{usines alternatives} à la Chine (stratégie dite \emph{China + 1} des multinationales), pour l'électronique, le textile et la chaussure.
\item \textbf{Le Golfe} : les pétrodollars sont recyclés dans des fonds souverains géants, des hubs logistiques (Dubaï, Doha) et aériens (Emirates, Qatar Airways) qui connectent les continents.
\item \textbf{L'Afrique} : réservoir de ressources (cobalt de RDC, pétrole du Nigeria, cacao de Côte d'Ivoire et du Ghana) et surtout \emph{marché d'avenir} : sa population devrait atteindre 2,5 milliards en 2050.
\end{itemize}

\begin{exemplebox}[L'Afrique, encore à la marge des flux]
La part de l'Afrique dans le commerce mondial reste inférieure à \textbf{3 \%}. Elle capte environ \textbf{5 \%} des IDE mondiaux entrants, concentrés sur les ressources extractives (pétrole, mines). Le commerce \emph{intra-africain} ne représente qu'environ 15 \% du commerce du continent (contre plus de 60 \% en Europe). Selon la Banque mondiale, la ZLECAf pourrait accroître le commerce intra-africain de plus de \textbf{50 \%} d'ici 2035 si elle est pleinement mise en œuvre. Ordres de grandeur à connaître, sources : CNUCED, Banque mondiale.
\end{exemplebox}

\begin{definition}[La ZLECAf]
La \cle{Zone de Libre-Échange Continentale Africaine} (ZLECAf), entrée en vigueur en 2021 et dont le siège est à Accra (Ghana), vise à créer un marché unique africain de plus de \textbf{1,3 milliard d'habitants} et d'environ \textbf{3 400 milliards de dollars} de PIB cumulé. C'est l'outil qui pourrait transformer l'Afrique de simple fournisseur de matières premières en \emph{relais}, voire en futur foyer moteur. Son succès dépend de conditions précises : infrastructures de transport, système de paiement panafricain (PAPSS), règles d'origine harmonisées, stabilité politique et paix.
\end{definition}

\begin{savaistu}[Le câble 2Africa : l'autoroute numérique du continent]
Le câble sous-marin \emph{2Africa}, déployé par un consortium réunissant Meta, Orange, MTN et d'autres opérateurs, mesure environ \textbf{45 000 km} : c'est le plus long câble sous-marin du monde. Mis en service progressivement depuis 2023--2024, il relie 33 pays sur trois continents (Europe, Afrique, Asie) et ceinture presque tout le littoral africain, avec des points d'atterrissage au Cameroun (Kribi, Douala). Il doit multiplier la connectivité internet du continent — condition invisible mais décisive de l'intégration africaine dans l'économie mondiale.
\end{savaistu}

\courssub{La gouvernance mondiale : qui décide ?}

La mondialisation économique possède des instances de régulation, mais leur légitimité est contestée :

\begin{itemize}
\item Le \textbf{G7} (États-Unis, Japon, Allemagne, France, Royaume-Uni, Italie, Canada) rassemble la Triade historique : moins de 10 \% de la population mondiale, mais encore environ 40--45 \% du PIB mondial.
\item Le \textbf{G20}, créé en 1999 et promu au sommet après la crise de 2008, intègre la Chine, la Russie, l'Inde, le Brésil, l'Afrique du Sud : il reconnaît le déplacement du centre de gravité économique vers le Sud.
\item Les \textbf{BRICS+} (Brésil, Russie, Inde, Chine, Afrique du Sud, élargis en 2024 à l'Égypte, l'Éthiopie, l'Iran, l'Arabie saoudite, les Émirats arabes unis) se présentent comme une alternative du \emph{Sud global} à la gouvernance occidentale.
\item L'\textbf{OMC}, censée arbitrer le commerce mondial, est affaiblie : son \cle{Organe d'appel} est bloqué depuis 2019 par le refus américain de nommer de nouveaux juges, ce qui paralyse le règlement des différends.
\end{itemize}

\begin{aretenir}[L'essentiel sur la gouvernance]
La gouvernance mondiale reflète le déplacement des puissances : le G7 incarne l'ordre ancien de la Triade, le G20 l'ordre multipolaire actuel, les BRICS+ la contestation venue du Sud. L'affaiblissement de l'OMC marque le passage d'une mondialisation \emph{régulée par le droit} à une mondialisation de plus en plus \emph{géopolitisée} (sanctions, guerres commerciales, blocs).
\end{aretenir}

\courssub{La place du Cameroun dans ce système mondial}

Situation concrète : au port de Douala, on charge du bois, du cacao, du coton et de l'aluminium à destination de l'Europe et de la Chine ; on décharge des conteneurs remplis de machines, de véhicules, de téléphones et de produits manufacturés. Ce constat résume la position camerounaise : \textbf{fournisseur de matières premières peu transformées, importateur de biens d'équipement et de consommation}. Résultat : une balance commerciale structurellement déficitaire et une faible capture de la valeur ajoutée.

Les pistes pour monter en gamme sont connues :

\begin{enumerate}
\item la \textbf{transformation locale} (usines de chocolat, sciages, filature, raffinage) soutenue par la ZLECAf ;
\item l'\textbf{attraction d'IDE manufacturiers} grâce aux zones économiques spéciales (Kribi, Dibamba) ;
\item la \textbf{négociation d'accords commerciaux} : Accord de Partenariat Économique (APE) avec l'Union européenne, en vigueur pour le Cameroun depuis 2016, et partenariats avec la Chine dans le cadre du Forum sur la Coopération sino-africaine (FOCAC).
\end{enumerate}

\begin{attention}[Précision sur l'AGOA]
L'\textbf{AGOA} (\emph{African Growth and Opportunity Act}) est la loi américaine qui ouvre le marché des États-Unis aux produits africains en franchise de droits de douane. Le Cameroun en a été \textbf{exclu en janvier 2020} par décision américaine, en raison de violations des droits humains. Ne cite donc pas l'AGOA comme un atout actuel du Cameroun : c'est au contraire un exemple de \emph{conditionnalité politique} des échanges, qui montre que l'accès aux marchés des foyers moteurs n'est jamais garanti.
\end{attention}

\begin{exoresolu}[Commenter un tableau statistique : l'Afrique dans les flux mondiaux]
\textbf{Énoncé.} On donne : part de l'Afrique dans le commerce mondial de biens $\approx$ 2,5 \% ; part dans les IDE entrants mondiaux $\approx$ 5 \% ; part du commerce intra-africain dans le commerce total du continent $\approx$ 15 \% (contre $\approx$ 60 \% pour l'Europe). À partir de ces données et de tes connaissances, montre en une vingtaine de lignes que l'Afrique reste marginalisée dans la mondialisation mais dispose de leviers pour devenir un relais.
\tcblower
\textbf{Corrigé type.} \emph{Introduction} : les chiffres montrent un paradoxe — continent de 1,4 milliard d'habitants, l'Afrique pèse moins de 3 \% du commerce mondial. \emph{Développement} : (1) la marginalisation tient à la nature des échanges (exportations de matières premières brutes, importations de produits manufacturés), à la faiblesse des IDE (5 \%, concentrés sur l'extraction) et surtout à la désintégration interne : 15 \% de commerce intra-africain contre 60 \% en Europe, ce que la loi de gravité explique par des masses économiques nationales faibles et des distances internes coûteuses (infrastructures insuffisantes). (2) Pourtant des leviers existent : la ZLECAf (marché de 1,3 milliard de consommateurs), la connectivité (câble 2Africa), la démographie (2,5 milliards d'habitants en 2050) et les ressources stratégiques (cobalt, cacao, pétrole). \emph{Conclusion} : l'Afrique peut passer du statut de périphérie fournisseuse à celui de relais, à condition de transformer localement et de s'intégrer régionalement — enjeu central pour le Cameroun.
\end{exoresolu}

\courssub{Construire le croquis de synthèse : « Les foyers moteurs de l'économie mondiale et la place de l'Afrique »}

\begin{methode}[Construire le croquis de synthèse en sept étapes (méthode d'examen)]
\begin{enumerate}
\item \textbf{Fond de carte} : trace un planisphère simplifié (projection type Robinson ou Peters), continents en blocs schématiques, avec la flèche du nord et une échelle indicative.
\item \textbf{Figer les cinq foyers} : colore chaque pôle d'une couleur distincte (UE, Amérique du Nord, Asie de l'Est, Russie) ; tu peux hachurer selon la densité du PIB.
\item \textbf{Tracer les flux majeurs} avec des flèches d'\emph{épaisseur proportionnelle} : IDE et finance (Nord $\rightarrow$ Sud et intra-Triade), biens manufacturés (Asie $\rightarrow$ Occident), énergie et matières premières (Russie, Golfe, Afrique $\rightarrow$ pôles), données (câbles).
\item \textbf{Placer les interfaces} : les grands ports mondiaux (Shanghai, Singapour, Rotterdam, Los Angeles...), les câbles sous-marins, les hubs aériens (Dubaï, Doha).
\item \textbf{Signaler les relais émergents} (Inde, ASEAN, Golfe, Afrique) avec un symbole distinct des pôles historiques.
\item \textbf{Organiser la légende} en quatre rubriques : Pôles / Flux / Interfaces / Relais.
\item \textbf{Finaliser} : titre précis, nom de l'auteur, sources, date.
\end{enumerate}
\end{methode}

\begin{popfigure}
\includegraphics[width=\linewidth,height=9.5cm,keepaspectratio]{N13/brics.png}
\legende{Fond de carte pour le croquis de synthèse : les membres des BRICS (en bleu : Brésil, Russie, Inde, Chine, Afrique du Sud, Égypte, Éthiopie, Iran, Indonésie, Émirats arabes unis) sur un planisphère gris. On y lit la place de l'Afrique : seuls trois pays africains (Afrique du Sud, Égypte, Éthiopie) y figurent. \textbf{Le croquis attendu situe :} pôles (Amérique du Nord, UE, Asie de l'Est, Russie) ; relais (Inde, ASEAN, Golfe, Afrique) ; flux (biens manufacturés, énergie et matières premières, IDE du Nord vers le Sud, câble sous-marin 2Africa) ; interfaces (grands ports mondiaux). \\ Source : Wikimedia Commons, « BRICS.svg », Cflm001, domaine public.}
\end{popfigure}

\begin{attention}[Les trois erreurs fréquentes au croquis]
\begin{enumerate}
\item \textbf{Oublier les flux invisibles} : les flux financiers et numériques (données, câbles sous-marins) sont majoritaires en valeur ; un croquis sans eux est incomplet.
\item \textbf{Ne pas hiérarchiser les flèches} : si toutes les flèches ont la même épaisseur, le correcteur ne voit ni les flux majeurs (Asie $\rightarrow$ Occident) ni les flux mineurs (Afrique $\rightarrow$ monde). L'épaisseur doit être proportionnelle à l'importance du flux.
\item \textbf{Confondre pôles et relais} : les foyers historiques (Triade) et les foyers émergents (Inde, ASEAN, Afrique) doivent être distingués visuellement (couleurs pleines contre symboles ou teintes claires).
\end{enumerate}
\end{attention}

\begin{exercice}[Pour s'entraîner]
Sur une feuille blanche, construis en 30 minutes le croquis « Les foyers moteurs de l'économie mondiale et la place de l'Afrique » en appliquant la méthode en sept étapes. Rédige ensuite un commentaire de 15 lignes qui dégage deux idées : (a) la persistance de la Triade face à la montée de l'Asie de l'Est ; (b) les conditions pour que l'Afrique, et le Cameroun en particulier, deviennent un relais de la mondialisation.
\end{exercice}

\begin{corrige}[Exercice]
Attendus du croquis : fond de planisphère orienté ; cinq foyers colorés et distingués des relais ; au moins trois familles de flux hiérarchisés (biens, énergie/matières, IDE/données) ; interfaces (ports, câble 2Africa) ; légende en quatre rubriques ; titre, auteur, sources, date. Attendus du commentaire : (a) la Triade conserve la richesse, la R\&D et la finance, mais la Chine est devenue premier exportateur mondial de biens et l'Asie de l'Est l'atelier et le laboratoire du monde ; (b) pour l'Afrique : transformation locale des matières premières, mise en œuvre effective de la ZLECAf, infrastructures et connectivité, stabilité politique ; pour le Cameroun : valoriser le corridor Douala--Kribi, l'APE avec l'Union européenne et le partenariat avec la Chine (FOCAC) pour attirer des IDE manufacturiers plutôt qu'extractifs, tout en sachant que l'accès aux marchés peut être conditionné (exclusion de l'AGOA en 2020).
\end{corrige}

\begin{aretenir}[Conclusion de la leçon]
L'économie mondiale est animée par une poignée de \cle{foyers moteurs} hiérarchisés : la Triade historique (Amérique du Nord, Union européenne, Japon élargi à la Corée du Sud), rejointe et bousculée par la Chine, tandis que la Russie joue une carte énergétique et militaire. Ces pôles sont liés par des interdépendances profondes (chaînes de valeur, finance, données) mais fragiles, comme l'ont montré la Covid-19 et la guerre en Ukraine. De nouveaux relais émergent (Inde, ASEAN, Golfe, Afrique), et la gouvernance mondiale se multipolarise (G7, G20, BRICS+, OMC affaiblie). Pour le Cameroun, l'enjeu est clair : échapper au statut de simple fournisseur de matières premières par la transformation locale et l'intégration africaine. Retiens la formule : \emph{la mondialisation a un moteur à plusieurs cylindres — et l'Afrique cherche encore sa place sous le capot}.
\end{aretenir}

\newpage

\coursec{Activité d'intégration}

\courssub{Objectifs de l'activité}

Cette activité d'intégration clôt la leçon sur les foyers moteurs de l'économie mondiale en te plaçant dans une situation professionnelle authentique : conseiller le Ministère des Mines, de l'Industrie et du Développement technologique du Cameroun sur une filière stratégique de la mondialisation, celle des \cle{batteries} pour véhicules électriques. À l'issue de ce dossier guidé de 2 heures, tu devras être capable de :

\begin{itemize}
\item \cle{Inventorier et classer} des documents géographiques variés (carte, tableau statistique, graphique, article de presse, schéma) en précisant leur nature, leur source, leur date et l'information clé qu'ils livrent ;
\item \cle{Observer, localiser et décrire} la répartition mondiale des ressources critiques (lithium, cobalt, nickel, manganèse, graphite) et des gigafactories ;
\item \cle{Interpréter} des données chiffrées (subventions publiques, parts de marché) pour comparer les stratégies des cinq foyers étudiés : Union européenne, Amérique du Nord, Asie de l'Est (Chine, Japon, Corée du Sud) et Russie ;
\item \cle{Représenter et schématiser} une chaîne de valeur et les flux qui la structurent ;
\item \cle{Argumenter} par écrit une proposition de politique publique pour le Cameroun, en mobilisant le vocabulaire de la mondialisation (chaîne de valeur, valeur ajoutée, dépendance, interdépendance, transition énergétique).
\end{itemize}

\courssub{Situation}

\textbf{Communiqué du Cabinet du Ministre des Mines, Yaoundé.} En mars 2024, le gouvernement camerounais reçoit une délégation d'investisseurs chinois et européens intéressés par les gisements de \cle{cobalt}, de \cle{nickel} et de \cle{manganèse} du pays (gisements de Nkamouna et Mada dans la région de l'Est, indices de nickel-cobalt à Lomie). La demande mondiale de batteries lithium-ion explose : la part des véhicules électriques dans les ventes automobiles mondiales est passée d'environ 4 \% en 2020 à environ 18 \% en 2023 (ordre de grandeur, source AIE). Or le Cameroun, comme la plupart des pays africains, exporte ses minerais \emph{bruts} et capte moins de 10 \% de la valeur ajoutée de la filière. Le Ministre constitue un groupe d'experts — dont tu fais partie — pour rédiger une note stratégique. Il met à ta disposition un dossier de six documents.

\textbf{Document 1 — Carte mondiale des ressources critiques et du raffinage (2023).} Le lithium est extrait surtout en Australie (environ la moitié de la production mondiale), au Chili et en Chine ; le cobalt en République démocratique du Congo (environ 70 \% de la production mondiale) ; le nickel en Indonésie, aux Philippines et en Russie ; le graphite naturel et le manganèse en Chine et en Afrique australe. Surtout, la carte montre que la \cle{Chine} assure environ 60 à 75 \% du \emph{raffinage} mondial de ces métaux et près de 90 \% de la production d'anodes en graphite.

\textbf{Document 2 — Tableau des subventions publiques à la filière batteries (annonces 2020-2030, ordres de grandeur).}

\begin{center}
\begin{tabularx}{\linewidth}{|l|X|X|}\hline
\rowcolor{popblueL} \textbf{Foyer} & \textbf{Instrument} & \textbf{Montant annoncé} \\ \hline
États-Unis & Inflation Reduction Act (IRA, 2022) : crédits d'impôt par kWh produit localement & environ 370 milliards de dollars pour le climat, dont une part majeure batteries \\ \hline
Union européenne & Pacte vert (Green Deal), Important Projects of Common European Interest & plusieurs dizaines de milliards d'euros ; objectif : quasi-autosuffisance en cellules vers 2030 \\ \hline
Chine & Plans quinquennaux, subventions à l'achat de véhicules électriques & plus de 100 milliards de dollars cumulés depuis 2009 (estimations) \\ \hline
Japon et Corée du Sud & Plans nationaux batteries, soutien à Panasonic, LG, Samsung SDI, SK On & plusieurs dizaines de milliards de dollars \\ \hline
Russie & Stratégie énergétique : rente hydrocarbures, projets gigafactories à Kaliningrad et Moscou & montants modestes (quelques milliards de dollars) \\ \hline
\end{tabularx}
\end{center}

\textbf{Document 3 — Parts de marché mondiales des cellules de batteries (2023, ordres de grandeur).} CATL (Chine) : environ 35 \% ; BYD (Chine) : environ 16 \% ; LG Energy Solution (Corée du Sud) : environ 14 \% ; Panasonic (Japon) : environ 6 \% ; Samsung SDI et SK On (Corée du Sud) : environ 9 \% cumulés ; Northvolt (Suède, UE) : moins de 2 \%. Les entreprises chinoises et sud-coréennes contrôlent donc à elles seules environ 80 \% du marché mondial des cellules.

\textbf{Document 4 — Extrait de presse (Investir au Cameroun, 2024).} « Des élus de la région de l'Est et des organisations de la société civile plaident pour que le cobalt et le nickel de Nkamouna soient transformés \emph{sur place} en précurseurs de cathodes, au lieu d'être exportés bruts vers la Chine. Ils rappellent que la RDC, premier producteur mondial de cobalt, ne raffine qu'une faible part de sa production et capte moins de 3 \% de la valeur finale d'une batterie. »

\textbf{Document 5 — Carte des flux de composants vers les gigafactories.} Les minerais africains et australiens partent en majorité vers les raffineries chinoises ; les cellules et composants alimentent ensuite les gigafactories de Chine (plus de 200 sites), d'Amérique du Nord (Tesla au Nevada et au Texas, usines LG et SK au Michigan et en Géorgie) et d'Europe (Northvolt en Suède, ACC en France, gigafactory Tesla près de Berlin).

\textbf{Document 6 — Schéma de la chaîne de valeur d'une batterie.} Extraction du minerai $\rightarrow$ raffinage métallurgique $\rightarrow$ production de précurseurs et de matériaux actifs (cathodes, anodes) $\rightarrow$ assemblage des cellules $\rightarrow$ intégration dans les véhicules $\rightarrow$ recyclage. La valeur ajoutée croît à chaque étape : l'extraction ne représente qu'environ 10 à 15 \% de la valeur finale, le raffinage et les matériaux actifs environ 40 \%, l'assemblage et l'intégration le reste.

\courssub{Consignes}

Travail en groupes de quatre, puis mise en commun collective, puis production écrite individuelle.

\begin{enumerate}
\item \textbf{Inventaire des documents (20 min).} Pour chacun des six documents, précise sa nature (carte, tableau, graphique, texte, schéma), sa source ou son origine probable, sa date et l'information essentielle qu'il apporte. Classe ensuite les documents en deux catégories : ceux qui décrivent la \emph{répartition spatiale} des activités et ceux qui décrivent les \emph{stratégies des acteurs}.
\item \textbf{Analyse d'un foyer (40 min).} Chaque groupe se voit attribuer un foyer (UE, États-Unis/Amérique du Nord, Chine, Japon-Corée du Sud, Russie). Réponds aux questions : (a) Quelle stratégie ce foyer adopte-t-il sur la filière batteries ? Appuie-toi sur au moins deux documents. (b) Quels sont ses atouts et ses faiblesses (ressources, technologies, marché, capitaux) ? (c) Quelle place cette stratégie laisse-t-elle à l'Afrique et au Cameroun ?
\item \textbf{Mise en commun (30 min).} Construis collectivement au tableau un tableau comparatif des cinq stratégies (colonnes : foyer ; instrument principal ; atouts ; faiblesses ; degré de dépendance extérieure).
\item \textbf{Cartographie schématique.} Sur ton croquis de synthèse mondiale, reporte : les zones d'extraction des métaux critiques, le pôle chinois de raffinage, et les trois grands bassins de gigafactories, en reliant les flux par des flèches d'épaisseur proportionnelle.
\item \textbf{Rédaction individuelle (30 min, environ 200 mots).} « En tant que conseiller au Ministère des Mines du Cameroun, propose trois mesures prioritaires pour que le pays capte davantage de valeur ajoutée sur la filière batteries. Justifie chaque mesure par au moins une donnée du dossier. »
\end{enumerate}

\courssub{Ressources à mobiliser}

\begin{itemize}
\item Les notions de la leçon : \cle{Triade} (Amérique du Nord, Europe occidentale, Asie de l'Est), foyers moteurs, interdépendances, émergence de la Chine, rente énergétique russe, intégration régionale (UE, USMCA).
\item Les notions transversales : \cle{chaîne de valeur}, \cle{valeur ajoutée}, \cle{ressources critiques}, transition énergétique et numérique, dépendance stratégique.
\item Les savoir-faire : lecture de carte (légende, localisation, flux), calcul de parts de marché et de taux de variation, construction d'un croquis légendé avec flèche du nord et échelle indicative, rédaction d'un paragraphe argumenté.
\item Rappel de calcul : la part de marché cumulée s'obtient par addition des parts ; par exemple Chine + Corée du Sud $\approx 35 + 16 + 14 + 9 = 74$ à $80\,\%$ selon les périmètres retenus.
\end{itemize}

\courssub{Démarche et corrigé commenté}

\begin{exoresolu}[Corrigé guidé du dossier « Batteries »]
Énoncé : à partir des six documents, comparer les stratégies des cinq foyers moteurs sur la filière batteries et formuler trois mesures prioritaires pour le Cameroun.

\tcblower

\textbf{Étape 1 — Inventaire et classement des documents.} Documents 1 et 5 sont des cartes : ils décrivent la \emph{répartition spatiale} (ressources, raffinage, gigafactories, flux). Document 2 est un tableau statistique et document 3 un graphique de parts de marché : ils décrivent les \emph{stratégies et rapports de force} entre acteurs. Document 4 est un article de presse camerounais (source secondaire, à croiser avec les données chiffrées) ; document 6 est un schéma de filière. Erreur fréquente : confondre la nature du document (ce qu'il est) et sa fonction (ce qu'il prouve).

\textbf{Étape 2 — Analyse par foyer.} \emph{Chine} : stratégie de maîtrise complète de la chaîne de valeur, de l'extraction (investissements en RDC et en Indonésie) au raffinage (60 à 75 \% mondial) et aux cellules (CATL + BYD $\approx 51\,\%$ du marché). Atouts : coûts, échelle, avance technologique ; faiblesse : dépendance aux minerais importés. \emph{États-Unis} : stratégie de \cle{réindustrialisation} par le protectionnisme incitatif (IRA, crédits d'impôt conditionnés à la production locale) ; atouts : capitaux, innovation (Tesla), marché ; faiblesse : quasi-absence de raffinage. \emph{Union européenne} : stratégie d'autonomie stratégique par le droit et le marché unique (Green Deal, Northvolt, ACC) ; faiblesse : retard industriel (moins de 2 \% des cellules pour Northvolt). \emph{Japon et Corée du Sud} : champions technologiques (Panasonic, LG, Samsung, SK) mais dépendants à 100 \% des matières premières importées. \emph{Russie} : stratégie de rente énergétique et de pivot eurasiatique ; elle dispose de nickel mais reste marginale sur les cellules, et les sanctions depuis 2022 freinent son intégration.

\textbf{Étape 3 — Tableau comparatif (extrait).} La Chine domine par l'intégration verticale ; les États-Unis et l'UE tentent de rattraper par des subventions massives ; le Japon et la Corée misent sur la technologie ; la Russie reste un fournisseur de matières et d'énergie. Point commun : les cinq stratégies traitent l'Afrique comme un \emph{réservoir de matières premières}, ce que confirme le document 4 (moins de 3 \% de la valeur finale captée par la RDC).

\textbf{Étape 4 — Croquis.} Trois flèches épaisses partent d'Afrique centrale et australe, d'Australie et d'Indonésie vers la Chine (raffinage) ; trois flèches moyennes vont de la Chine vers les bassins de gigafactories nord-américain, européen et est-asiatique ; une flèche fine symbolise le recyclage en boucle. Légende numérotée, nord, titre obligatoires.

\textbf{Étape 5 — Note au Ministre (modèle de réponse).} Mesure 1 : conditionner tout permis d'exploitation de Nkamouna à une \emph{transformation locale} minimale (raffinage ou précurseurs), car le document 6 montre que ces étapes concentrent environ 40 \% de la valeur ajoutée. Mesure 2 : négocier des partenariats diversifiés (UE, Corée, Chine) pour éviter la dépendance à un seul acheteur, à l'image de la concurrence que se livrent les cinq foyers (document 2). Mesure 3 : investir dans la formation technique et l'énergie (barrages) pour rendre l'industrie de transformation compétitive, car sans électricité abondante et compétences, la transformation locale reste un vœu pieux. Conclusion : le Cameroun peut passer du statut de réservoir à celui de maillon de la chaîne de valeur, à condition d'une stratégie d'État cohérente.
\end{exoresolu}

\courssub{Bilan de l'activité}

Ce dossier a permis de vérifier que les cinq foyers moteurs ne se ressemblent pas : la Chine domine par l'intégration de toute la chaîne de valeur, les États-Unis par la puissance financière et technologique, l'Union européenne par le droit et le marché unique, le Japon et la Corée du Sud par l'excellence industrielle, la Russie par sa rente énergétique et minière. Il a surtout montré que la \cle{mondialisation} est un système d'interdépendances hiérarchisées : la transition énergétique, présentée comme un progrès partagé, reproduit une division internationale du travail où l'Afrique fournit les minerais et où la valeur ajoutée se concentre dans la Triade élargie à la Chine. Enfin, l'activité a entraîné l'ensemble des savoir-faire du programme — inventorier, localiser, interpréter des statistiques, construire un croquis, argumenter — et a donné au cours de géographie une portée concrète : celle d'éclairer les choix de développement du Cameroun.

\begin{attention}[Critères d'évaluation]
La grille critériée porte sur quatre points : exploitation rigoureuse des documents (citer le document et la donnée à l'appui de chaque affirmation), vocabulaire géographique exact (chaîne de valeur, valeur ajoutée, Triade, foyer moteur), qualité de l'argumentation (mesures justifiées, pas de liste sans explication), correction du croquis (titre, légende, nord, flèches proportionnelles). Une donnée inventée ou non sourcée fait perdre des points : quand un chiffre est incertain, donne un ordre de grandeur et dis-le.
\end{attention}

\newpage

\coursec{Méthodes et exercices résolus}

\coursec{Cadre conceptuel : la Triade et l'émergence de nouveaux moteurs}

\begin{savaistu}[Situation déclenchante : le Cameroun dans les flux mondiaux]
Imagine un commerçant au marché Central de Douala : il vend des téléphones \textbf{Chinois} (Asie de l'Est), des tissus \textbf{Européens} (UE), boit un café issu de grains \textbf{Brésiliens} transportés par un navire battant pavillon \textbf{Américain} ou \textbf{Asiatique}, et paie en \textbf{Francs CFA} arrimé à l'\textbf{Euro}. Le port de Kribi, en eau profonde, reçoit des porte-conteneurs de 10 000 EVP venus de Shangaï ou de Rotterdam, et exporte du bois, du cacao, du coton vers la Triade. Le Cameroun importe aussi des engrais \textbf{Russes} pour ses plantations (CDC, SOCAPALM). Cette scène quotidienne illustre la place du Cameroun comme \textbf{interface} périphérique reliée aux \textbf{quatre foyers moteurs} qui structurent l'économie mondiale.
\end{savaistu}

\begin{definition}[Triade et Triade élargie]
La \cle{Triade} désigne, au sens strict, les trois pôles historiques dominants de l'économie mondiale : l'\textbf{Amérique du Nord} (États-Unis, Canada, Mexique), l'\textbf{Union européenne} (27 États) et le \textbf{Japon}. La \cle{Triade élargie} intègre la \textbf{Chine} (deuxième PIB mondial, premier exportateur) et les \textbf{Nouveaux Pays Industrialisés} (Corée du Sud, Taïwan, Singapour, Hong Kong) au pôle asiatique. Ces ensembles concentrent environ \textbf{deux tiers du PIB mondial} (PPA), \textbf{75 \% des échanges de marchandises} et l'essentiel des flux d'IDE et d'innovation.
\end{definition}

\begin{propriete}[Localisation des foyers moteurs : une constante nordique]
\begin{enumerate}
\item \textbf{Hémisphère Nord exclusivement :} les quatre foyers (Amérique du Nord, UE, Asie de l'Est, Russie) se situent au nord de l'équateur (entre 30° N et 60° N pour l'essentiel).
\item \textbf{Position littorale et insulaire :} trois foyers sur quatre (Amérique du Nord, UE, Asie de l'Est) sont ouverts sur les deux grands océans (Atlantique, Pacifique) qui supportent 90 \% du commerce maritime mondial. La Russie est un foyer \textbf{continental} (pivot eurasiatique) avec des façades arctiques et pacifiques.
\item \textbf{Concentration côtière :} la richesse et la population se massent sur les façades maritimes (mégalopoles, deltas, ports) : \emph{BosWash} (USA), \emph{dorsale européenne} (UE), \emph{façade pacifique} (Chine, Japon, Corée).
\end{enumerate}
\end{propriete}

\courssub{Méthode 1 : Observer, localiser et décrire un foyer moteur sur une carte}

\begin{methode}[Décrire la localisation d'un foyer moteur]
Pour décrire la localisation d'un foyer de l'économie mondiale (UE, Amérique du Nord, Asie de l'Est, Russie), suis toujours la même démarche en quatre étapes :
\begin{enumerate}
\item \cle{Situer} le foyer dans le monde : hémisphère (Nord pour tous les foyers moteurs), continent, position par rapport aux mers et aux grandes routes maritimes.
\item \cle{Délimiter} : nommer les États ou ensembles régionaux concernés (ex. : les 27 États membres de l'UE ; les 3 États de l'USMCA).
\item \cle{Identifier les points de concentration} : mégalopoles, littorales, façades portuaires (ex. : mégalopole américaine BosWash, façade pacifique chinoise, dorsale européenne).
\item \cle{Conclure par l'essentiel} : une phrase qui relie la localisation à la puissance (position d'interface, ouverture maritime, densité urbaine).
\end{enumerate}
Réflexe : une description va toujours \textbf{du général au particulier} (monde $\rightarrow$ région $\rightarrow$ pôles précis) et se termine par une phrase d'interprétation.
\end{methode}

\begin{exoresolu}[Localiser et décrire la Triade]
\textbf{Énoncé (type Bac).} À partir de vos connaissances et d'une carte mondiale des principaux foyers économiques : (a) définis la Triade ; (b) décris la localisation de ses trois pôles ; (c) explique en une phrase pourquoi on parle de « Triade élargie ».
\tcblower
\textbf{Solution détaillée.}

(a) \emph{Définition.} La \cle{Triade} désigne les trois pôles dominants de l'économie mondiale : l'\textbf{Amérique du Nord} (États-Unis, Canada, Mexique), l'\textbf{Union européenne} et l'\textbf{Asie de l'Est} (Japon, Chine, Corée du Sud). Ces trois ensembles concentrent l'essentiel des richesses produites, des échanges, des flux d'investissement et de l'innovation mondiale (ordre de grandeur : environ les deux tiers du PIB mondial, les parts exactes variant selon les années et les sources).

(b) \emph{Localisation (du général au particulier).}
\begin{itemize}
\item Les trois pôles sont tous situés dans l'\textbf{hémisphère Nord}, en position \textbf{littorale}, ouverts sur les deux grands océans (Atlantique et Pacifique) que parcourent les principales routes maritimes mondiales.
\item L'\textbf{Amérique du Nord} s'étend entre l'Atlantique et le Pacifique ; sa richesse se concentre dans la mégalopole du Nord-Est américain (\emph{BosWash}, de Boston à Washington), sur les Grands Lacs et sur la façade californienne.
\item L'\textbf{Union européenne} occupe l'ouest du continent européen ; son cœur économique est la \emph{dorsale européenne} (ou « banane bleue »), qui s'étire des Midlands britanniques à l'Italie du Nord en passant par le Benelux, la vallée du Rhin et la Suisse.
\item L'\textbf{Asie de l'Est} borde le Pacifique occidental ; la production s'y concentre sur les littoraux : façade pacifique du Japon (de Tokyo à Osaka), littoral chinois (Shangaï, delta du Yangzi, delta de la Rivière des Perles avec Shenzhen et Canton), Séoul et sa région en Corée du Sud.
\end{itemize}

(c) \emph{Triade élargie.} On parle de « Triade élargie » parce que la \textbf{Chine}, absente de la Triade initiale (qui ne comptait que le Japon en Asie), est devenue la deuxième puissance économique mondiale et le premier exportateur de marchandises : le pôle asiatique ne se réduit plus au Japon.

\emph{Conclusion.} La Triade structure l'économie mondiale autour de trois façades maritimes de l'hémisphère Nord, véritables interfaces de la mondialisation.
\end{exoresolu}

\courssub{Méthode 3 : Inventorier, classer et interpréter un tableau statistique}

\begin{methode}[Exploiter un tableau de données économiques]
Devant un tableau statistique (PIB, parts dans le commerce mondial, population) :
\begin{enumerate}
\item \cle{Lire le titre complet} : quoi (l'indicateur), où (les pays), quand (l'année), dans quelle unité (milliards de dollars, \%).
\item \cle{Classer} les données : repérer le maximum, le minimum, et ordonner mentalement les valeurs.
\item \cle{Calculer} pour transformer la lecture : part en pourcentage $\text{part} = \dfrac{\text{valeur de la partie}}{\text{valeur totale}} \times 100$ ; taux de variation $\text{TV} = \dfrac{V_f - V_i}{V_i} \times 100$ ; rapport entre extrêmes (écart multiplicatif).
\item \cle{Interpréter} : dégager une ou deux idées forces (hiérarchie, évolution, concentration) et les relier au sujet — jamais de simple récitation des chiffres.
\end{enumerate}
Réflexe : annonce toujours le chiffre \textbf{avec son unité et son année}, et accompagne chaque calcul du résultat arrondi de façon cohérente (une décimale suffit pour des pourcentages).
\end{methode}

\begin{exoresolu}[Hiérarchie des puissances économiques]
\textbf{Énoncé (type Bac).} Le tableau suivant donne le PIB en parité de pouvoir d'achat (ordres de grandeur indicatifs, estimations FMI/Banque Mondiale 2023-2024) :

\begin{center}
\begin{tabularx}{\linewidth}{|l|X|X|}
\hline
\rowcolor{popblueL} \textbf{Ensemble} & \textbf{PIB (milliards de \$ PPA)} & \textbf{Population (millions)} \\
\hline
États-Unis & \num{25000} & 335 \\
\hline
Chine & \num{27000} & 1410 \\
\hline
Union européenne (27) & \num{21000} & 450 \\
\hline
Japon & \num{6000} & 125 \\
\hline
Russie & \num{4500} & 145 \\
\hline
Corée du Sud & \num{2500} & 52 \\
\hline
\end{tabularx}
\end{center}

(a) Calcule la part de chacun des trois pôles de la Triade élargie (Amérique du Nord assimilée ici aux seuls États-Unis, UE, Asie de l'Est = Chine + Japon + Corée du Sud) dans le total des six ensembles. (b) Calcule le rapport entre le PIB de la Chine et celui de la Russie. (c) Tire deux conclusions géographiques.
\tcblower
\textbf{Solution détaillée.}

\emph{Étape 1 : total des six ensembles.}
\[ 25\,000 + 27\,000 + 21\,000 + 6\,000 + 4\,500 + 2\,500 = 86\,000 \text{ milliards de dollars.} \]

\emph{Étape 2 : parts de la Triade élargie.}
\begin{align*}
\text{Asie de l'Est} &= 27\,000 + 6\,000 + 2\,500 = 35\,500 \text{ Mds \$} \\
\text{Part Asie de l'Est} &= \frac{35\,500}{86\,000} \times 100 \approx 41{,}3\,\% \\
\text{Part États-Unis} &= \frac{25\,000}{86\,000} \times 100 \approx 29{,}1\,\% \\
\text{Part UE} &= \frac{21\,000}{86\,000} \times 100 \approx 24{,}4\,\%
\end{align*}
Vérification : $41{,}3 + 29{,}1 + 24{,}4 = 94{,}8\,\%$ ; le reste (Russie, $4\,500/86\,000 \approx 5{,}2\,\%$) complète bien à $100\,\%$.

\emph{Étape 3 : rapport Chine/Russie.}
\[ \frac{27\,000}{4\,500} = 6. \]
Le PIB de la Chine est environ \textbf{6 fois} celui de la Russie.

\emph{Étape 4 : conclusions.}
\begin{itemize}
\item La Triade élargie concentre à elle seule près de \textbf{95 \%} de la richesse des six ensembles étudiés : l'économie mondiale reste polarisée par trois grands foyers, mais le centre de gravité se déplace vers l'\textbf{Asie de l'Est}, premier pôle devant les États-Unis.
\item La \textbf{Russie} pèse économiquement peu (environ 5 \%) malgré son immensité territoriale : sa qualité de foyer moteur repose moins sur son PIB que sur ses \textbf{ressources énergétiques} (gaz, pétrole) et sa \textbf{puissance militaire}, ce qui en fait un foyer d'un type particulier.
\end{itemize}

\emph{Phrase de conclusion.} Les chiffres confirment la hiérarchie mondiale : une Triade élargie dominée désormais par l'Asie de l'Est, tandis que la Russie compense sa faiblesse économique par son rôle de pivot énergétique et stratégique.
\end{exoresolu}

\coursec{L'Union européenne : le premier foyer intégré par le droit et le marché unique}

\begin{definition}[Union européenne (UE) et Marché unique]
L'\cle{Union européenne} est une union politico-économique de \textbf{27 États membres} (depuis le Brexit, 2020) peuplée de \textbf{450 millions d'habitants}. Son cœur est le \cle{Marché unique} (libre circulation des biens, services, capitaux, personnes) et l'\cle{Espace Schengen} (suppression des contrôles aux frontières intérieures pour 29 pays). \textbf{20 États} partagent la \cle{monnaie unique}, l'\textbf{euro} (zone euro), gérée par la \textbf{BCE} (Banque centrale européenne, siège à Francfort).
\end{definition}

\begin{propriete}[Forces structurelles de l'UE]
\begin{itemize}
\item \textbf{Puissance commerciale :} premier ou deuxième ensemble commercial mondial (selon indicateur), excédent structurel dans les services et l'industrie de haute technologie.
\item \textbf{Puissance normative :} « Bruxelles effect » — les normes européennes (environnement, numérique RGPD, concurrence) deviennent des standards mondiaux.
\item \textbf{Intégration par le droit :} institutions supranationales (Commission, Parlement, Cour de justice) : les États partagent leur souveraineté, unique au monde.
\item \textbf{Réseau urbain de rang mondial :} Londres (finance, hors UE mais connectée), Paris, Francfort (BCE), Milan, Amsterdam, Madrid — dorsale européenne = « banane bleue » (60 millions d'habitants, PIB dense).
\end{itemize}
\end{propriete}

\begin{exemplebox}[Exemples d'acteurs européens mondiaux]
\textbf{Airbus} (aéronautique, concurrent de Boeing), \textbf{Volkswagen / Stellantis} (automobile), \textbf{TotalEnergies} (énergie), \textbf{LVMH} (luxe), \textbf{ASML} (Pays-Bas, machines pour semi-conducteurs, monopole mondial EUV), \textbf{SAP} (logiciels), \textbf{Sanofi} (pharma).
\end{exemplebox}

\begin{aretenir}[Fractures et limites de l'UE]
\begin{enumerate}
\item \textbf{Disparités Nord/Sud et Est/Ouest :} PIB/hab. Luxembourg (130 000 \$) vs Bulgarie (28 000 \$) — écart 1 à 4,5. Fonds structurels (FEDER, FSE) et Plan de relance post-Covid (750 Mds €) tentent de réduire l'écart.
\item \textbf{Dépendance énergétique :} brutale prise de conscience en 2022 (gaz russe) ; diversification accélérée (GNL USA, Norvège, renouvelables).
\item \textbf{Vieillissement démographique :} taux de fécondité ~1,5 ; besoin d'immigration choisie ; coût des retraites.
\item \textbf{Fragmentation politique :} montée des euroscepticismes, règle de l'unanimité en politique étrangère/fiscale (blocages).
\end{enumerate}
\end{aretenir}

\courssub{Méthode 5 : Rédiger un commentaire de document (carte ou texte) sur un foyer moteur}

\begin{methode}[Commenter un document sur les foyers moteurs]
\begin{enumerate}
\item \cle{Présenter le document} en deux lignes : nature (carte, texte, tableau), thème, auteur/source, date.
\item \cle{Dégager la problématique} : quelle question le document permet-il de traiter ? (ex. : « En quoi l'UE est-elle un foyer intégré mais inégalement développé ? »).
\item \cle{Analyser en deux ou trois parties} annoncées : pour chaque partie, partir du document (citer un élément précis : un chiffre, un lieu, une flèche) puis l'expliquer par tes \textbf{connaissances} (marché unique, USMCA, atelier du monde, hydrocarbures russes).
\item \cle{Conclure} en répondant à la problématique et en ouvrant brièvement (limites du foyer, défis : vieillissement au Japon, tensions commerciales, guerre en Ukraine).
\end{enumerate}
Réflexe : alterne toujours « je vois » (le document) et « je sais » (le cours) ; un commentaire sans aucune connaissance personnelle est sanctionné, tout comme une paraphrase du document.
\end{methode}

\begin{exoresolu}[Commentaire : l'Union européenne, premier foyer intégré]
\textbf{Énoncé (type Bac).} « Fondée sur un marché unique de 450 millions de consommateurs, dotée d'une monnaie commune partagée par 20 États (zone euro) et d'institutions supranationales (Commission, Parlement, BCE), l'Union européenne est le premier foyer intégré du monde par le droit. Mais la richesse reste concentrée dans la dorsale européenne, des Midlands à l'Italie du Nord, tandis que les périphéries orientales et méridionales accusent un retard de développement. » Commentez ce texte en montrant les forces et les limites de l'UE comme foyer moteur.
\tcblower
\textbf{Solution détaillée (plan et éléments de rédaction).}

\emph{Présentation.} Il s'agit d'un texte de synthèse géographique portant sur l'Union européenne, qui pose la question suivante : en quoi l'intégration fait-elle de l'UE un foyer majeur, et quelles fractures internes limitent cette puissance ?

\emph{I. Un foyer moteur fondé sur une intégration inédite.} Le texte indique un « marché unique de 450 millions de consommateurs » : c'est le résultat de la construction européenne (depuis le traité de Rome, 1957), qui garantit la libre circulation des biens, des capitaux, des services et des personnes (espace Schengen). La monnaie unique, l'\textbf{euro} (20 États de la zone euro), élimine les risques de change et affirme la puissance monétaire de l'UE, deuxième monnaie de réserve mondiale derrière le dollar. Les institutions supranationales (Commission européenne, Parlement, Banque centrale européenne à Francfort) font de l'UE le seul foyer intégré « par le droit » : les États ont accepté de partager une partie de leur souveraineté. L'UE est ainsi le premier ou deuxième ensemble commercial mondial selon les indicateurs, avec des pôles de rang mondial (Londres, Paris, Francfort, Milan) et de grandes firmes (Airbus, Volkswagen, TotalEnergies).

\emph{II. Un foyer aux disparités internes marquées.} Le texte souligne la concentration de la richesse dans la « dorsale européenne » : cette mégalopole européenne (ou « banane bleue ») cumule densités de population, industries de pointe et services supérieurs. À l'inverse, les périphéries — Europe de l'Est (Pologne, Roumanie, Bulgarie, intégrées après 2004) et du Sud (Grèce, sud de l'Italie) — présentent des PIB par habitant inférieurs, malgré les fonds structurels européens destinés à réduire ces écarts. S'ajoutent d'autres limites non mentionnées mais connues : dépendance énergétique (au gaz russe, brutalement révélée par la guerre en Ukraine en 2022), vieillissement démographique, concurrence chinoise.

\emph{Conclusion.} L'UE est bien un foyer moteur de premier plan, original par son intégration juridique et monétaire, mais sa puissance est inégalement répartie et fragilisée par ses dépendances extérieures — ce qui la distingue de l'hégémonie américaine.
\end{exoresolu}

\coursec{L'Amérique du Nord : la puissance hégémonique (USA) et ses partenaires intégrés (USMCA)}

\begin{definition}[USMCA et hégémonie américaine]
L'\cle{USMCA} (Accord États-Unis–Mexique–Canada, 2020, succédant à l'ALENA) crée une zone de libre-échange de \textbf{500 millions d'habitants} et \textbf{28 000 Mds \$ de PIB}. Les \textbf{États-Unis} en sont le moteur hégémonique : \textbf{première puissance mondiale} (PIB nominal ~27 000 Mds \$, dollar = 58 \% des réserves mondiales, première puissance militaire, berceau de l'innovation numérique).
\end{definition}

\begin{propriete}[Les trois piliers de la puissance nord-américaine]
\begin{itemize}
\item \textbf{Financière et monétaire :} Wall Street, \textbf{dollar} (monnaie de facturation du pétrole, des matières premières, dette refuge), FED (banque centrale mondiale de fait).
\item \textbf{Technologique et immatérielle :} \textbf{GAFAM} (Google/Alphabet, Apple, Facebook/Meta, Amazon, Microsoft) + semi-conducteurs (Nvidia, Intel), biotech, IA. \textbf{Silicon Valley} (Californie) = cœur mondial de l'innovation.
\item \textbf{Intégration régionale asymétrique :} Canada (matières premières, énergie, services) et Mexique (\emph{maquiladoras}, industrie automobile/électronique, main-d'œuvre compétitive) sont profondément imbriqués dans les chaînes de valeur US (75 \% des exportations mexicaines vont aux USA).
\end{itemize}
\end{propriete}

\begin{aretenir}[Faiblesses et défis]
\textbf{Désindustrialisation partielle} (Rust Belt), \textbf{déficit commercial massif} (surtout avec la Chine, ~ -1 000 Mds \$/an), \textbf{inégalités sociales} records (Gini > 0,48), \textbf{polarisation politique} freinant les investissements d'infrastructure, \textbf{endettement public} > 120 \% PIB.
\end{aretenir}

\coursec{L'Asie de l'Est : l'atelier du monde et l'innovation (Chine, Japon, Corée du Sud)}

\begin{definition}[Asie de l'Est : Chine, Japon, Corée du Sud]
L'\cle{Asie de l'Est} (au sens du programme) regroupe la \textbf{Chine} (1,41 Md hab., PIB PPA ~27 000 Mds \$), le \textbf{Japon} (125 M hab., PIB PPA ~6 000 Mds \$) et la \textbf{Corée du Sud} (52 M hab., PIB PPA ~2 500 Mds \$). C'est le \textbf{premier pôle économique mondial} en PPA (41 \% de l'échantillon Triade+Russie), le \textbf{premier exportateur de marchandises} et le cœur des \textbf{chaînes de valeur mondiales} (électronique, textile, automobile, batteries).
\end{definition}

\begin{propriete}[Modèles de croissance différenciés mais complémentaires]
\begin{itemize}
\item \textbf{Chine :} « Atelier du monde » $\rightarrow$ montée en gamme (\emph{Made in China 2025}) : 1er exportateur high-tech (Huawei, BYD véhicules électriques, CATL batteries), \textbf{Nouvelles Routes de la Soie (BRI)} : investissements portuaires/ferroviaires en Afrique (Kribi, Lagos, Mombasa), Asie, Europe.
\item \textbf{Japon :} Puissance \textbf{technologique mature} (Toyota, Sony, Keyence, robotique, composants de pointe), excédents commerciaux structurels, \textbf{vieillissement extrême} (médiane 49 ans), dette publique > 260 \% PIB.
\item \textbf{Corée du Sud :} Modèle \textbf{Chaebols} (Samsung, Hyundai, SK Hynix, LG) — champions mondiaux puces mémoire, smartphones, navale, batteries. Forte dépendance aux exportations (70 \% PIB), vieillissement le plus rapide du monde.
\end{itemize}
\end{propriete}

\begin{aretenir}[Intégration régionale et tensions]
\textbf{RCEP} (Partenariat économique régional global, 2022) : plus grand accord de libre-échange du monde (15 pays, Chine, Japon, Corée, ASEAN, Australie, NZ) — réduit les droits de douane, harmonise les règles d'origine. \textbf{Tensions géopolitiques :} Mer de Chine méridionale, Taïwan, péninsule coréenne, rivalité technologique US/Chine (puces, IA).
\end{aretenir}

\coursec{La Russie : foyer énergétique, militaire et pivot eurasiatique}

\begin{definition}[La Russie : une puissance asymétrique]
La \cle{Fédération de Russie} (145 M hab., PIB PPA ~4 500 Mds \$, ~6 fois < Chine) est un foyer moteur \textbf{sui generis} : sa puissance ne repose pas sur la production de richesse manufacturière ou de services, mais sur la \textbf{rente énergétique} et la \textbf{dissuasion nucléaire}.
\end{definition}

\begin{propriete}[Les trois piliers de la puissance russe]
\begin{enumerate}
\item \textbf{Énergétique :} 1er ou 2e exportateur mondial de \textbf{gaz naturel} (Gazprom, gazoducs \emph{Nord Stream} — sabotés 2022 —, \emph{TurkStream}, \emph{Force de Sibérie} vers Chine), grand exportateur de \textbf{pétrole} (Rosneft, Lukoil) et de produits raffinés. Pivot vers l'Asie (Chine, Inde) depuis sanctions 2022.
\item \textbf{Militaire et stratégique :} 1er arsenal nucléaire mondial (6 000 ogives), siège permanent CSNU, projection en Syrie, Afrique (Wagner/Africa Corps), Arctique.
\item \textbf{Géographique :} Territoire immense (17 M km²), ressources minières (diamants, métaux rares, uranium), position de \textbf{pivot eurasiatique} (UEEA, OCS, BRICS+).
\end{enumerate}
\end{propriete}

\begin{aretenir}[Faiblesses structurelles]
Économie \textbf{peu diversifiée} (pétrole/gaz = 30-40 \% PIB, 60 \% exportations), \textbf{démographie déclinante} (natalité basse, espérance de vie hommes ~67 ans), \textbf{sanctions occidentales} (gel réserves BC, exclusion SWIFT, embargo technologique), \textbf{fuite des cerveaux} et capitaux. Dépendance croissante vis-à-vis de la Chine (asymétrie commerciale).
\end{aretenir}

\coursec{Synthèse : Hiérarchie, interdépendances et croquis de synthèse mondiale}

\begin{aretenir}[Hiérarchie des quatre foyers moteurs (synthèse)]
\begin{center}
\begin{tabularx}{\linewidth}{|l|X|X|X|}
\hline
\rowcolor{popblueL} \textbf{Foyer} & \textbf{Nature dominante de la puissance} & \textbf{Moteur principal} & \textbf{Faiblesse majeure} \\
\hline
\textbf{USA (Amérique du Nord)} & Hégémonique (globale) & Finance, Dollar, Innovation (IA, Numérique), Armée & Déficits jumeaux, inégalités, polarisation \\
\hline
\textbf{UE} & Normative, Commerciale, Juridique & Marché unique, Euro, Normes, Haute technologie & Fragmentation, Démographie, Dépendance énergie \\
\hline
\textbf{Asie de l'Est (Chine, Japon, Corée)} & Productive, Commerciale, Technologique & Usine monde, Exportations, Innovation (puces, EV, Robotique) & Vieillissement, Dépendance énergie/tech US, Tensions géopolitiques \\
\hline
\textbf{Russie} & Énergétique, Militaire, Géographique & Gaz/Pétrole, Armement nucléaire, Territoire & Économie mono-produit, Démographie, Sanctions, Isolement relatif \\
\hline
\end{tabularx}
\end{center}
\end{aretenir}

\courssub{Méthode 2 : Construire et légender un croquis de synthèse mondiale}

\begin{methode}[Réussir un croquis « foyers moteurs de l'économie mondiale »]
Le croquis est l'exercice reine du Bac de géographie. Pour un croquis mondial des foyers moteurs :
\begin{enumerate}
\item \cle{Tracer le fond de carte} : contours très simplifiés des continents en blocs (pas de précision inutile), avec l'équateur et éventuellement le tropique du Cancer en pointillés.
\item \cle{Placer les foyers} par zones colorées ou hachurées : Amérique du Nord, UE, Asie de l'Est, Russie. Ajoute les pôles précis par des points nommés (New York, Los Angeles, Londres, Francfort, Tokyo, Shangaï, Séoul, Moscou).
\item \cle{Figurer les flux} par des flèches d'épaisseur proportionnelle : routes maritimes transpacifiques et transatlantiques, flux d'hydrocarbures de la Russie vers l'Europe et la Chine.
\item \cle{Soigner la présentation obligatoire} : un \textbf{titre} précis, une \textbf{légende numérotée} organisée (de l'espace le plus intégré au moins intégré), une \textbf{flèche du nord} et une échelle indicative.
\end{enumerate}
Réflexe : la légende doit être \textbf{hiérarchisée} (2 à 4 rubriques maximum) et chaque élément de la légende doit apparaître sur la carte, et réciproquement.
\end{methode}

\begin{exoresolu}[Croquis : les foyers moteurs et leurs échanges]
\textbf{Énoncé (type Bac).} Construis un croquis intitulé « Les foyers moteurs de l'économie mondiale : puissance et interdépendances ». Tu y figureras les quatre foyers étudiés, leurs principaux pôles urbains et les grands flux qui les relient.
\tcblower
\textbf{Solution détaillée.} Le croquis attendu prend la forme suivante :

\begin{popfigure}
\includegraphics[width=\linewidth,height=9.5cm,keepaspectratio]{N13/dev_emerging.png}
\legende{Fond de carte pour le croquis « Les foyers moteurs de l'économie mondiale » : rose = marchés développés (Amérique du Nord, Europe, Japon, Australie), bleu = marchés émergents (Chine, Inde, Russie, Brésil, Mexique, Afrique du Sud, etc.), gris = autres pays. \textbf{Le croquis attendu situe :} 1. les quatre foyers (Amérique du Nord et USMCA, Union européenne, Asie de l'Est, Russie) ; 2. les principaux pôles urbains et économiques (points noirs) ; 3. les grands flux maritimes de marchandises (transatlantique, transpacifique, Europe--Asie) ; 4. les flux d'hydrocarbures russes vers l'Europe et la Chine (flèches violettes discontinues). \\ Source : Wikimedia Commons, « Developed and Emerging markets », AlexCovarrubias, domaine public.}
\end{popfigure}

\emph{Étapes de construction.} (1) Fond de carte en blocs continentaux avec l'équateur en repère. (2) Quatre foyers colorés, tous dans l'hémisphère Nord : on remarque qu'aucun foyer moteur n'est au sud de l'équateur, ce qui illustre la concentration de la puissance au Nord. (3) Pôles urbains placés sur les littoraux (New York, Shangaï, Tokyo), sauf Moscou, foyer continental. (4) Flux : les deux artères majeures sont la route transpacifique (Asie de l'Est $\rightarrow$ Amérique du Nord, la plus chargée du monde) et la route transatlantique ; les flux énergétiques russes, figurés en pointillés, montrent la dépendance de l'Europe au gaz russe.

\emph{Conclusion.} Le croquis met en évidence une économie mondiale polarisée par quatre foyers nordiques, reliés par des flux maritimes et énergétiques intenses : puissance et interdépendance vont de pair.
\end{exoresolu}

\courssub{Méthode 4 : Construire un diagramme statistique (barres, circulaire)}

\begin{methode}[Choisir et construire le bon diagramme]
\begin{enumerate}
\item \cle{Choisir le type de diagramme} adapté à la question : \textbf{barres} pour comparer des valeurs entre elles (PIB des puissances) ; \textbf{circulaire} (camembert) pour montrer des parts d'un total (répartition en \%) ; \textbf{courbe} pour une évolution dans le temps.
\item \cle{Préparer les données} : pour un diagramme circulaire, convertis chaque part en angle : $\text{angle} = \dfrac{\text{part} \times 360}{100}$ (en degrés).
\item \cle{Tracer avec soin} : axes gradués et légendés (grandeur + unité) pour les barres ; secteurs proportionnels pour le circulaire.
\item \cle{Finaliser} : titre précis (quoi, où, quand), légende, source éventuelle.
\end{enumerate}
Réflexe : vérifie que la somme des angles d'un diagramme circulaire fait bien $360°$ et que les hauteurs de barres sont proportionnelles aux valeurs.
\end{methode}

\begin{exoresolu}[Diagramme circulaire : les parts de la Triade élargie]
\textbf{Énoncé (type Bac).} À partir des parts calculées dans l'exercice précédent (Asie de l'Est : 41,3 \% ; États-Unis : 29,1 \% ; UE : 24,4 \% ; Russie : 5,2 \%), construis un diagramme circulaire et commente-le en trois lignes.
\tcblower
\textbf{Solution détaillée.}

\emph{Étape 1 : conversion des pourcentages en angles.}
\begin{align*}
\text{Asie de l'Est} &: \frac{41{,}3 \times 360}{100} \approx 149° \\
\text{États-Unis} &: \frac{29{,}1 \times 360}{100} \approx 105° \\
\text{UE} &: \frac{24{,}4 \times 360}{100} \approx 88° \\
\text{Russie} &: \frac{5{,}2 \times 360}{100} \approx 19°
\end{align*}
Vérification : $149 + 105 + 88 + 19 = 361° \approx 360°$ (écart dû aux arrondis, acceptable).

\emph{Étape 2 : tracé.}

\begin{popfigure}
\begin{tikzpicture}[scale=1.1]
% Secteurs : Asie 149°, USA 105°, UE 88°, Russie 19° (cumuls : 149, 254, 342, 361)
\fill[popgreen!70] (0,0) -- (90:2) arc (90:-59:2) -- cycle;
\fill[popblue!70] (0,0) -- (-59:2) arc (-59:-164:2) -- cycle;
\fill[poporange!80] (0,0) -- (-164:2) arc (-164:-252:2) -- cycle;
\fill[poppurple!60] (0,0) -- (-252:2) arc (-252:-271:2) -- cycle;
\draw[popink] (0,0) circle (2);
% Étiquettes
\node[font=\scriptsize] at (15:1.3) {\textbf{Asie de l'Est}};
\node[font=\scriptsize] at (15:0.85) {41,3 \%};
\node[font=\scriptsize] at (-111:1.35) {\textbf{États-Unis}};
\node[font=\scriptsize] at (-111:0.9) {29,1 \%};
\node[font=\scriptsize] at (152:1.35) {\textbf{UE}};
\node[font=\scriptsize] at (152:0.95) {24,4 \%};
\node[font=\scriptsize, right] at (2.15,-1.55) {Russie : 5,2 \%};
\draw[popmuted, thin] (-261.5:2) -- (2.1,-1.5);
\end{tikzpicture}
\legende{Diagramme circulaire : répartition du PIB (PPA) des six grands ensembles étudiés, début des années 2020 (ordres de grandeur indicatifs). Secteurs : vert = Asie de l'Est ; bleu = États-Unis ; orange = Union européenne ; violet = Russie.}
\end{popfigure}

\emph{Étape 3 : commentaire (trois lignes).} Le diagramme montre l'écrasante domination de la Triade élargie, qui réunit près de 95 \% du total. L'Asie de l'Est forme le secteur le plus large (près de la moitié du cercle), signe du basculement de l'économie mondiale vers le Pacifique. La Russie n'apparaît que comme un secteur étroit, confirmant que son statut de foyer repose sur l'énergie et l'armement plus que sur la richesse produite.
\end{exoresolu}

\courssub{Méthode 6 : Mener une étude de cas comparée (Asie de l'Est, Amérique du Nord, Russie)}

\begin{methode}[Comparer les foyers moteurs sans les confondre]
\begin{enumerate}
\item \cle{Définir les critères de comparaison} avant de rédiger : nature de la puissance (productive, commerciale, financière, énergétique, militaire), moteurs de croissance, faiblesses.
\item \cle{Traiter chaque foyer avec les mêmes critères} : c'est ce qui rend la comparaison valide.
\item \cle{Mobiliser les exemples précis} : États-Unis = dollar, GAFAM, première puissance financière et militaire, intégration USMCA avec Canada et Mexique ; Chine = « atelier du monde », premier exportateur, nouvelles routes de la soie ; Japon et Corée du Sud = innovation (Toyota, Samsung, Sony) mais vieillissement ; Russie = gaz et pétrole (Gazprom), arsenal militaire, pivot eurasiatique.
\item \cle{Conclure par une hiérarchie nuancée} : puissance globale (États-Unis), puissance productive et commerciale ascendante (Chine), puissance technologique (Japon, Corée), puissance énergétique et stratégique (Russie).
\end{enumerate}
Réflexe : une comparaison peut s'appuyer sur un tableau à double entrée (critères en lignes, foyers en colonnes) avant la rédaction.
\end{methode}

\begin{exoresolu}[Étude de cas : quelle nature de puissance pour chaque foyer ?]
\textbf{Énoncé (type Bac).} « Les foyers moteurs de l'économie mondiale ne sont pas puissants de la même manière. » Illustre cette affirmation en comparant l'Amérique du Nord, l'Asie de l'Est et la Russie, à l'aide d'exemples précis.
\tcblower
\textbf{Solution détaillée.}

\emph{Introduction.} Si la puissance se mesure d'abord au PIB, elle repose en réalité sur des fondements différents selon les foyers : financière et technologique en Amérique du Nord, productive et commerciale en Asie de l'Est, énergétique et militaire pour la Russie.

\emph{1. L'Amérique du Nord : la puissance hégémonique.} Les États-Unis cumulent tous les attributs : première puissance financière (Wall Street, rôle central du \textbf{dollar} dans les échanges et les réserves mondiales), suprématie technologique (les GAFAM : Google, Apple, Facebook/Meta, Amazon, Microsoft ; la Silicon Valley), première puissance militaire. L'accord \textbf{USMCA} (succédant à l'ALENA en 2020) intègre le Canada (matières premières) et le Mexique (main-d'œuvre industrielle, maquiladoras) dans une économie de plus de 500 millions d'habitants. Faiblesse : désindustrialisation partielle et déficit commercial massif avec la Chine.

\emph{2. L'Asie de l'Est : la puissance productive et l'innovation.} La \textbf{Chine} est devenue « l'atelier du monde » : premier exportateur mondial de marchandises, usine des produits électroniques et textiles, elle monte désormais en gamme (Huawei, véhicules électriques BYD) et projette sa puissance par les \textbf{nouvelles routes de la soie} (Belt and Road Initiative). Le \textbf{Japon} (Toyota, Sony) et la \textbf{Corée du Sud} (Samsung, Hyundai) sont des champions de l'innovation et de la robotique. Faiblesses : dépendance énergétique, vieillissement démographique aigu (Japon, Corée), tensions géopolitiques (mer de Chine, péninsule coréenne).

\emph{3. La Russie : la puissance énergétique et stratégique.} Avec un PIB modeste (environ 6 fois inférieur à celui de la Chine, d'après l'exercice statistique), la Russie est un foyer d'un type particulier : premier ou deuxième exportateur mondial de gaz naturel (Gazprom) et grand exportateur de pétrole, elle alimente l'Europe et de plus en plus la Chine ; elle dispose du premier arsenal nucléaire mondial et d'un territoire eurasiatique immense, pivot entre l'Europe et l'Asie. Faiblesses : économie peu diversifiée, démographie déclinante, sanctions occidentales depuis 2022.

\emph{Conclusion.} La puissance mondiale est donc plurielle : hégémonique pour les États-Unis, manufacturière et ascendante pour l'Asie de l'Est, énergétique et militaire pour la Russie. Cette complémentarité explique les interdépendances — et les rivalités — qui structurent la mondialisation.
\end{exoresolu}

\begin{attention}[Les 5 erreurs qui coûtent des points au Bac]
\begin{enumerate}
\item \textbf{Confondre les ensembles.} Ne confonds jamais \cle{Triade} (Amérique du Nord, UE, Asie de l'Est), \cle{USMCA} (États-Unis, Mexique, Canada) et \cle{Union européenne} (27 États depuis le Brexit). Écrire « l'UE compte 28 membres » ou « le Royaume-Uni est dans l'UE » est une erreur factuelle sanctionnée.
\item \textbf{Oublier les éléments obligatoires du croquis.} Un croquis sans \textbf{titre}, sans \textbf{légende organisée} ou sans \textbf{flèche du nord} perd systématiquement des points, même s'il est exact. La légende doit être hiérarchisée et chaque symbole doit figurer sur la carte.
\item \textbf{Se tromper dans les calculs de pourcentages.} Erreurs classiques : oublier de multiplier par 100 ; calculer une part sur un mauvais total ; additionner des pourcentages sans vérifier qu'ils font 100 \%. Pose toujours la formule : $\text{part} = \dfrac{\text{partie}}{\text{total}} \times 100$, puis vérifie la cohérence du résultat.
\item \textbf{Paraphraser le document au lieu de le commenter.} Recopier les chiffres du texte sans les expliquer par tes connaissances (marché unique, atelier du monde, hydrocarbures) ne vaut rien : chaque élément cité doit être relié à une notion du cours.
\item \textbf{Présenter la Russie comme une grande puissance économique.} C'est l'erreur d'appréciation la plus fréquente : la Russie est un foyer \textbf{énergétique et militaire}, mais son PIB est modeste (environ 6 fois inférieur à celui de la Chine). Nuancer la nature des puissances est précisément ce qui est attendu.
\end{enumerate}
\end{attention}

\newpage

\coursec{Exercices}

\courssub{Je vérifie mes connaissances}

\begin{exercice}[QCM : vocabulaire et repères]
Pour chaque question, une seule réponse est exacte. Recopie la lettre correspondante.
\begin{enumerate}
\item La \cle{Triade} désigne :
\begin{enumerate}
\item l'Amérique du Nord, l'Europe occidentale et le Japon ;
\item les États-Unis, la Chine et la Russie ;
\item l'ONU, le FMI et l'OMC ;
\item les trois plus grandes places boursières mondiales.
\end{enumerate}
\item L'accord \cle{USMCA} (2018) lie :
\begin{enumerate}
\item les États-Unis, le Mexique et le Canada ;
\item les États-Unis, le Mexique et Cuba ;
\item les États-Unis, le Maroc et le Canada ;
\item les pays d'Amérique centrale.
\end{enumerate}
\item Le premier exportateur mondial de marchandises est :
\begin{enumerate}
\item les États-Unis ; \item l'Allemagne ; \item la Chine ; \item le Japon.
\end{enumerate}
\item L'\cle{effet Bruxelles} désigne :
\begin{enumerate}
\item la concentration des sièges de FMN en Belgique ;
\item la capacité de l'UE à imposer ses normes au reste du monde par son marché ;
\item l'afflux de migrants vers Bruxelles ;
\item la politique agricole commune.
\end{enumerate}
\item La \cle{BRI} (\og nouvelles routes de la soie \fg{}) est une stratégie :
\begin{enumerate}
\item japonaise ; \item coréenne ; \item chinoise ; \item russe.
\end{enumerate}
\item Un \cle{choke point} est :
\begin{enumerate}
\item un point de passage stratégique et vulnérable des flux mondiaux (détroit, canal) ;
\item un paradis fiscal ;
\item une zone de libre-échange ;
\item un centre de recherche.
\end{enumerate}
\end{enumerate}
\end{exercice}

\begin{exercice}[Vrai ou faux ? Justifie]
Pour chaque affirmation, indique si elle est vraie ou fausse et justifie en deux lignes maximum.
\begin{enumerate}
\item La Russie est un foyer moteur de l'économie mondiale au même titre que les États-Unis ou la Chine.
\item L'Union européenne est le premier marché intégré du monde par le droit et le marché unique.
\item Le Japon reste la deuxième puissance économique mondiale derrière les États-Unis.
\item Le \cle{friend-shoring} consiste à relocaliser ses chaînes de valeur dans des pays jugés politiquement fiables.
\item La Corée du Sud est spécialisée dans les semi-conducteurs et les écrans (Samsung, SK Hynix).
\item L'OPEP+, qui réunit l'OPEP et la Russie, n'a aucune influence sur les prix mondiaux de l'énergie.
\end{enumerate}
\end{exercice}

\begin{exercice}[Définitions éclair]
Définis précisément, en deux ou trois lignes chacun, les notions suivantes : \cle{Triade} ; \cle{IDE} (investissements directs à l'étranger) ; \cle{DIT} (division internationale du travail) ; \cle{corridors TEN-T} ; \cle{Green Deal} européen ; \cle{ZLECAf}.
\end{exercice}

\begin{exercice}[Calculs sur indicateurs]
On donne les valeurs suivantes (ordres de grandeur, 2023) :
\begin{center}
\begin{tabularx}{\linewidth}{|l|X|X|X|}
\hline
\rowcolor{popblueL} \textbf{Foyer} & \textbf{PIB (milliards USD)} & \textbf{Population (millions)} & \textbf{Dépenses R\&D (\% du PIB)} \\
\hline
États-Unis & \num{27000} & 335 & 3,5 \\
\hline
Chine & \num{17800} & \num{1410} & 2,4 \\
\hline
UE (27) & \num{18300} & 450 & 2,3 \\
\hline
Japon & \num{4200} & 125 & 3,3 \\
\hline
Corée du Sud & \num{1700} & 52 & 4,9 \\
\hline
Russie & \num{2000} & 144 & 1,1 \\
\hline
\end{tabularx}
\end{center}
\begin{enumerate}
\item Calcule la part de la Chine dans le PIB cumulé des six foyers (arrondis à 0,1 \%).
\item Calcule le PIB par habitant des États-Unis et de la Russie (en USD, arrondis au millier).
\item Quel foyer investit proportionnellement le plus dans la recherche ? Quel enseignement en tires-tu pour la hiérarchie future des foyers ?
\end{enumerate}
\end{exercice}

\courssub{Je m'entraîne}

\begin{exercice}[Localiser sur un fond de carte mondial]
Sur un planisphère vierge, place et nomme : les trois pôles historiques de la Triade ; les détroits d'Ormuz, de Malacca et de Taïwan ; les canaux de Suez et de Panama ; les ports de Shanghai, Singapour, Rotterdam et Los Angeles ; les mégapoles de Tokyo, Séoul, Shenzhen et la mégalopole américaine (BosWash). Rédige une légende organisée en trois rubriques.
\end{exercice}

\begin{exercice}[Classer et interpréter un tableau statistique]
On donne les parts dans les exportations mondiales de marchandises (en \%, ordres de grandeur) :
\begin{center}
\begin{tabularx}{\linewidth}{|l|X|X|X|}
\hline
\rowcolor{popblueL} \textbf{Foyer} & \textbf{2000} & \textbf{2010} & \textbf{2022} \\
\hline
UE (27, extra-UE) & 18 & 15 & 14 \\
\hline
États-Unis & 12 & 8 & 8 \\
\hline
Chine & 4 & 10 & 14 \\
\hline
Japon & 7 & 5 & 3 \\
\hline
Corée du Sud & 2 & 3 & 3 \\
\hline
Russie & 1,5 & 2,5 & 2 \\
\hline
\end{tabularx}
\end{center}
\begin{enumerate}
\item Classe les foyers en trois catégories : parts en hausse, en stagnation, en baisse.
\item Calcule le taux de variation de la part chinoise entre 2000 et 2022.
\item Rédige un commentaire de 10 lignes intitulé : \og Le basculement du commerce mondial vers l'Asie de l'Est \fg{}.
\end{enumerate}
\end{exercice}

\begin{exercice}[Schématiser une filière]
À l'aide d'un schéma fléché (cases et flèches), représente la chaîne de valeur mondiale d'un smartphone : conception (Californie), puces (Taïwan, Corée du Sud), écrans (Corée du Sud, Japon), assemblage (Chine, Vietnam), distribution (monde). Pour chaque étape, indique la nature de la valeur ajoutée (faible, moyenne, forte) et conclus en quatre lignes sur la DIT en Asie de l'Est.
\end{exercice}

\begin{exercice}[Commentaire guidé d'une photographie]
On te donne la description d'une photographie satellite de nuit : \og De vastes auras lumineuses continues éclairent le littoral chinois de Shanghai à Shenzhen, la plaine du Kanto (Tokyo) et l'axe Séoul--Pusan ; en Europe, une tache lumineuse dense s'étend de Londres à Milan en passant par la Ruhr ; en Amérique du Nord, deux façades brillent (BosWash à l'est, Californie à l'ouest) ; la Russie n'apparaît que par des points lumineux isolés (Moscou, Saint-Pétersbourg) et des torchères pétrolières de Sibérie. \fg{}
\begin{enumerate}
\item Décris méthodiquement le document (nature, ce qu'on y voit, oppositions).
\item Quels enseignements tires-tu sur la localisation des foyers moteurs et sur la nature particulière de la puissance russe ?
\end{enumerate}
\end{exercice}

\courssub{Je me prépare au Bac}

\begin{exercice}[Sujet type Bac — Croquis : l'Union européenne, foyer moteur]
\textbf{Consigne.} \og À l'aide des documents fournis et de tes connaissances, réalise un croquis commenté montrant que l'Union européenne est un foyer moteur de l'économie mondiale. \fg{}

\textbf{Document 1.} Indicateurs de l'UE (ordres de grandeur, 2022) : PIB $\approx$ \num{18300} milliards USD ; environ 14 \% des exportations mondiales de marchandises (hors échanges intra-UE) ; premier partenaire commercial du Cameroun ; 450 millions de consommateurs.

\textbf{Document 2.} Éléments à placer sur le fond de carte fourni : la \og dorsale \fg{} européenne (de Londres à Milan via la Ruhr et le Benelux) ; les ports de Rotterdam, Anvers, Hambourg, Marseille ; les corridors TEN-T (rhénan, méditerranéen, atlantique, Baltique--Adriatique) ; les métropoles (Paris, Berlin, Madrid, Milan, Amsterdam) ; les périphéries moins dynamiques (sud de l'Italie, ouest de l'Espagne, Grèce).

\textbf{Document 3.} Extraits : objectifs du Green Deal (neutralité carbone en 2050) ; stock d'IDE entrants de l'UE parmi les premiers mondiaux.

\begin{enumerate}
\item Propose un titre et une légende organisée (au moins trois rubriques : cœur intégré, interfaces portuaires, corridors et métropoles).
\item Réalise le croquis sur le fond de carte fourni (orientation, échelle indicative).
\item Rédige un commentaire de 15 lignes montrant les atouts du foyer européen mais aussi ses limites (vieillissement, dépendances énergétique et numérique).
\end{enumerate}
\end{exercice}

\begin{exercice}[Sujet type Bac — Dissertation]
\textbf{Sujet.} \og Les foyers moteurs de l'économie mondiale sont-ils en train de perdre le contrôle de la mondialisation ? \fg{}

Tu mobiliseras notamment : la Triade historique et son évolution ; l'émergence de la Chine et des nouveaux relais de croissance ; la fragmentation géopolitique (guerre en Ukraine, rivalité sino-américaine, sanctions) ; le pouvoir des firmes multinationales face aux États ; le rôle des normes (effet Bruxelles) et des infrastructures (BRI) ; la place des pays du Sud, dont le Cameroun et la ZLECAf.

\begin{enumerate}
\item Dégage la problématique et annonce un plan en deux ou trois parties.
\item Rédige l'introduction (accroche, définition des termes, problématique, annonce du plan).
\item Rédige le développement complet en t'appuyant sur au moins six exemples précis et localisés.
\item Rédige une conclusion ouverte (par exemple sur la mondialisation \og fragmentée \fg{} ou \og à la carte \fg{}).
\end{enumerate}
\end{exercice}

\begin{exercice}[Sujet type Bac — TP : construire et commenter des diagrammes]
À partir du tableau de données ci-dessous (ordres de grandeur, sources : Banque mondiale, OCDE) :

\begin{center}
\begin{tabularx}{\linewidth}{|l|X|X|X|X|}
\hline
\rowcolor{popblueL} \textbf{Indicateur} & \textbf{2000} & \textbf{2010} & \textbf{2023} & \textbf{Unité} \\
\hline
PIB États-Unis & \num{10200} & \num{15000} & \num{27000} & Mds USD \\
\hline
PIB Chine & \num{1200} & \num{6100} & \num{17800} & Mds USD \\
\hline
PIB UE (27) & \num{7300} & \num{14500} & \num{18300} & Mds USD \\
\hline
PIB Japon & \num{4900} & \num{5700} & \num{4200} & Mds USD \\
\hline
R\&D Corée (\% PIB) & 2,3 & 3,5 & 4,9 & \% \\
\hline
R\&D Chine (\% PIB) & 0,9 & 1,7 & 2,4 & \% \\
\hline
\end{tabularx}
\end{center}

\begin{enumerate}
\item Construis un graphique en courbes montrant l'évolution des PIB des quatre premiers foyers entre 2000 et 2023 (axes légendés avec grandeurs et unités).
\item Construis un diagramme en barres comparant les efforts de R\&D de la Corée du Sud et de la Chine aux trois dates.
\item Calcule le taux de croissance du PIB chinois entre 2000 et 2023 et compare-le à celui du Japon.
\item Rédige un commentaire synthétique de cinq lignes intitulé : \og Rattrapage asiatique et recompositions de la hiérarchie mondiale \fg{}.
\item Propose un troisième graphique pertinent que l'on pourrait construire avec des données d'IDE et justifie ton choix en trois lignes.
\end{enumerate}
\end{exercice}

\newpage

\coursec{Corrigés des exercices}

\begin{corrige}[Exercice 1 — QCM : vocabulaire et repères]
\begin{enumerate}
\item \textbf{Réponse a.} La \cle{Triade} désigne les trois pôles historiques de la mondialisation : l'Amérique du Nord, l'Europe occidentale et le Japon (Asie de l'Est élargie ensuite à la Corée du Sud). Les réponses b et c confondent pôles économiques et organisations internationales.
\item \textbf{Réponse a.} L'\cle{USMCA} (Accord États-Unis--Mexique--Canada, entré en vigueur en 2020, signé en 2018) succède à l'ALENA (1994). Cuba est sous embargo américain et le Maroc n'a aucun lien avec cet accord.
\item \textbf{Réponse c.} La \cle{Chine} est le premier exportateur mondial de marchandises (environ 14 \% des exportations mondiales en 2022), devant l'UE (hors échanges intra-UE) et les États-Unis.
\item \textbf{Réponse b.} L'\cle{effet Bruxelles} désigne la capacité normative de l'UE : ses normes (RGPD, normes environnementales, sécurité alimentaire) s'imposent de facto aux entreprises du monde entier qui veulent accéder à son marché de 450 millions de consommateurs.
\item \textbf{Réponse c.} La \cle{BRI} (\emph{Belt and Road Initiative}, lancée en 2013) est la stratégie d'infrastructures et d'influence de la Chine : routes, voies ferrées, ports (Le Pirée, Gwadar, Kribi en partie financé par des prêts chinois).
\item \textbf{Réponse a.} Un \cle{choke point} (point d'étranglement) est un passage obligé et vulnérable des flux mondiaux : détroits d'Ormuz, de Malacca, de Taïwan, canaux de Suez et de Panama. Leur blocage (ex. porte-conteneurs \emph{Ever Given} à Suez, mars 2021) perturbe le commerce mondial.
\end{enumerate}
\end{corrige}

\begin{corrige}[Exercice 2 — Vrai ou faux ? Justifie]
\begin{enumerate}
\item \textbf{Faux.} La Russie est un foyer \emph{énergétique, militaire et pivot eurasiatique}, mais son PIB (environ \num{2000} milliards USD) est comparable à celui de l'Italie ou de la Corée du Sud : elle n'est pas un moteur économique de même rang que les États-Unis (\num{27000} Mds USD) ou la Chine (\num{17800} Mds USD).
\item \textbf{Vrai.} L'UE est le marché intégré le plus abouti : marché unique (libre circulation des biens, services, capitaux, personnes), union douanière, monnaie commune pour 20 membres, droit supranational.
\item \textbf{Faux.} Le Japon a été dépassé par la Chine en 2010 ; il est aujourd'hui quatrième (environ \num{4200} Mds USD), derrière les États-Unis, la Chine et l'Allemagne (ordre de grandeur 2023).
\item \textbf{Vrai.} Le \cle{friend-shoring} consiste à implanter ses chaînes de valeur chez des partenaires politiquement fiables (ex. les États-Unis favorisent le Mexique ou le Vietnam plutôt que la Chine pour les semi-conducteurs).
\item \textbf{Vrai.} La Corée du Sud est un géant des semi-conducteurs et des écrans : Samsung et SK Hynix dominent le marché mondial des mémoires (DRAM, NAND).
\item \textbf{Faux.} L'OPEP+ (OPEP + Russie et alliés) contrôle une part majeure de la production mondiale de pétrole ; ses décisions de quotas influencent directement les prix mondiaux (ex. réductions de production en 2022--2023).
\end{enumerate}
\end{corrige}

\begin{corrige}[Exercice 3 — Définitions éclair]
\begin{itemize}
\item \textbf{\cle{Triade}} : ensemble des trois pôles qui dominent l'économie mondiale depuis les années 1980 — Amérique du Nord, Europe occidentale, Asie de l'Est (Japon puis Corée du Sud) — concentrent l'essentiel des richesses, des IDE et des échanges. Son périmètre s'est élargi avec l'émergence de la Chine.
\item \textbf{\cle{IDE}} (investissements directs à l'étranger) : investissements réalisés par une entreprise (souvent une FMN) pour créer, acquérir ou contrôler une filiale à l'étranger ; ils matérialisent la puissance et la stratégie de localisation des firmes.
\item \textbf{\cle{DIT}} (division internationale du travail) : répartition des étapes de production entre pays selon leurs avantages comparatifs (conception au Nord, assemblage dans les pays à bas coût de main-d'œuvre, matières premières au Sud).
\item \textbf{\cle{Corridors TEN-T}} : réseau transeuropéen de transport (neuf corridors : rhénan, atlantique, méditerranéen, Baltique--Adriatique, etc.) destiné à fluidifier et désenclaver le marché unique européen.
\item \textbf{\cle{Green Deal}} : pacte vert européen (2019) visant la neutralité carbone de l'UE en 2050 ; il fait de la transition écologique un levier de compétitivité et de normes (mécanisme d'ajustement carbone aux frontières).
\item \textbf{\cle{ZLECAf}} : Zone de libre-échange continentale africaine (en vigueur depuis 2021), qui vise à créer un marché unique africain de plus d'un milliard de consommateurs et à renforcer la place du continent, dont le Cameroun, dans la mondialisation.
\end{itemize}
\end{corrige}

\begin{corrige}[Exercice 4 — Calculs sur indicateurs]
\begin{enumerate}
\item PIB cumulé des six foyers :
\[ \num{27000} + \num{17800} + \num{18300} + \num{4200} + \num{1700} + \num{2000} = \num{71000} \text{ milliards USD}. \]
Part de la Chine :
\[ \frac{\num{17800}}{\num{71000}} \times 100 \approx 25{,}07\ \% \approx \textbf{25{,}1\ \%}. \]
\item PIB par habitant :
\begin{align*}
\text{États-Unis : } & \frac{\num{27000}\times 10^9}{335\times 10^6} \approx \num{80597} \approx \textbf{\num{81000} USD/hab.} \\
\text{Russie : } & \frac{\num{2000}\times 10^9}{144\times 10^6} \approx \num{13889} \approx \textbf{\num{14000} USD/hab.}
\end{align*}
Le PIB par habitant américain est donc près de six fois supérieur au russe : la puissance russe n'est pas d'abord économique.
\item La \textbf{Corée du Sud} investit proportionnellement le plus dans la R\&D (4,9 \% du PIB), devant les États-Unis (3,5 \%) et le Japon (3,3 \%). \textbf{Enseignement} : l'effort de recherche prépare la hiérarchie future — l'Asie de l'Est (Corée, Chine à 2,4 \% et en hausse rapide) monte en gamme technologique, tandis que la Russie (1,1 \%) hypothèque sa compétitivité de long terme.
\end{enumerate}
\begin{attention}[Points de vigilance]
Bien convertir milliards et millions avant de diviser ; arrondir seulement à la fin ; ne pas confondre \og part en \% \fg{} et \og taux de variation \fg{}.
\end{attention}
\end{corrige}

\begin{corrige}[Exercice 5 — Localiser sur un fond de carte mondial]
\textbf{Méthode de construction.} Sur le planisphère vierge : (1) colorier en trois teintes les trois pôles de la Triade (Amérique du Nord, Europe occidentale, Asie de l'Est) ; (2) marquer par des symboles en forme de cadenas ou de losanges les \emph{choke points} : détroits d'Ormuz (entre Iran et péninsule Arabique), de Malacca (entre Malaisie et Sumatra), de Taïwan (entre Chine et Taïwan), canaux de Suez (Égypte) et de Panama ; (3) placer par des points carrés les ports : Shanghai (est de la Chine), Singapour (pointe de Malacca), Rotterdam (embouchure du Rhin), Los Angeles (côte pacifique des États-Unis) ; (4) placer par des points ronds les mégapoles : Tokyo (plaine du Kanto, Honshu), Séoul (nord-ouest de la Corée du Sud), Shenzhen (delta de la rivière des Perles, au nord de Hong Kong), et hachurer la mégalopole BosWash (de Boston à Washington, façade atlantique des États-Unis).

\textbf{Légende organisée en trois rubriques :}
\begin{enumerate}
\item \textbf{Les pôles de la Triade} (aplats colorés) : Amérique du Nord ; Europe occidentale ; Asie de l'Est.
\item \textbf{Les points de passage stratégiques} (losanges) : détroits d'Ormuz, de Malacca, de Taïwan ; canaux de Suez et de Panama.
\item \textbf{Les interfaces et métropoles mondiales} : ports (carrés) : Shanghai, Singapour, Rotterdam, Los Angeles ; mégapoles (ronds) : Tokyo, Séoul, Shenzhen ; mégalopole BosWash (hachures).
\end{enumerate}
\begin{attention}[Points de vigilance]
Orienter la carte (flèche du nord), indiquer une échelle indicative, ne pas confondre Shenzhen (Chine) et Séoul (Corée du Sud), et distinguer clairement détroit (naturel) et canal (artificiel).
\end{attention}
\end{corrige}

\begin{corrige}[Exercice 6 — Classer et interpréter un tableau statistique]
\begin{enumerate}
\item \textbf{Classement} (2000 $\rightarrow$ 2022) :
\begin{itemize}
\item \textbf{En hausse} : Chine (4 $\rightarrow$ 14 \%), Corée du Sud (2 $\rightarrow$ 3 \%), Russie (1,5 $\rightarrow$ 2 \%, avec pic en 2010).
\item \textbf{En stagnation} : aucun foyer strictement ; la Corée du Sud stagne depuis 2010 (3 \%).
\item \textbf{En baisse} : UE (18 $\rightarrow$ 14 \%), États-Unis (12 $\rightarrow$ 8 \%), Japon (7 $\rightarrow$ 3 \%).
\end{itemize}
\item Taux de variation de la part chinoise :
\[ \frac{14 - 4}{4} \times 100 = \frac{10}{4} \times 100 = \textbf{+250\ \%}. \]
La part de la Chine a été multipliée par 3,5 entre 2000 et 2022.
\item \textbf{Commentaire (10 lignes)} : En vingt-deux ans, la géographie du commerce mondial s'est profondément recomposée. Les trois pôles historiques de la Triade reculent tous : l'UE passe de 18 \% à 14 \%, les États-Unis de 12 \% à 8 \%, et surtout le Japon s'effondre de 7 \% à 3 \%. À l'inverse, l'Asie de l'Est émergente s'affirme : la Chine, entrée à l'OMC en 2001, voit sa part bondir de 4 \% à 14 \% (+250 \%), devenant le premier exportateur mondial de marchandises, tandis que la Corée du Sud consolide sa place (3 \%) grâce aux semi-conducteurs et à l'électronique. Ce basculement traduit la nouvelle division internationale du travail : l'Asie de l'Est est devenue \og l'atelier du monde \fg{}, puis un pôle d'innovation. La Russie reste marginale (2 \%), exportant surtout de l'énergie. La mondialisation commerciale a donc changé de centre de gravité : de l'Atlantique vers le Pacifique.
\end{enumerate}
\end{corrige}

\begin{corrige}[Exercice 7 — Schématiser une filière]
\begin{popfigure}
\begin{center}
\begin{tikzpicture}[node distance=0.35cm, every node/.style={font=\small},
box/.style={draw=popblue, fill=popblueL, rounded corners, align=center, minimum width=2.6cm, minimum height=0.9cm},
va/.style={font=\footnotesize\itshape, text=popdark}]
\node[box, align=center] (c){Conception\\ Californie (Apple)};
\node[box, right=of c, align=center] (p){Puces\\ Taïwan (TSMC), Corée};
\node[box, right=of p, align=center] (e){Écrans\\ Corée, Japon};
\node[box, right=of e, align=center] (a){Assemblage\\ Chine, Vietnam};
\node[box, right=of a, align=center] (d){Distribution\\ Monde};
\draw[-{Stealth}, thick, poporange] (c) -- (p);
\draw[-{Stealth}, thick, poporange] (p) -- (e);
\draw[-{Stealth}, thick, poporange] (e) -- (a);
\draw[-{Stealth}, thick, poporange] (a) -- (d);
\node[va, below=0.1cm of c] {VA forte};
\node[va, below=0.1cm of p] {VA forte};
\node[va, below=0.1cm of e] {VA moyenne à forte};
\node[va, below=0.1cm of a] {VA faible};
\node[va, below=0.1cm of d] {VA moyenne};
\end{tikzpicture}
\end{center}
\legende{Chaîne de valeur mondiale d'un smartphone : la valeur ajoutée (VA) se concentre aux deux bouts de la chaîne (conception, composants de pointe), l'assemblage n'en capte qu'une faible part.}
\end{popfigure}

\textbf{Conclusion (4 lignes)} : La DIT en Asie de l'Est est hiérarchisée : la conception et les profits restent aux États-Unis (Californie), tandis que l'Asie de l'Est s'est hissée au cœur des étapes à forte valeur ajoutée (puces à Taïwan et en Corée, écrans en Corée et au Japon). La Chine et le Vietnam, longtemps cantonnés à l'assemblage à faible valeur ajoutée, cherchent à monter en gamme. La région forme ainsi un pôle productif intégré, mais dépendant de la conception américaine.
\end{corrige}

\begin{corrige}[Exercice 8 — Commentaire guidé d'une photographie]
\begin{enumerate}
\item \textbf{Description méthodique.} Nature : il s'agit d'une photographie satellite nocturne de la Terre, où les éclairages artificiels matérialisent l'urbanisation et l'activité économique. Ce qu'on y voit : trois immenses concentrations lumineuses — le littoral chinois de Shanghai à Shenzhen, la plaine du Kanto (Tokyo) et l'axe Séoul--Pusan en Asie de l'Est ; une tache dense de Londres à Milan via la Ruhr en Europe ; deux façades brillantes en Amérique du Nord (BosWash à l'est, Californie à l'ouest). Opposition majeure : à ces auras continues s'oppose la Russie, réduite à des points isolés (Moscou, Saint-Pétersbourg) et aux torchères pétrolières de Sibérie.
\item \textbf{Enseignements.} Les foyers moteurs se localisent dans l'hémisphère Nord, sur les littoraux et les grandes façades maritimes, organisés en mégapoles et en corridors continus : la lumière dessine presque exactement la Triade. La nature particulière de la puissance russe apparaît clairement : une puissance ponctuelle (deux métropoles) et extractive (torchères d'hydrocarbures), sans tissu urbain et productif dense — une puissance énergétique et militaire plus qu'un véritable moteur économique intégré.
\end{enumerate}
\begin{attention}[Points de vigilance]
Décrire avant d'interpréter ; ne jamais écrire \og on voit que la Russie est pauvre \fg{} : la faible luminosité traduit aussi la faible densité de population ; toujours annoncer la nature du document en première phrase.
\end{attention}
\end{corrige}

\begin{corrige}[Exercice 9 — Sujet type Bac : croquis de l'Union européenne]
\begin{enumerate}
\item \textbf{Titre proposé} : \og L'Union européenne : un foyer moteur organisé autour d'une dorsale et d'interfaces portuaires mondiales \fg{}.
\textbf{Légende organisée} :
\begin{itemize}
\item \textbf{Rubrique 1 — Le cœur intégré} : la dorsale européenne de Londres à Milan (via Benelux et Ruhr) : aplat coloré.
\item \textbf{Rubrique 2 — Les interfaces portuaires} : Rotterdam, Anvers, Hambourg (façade Nord) ; Marseille (façade méditerranéenne) : symboles d'ancre.
\item \textbf{Rubrique 3 — Corridors et métropoles} : corridors TEN-T (rhénan, atlantique, méditerranéen, Baltique--Adriatique) : flèches ; métropoles (Paris, Berlin, Madrid, Milan, Amsterdam) : points nommés ; périphéries moins dynamiques (Mezzogiorno, ouest espagnol, Grèce) : hachures.
\end{itemize}
\item \textbf{Croquis attendu} (modèle simplifié) :
\begin{popfigure}
\includegraphics[width=\linewidth,height=9.5cm,keepaspectratio]{N13/blue_banana.png}
\legende{Fond de carte pour le croquis de l'armature européenne : la « banane bleue » (aplat bleu), dorsale qui s'étire du sud-est de l'Angleterre à la Lombardie en passant par le Benelux, la vallée du Rhin et la Suisse, sur un fond de pays d'Europe en gris. \textbf{Le croquis attendu situe :} 1. la dorsale européenne (aplat de la carte) ; 2. les ports (Rotterdam, Anvers, Hambourg, Marseille) ; 3. les métropoles (Paris, Berlin, Madrid, Milan, Amsterdam) ; 4. les corridors TEN-T (flèches) ; 5. les périphéries moins dynamiques (Mezzogiorno, ouest espagnol, Grèce), qui se lisent hors de l'aplat bleu. \\ Source : Wikimedia Commons, « Blue Banana », ArnoldPlaton, CC BY-SA 3.0.}
\end{popfigure}
\item \textbf{Commentaire (15 lignes)} : L'Union européenne est un foyer moteur de premier plan : avec environ \num{18300} milliards USD de PIB et 450 millions de consommateurs, elle constitue le premier marché intégré du monde par le droit et le marché unique, et pèse environ 14 \% des exportations mondiales (hors échanges intra-UE). Le croquis montre une organisation polarisée : la richesse se concentre sur la dorsale qui s'étend de Londres à Milan via le Benelux et la Ruhr, cœur industriel et financier de l'Europe. Ce cœur s'ouvre sur le monde par de puissantes interfaces portuaires — Rotterdam, premier port d'Europe, Anvers et Hambourg sur la façade Nord, Marseille en Méditerranée — et par un réseau de métropoles mondiales (Paris, Berlin, Milan, Amsterdam). Les corridors TEN-T (rhénan, atlantique, méditerranéen, Baltique--Adriatique) relient ces pôles et désenclavent les périphéries, qui restent toutefois moins dynamiques (Mezzogiorno, ouest espagnol, Grèce) : l'intégration n'efface pas les inégalités internes. Premier partenaire commercial du Cameroun, l'UE rayonne aussi par ses normes (effet Bruxelles) et par le Green Deal, qui vise la neutralité carbone en 2050. Mais ce foyer a des limites : vieillissement démographique, dépendance énergétique (révélée par la guerre en Ukraine et la rupture avec le gaz russe) et dépendance numérique vis-à-vis des États-Unis et de l'Asie. L'UE reste donc un géant commercial et normatif, mais un nain technologique relatif, dont la puissance future dépend de sa capacité d'innovation.
\end{enumerate}
\textbf{Barème indicatif (sur 20)} : titre et légende organisée : 4 pts ; exactitude de la localisation : 6 pts ; qualité graphique (orientation, échelle, propreté) : 4 pts ; commentaire structuré (atouts/limites) : 6 pts.
\begin{attention}[Points de vigilance]
Ne jamais colorier toute l'UE en \og cœur \fg{} : la dorsale est une bande limitée ; citer au moins une limite dans le commentaire ; la légende doit être hiérarchisée, pas une simple liste.
\end{attention}
\end{corrige}

\begin{corrige}[Exercice 10 — Sujet type Bac : dissertation]
\begin{enumerate}
\item \textbf{Problématique} : Les pôles historiques de la mondialisation (Triade) conservent-ils la maîtrise des flux, des normes et des chaînes de valeur, ou bien l'émergence de nouvelles puissances et la fragmentation géopolitique leur échappent-elles ? \textbf{Plan en trois parties} : I. Une domination encore réelle de la Triade ; II. Mais un contrôle contesté par l'émergence chinoise et la fragmentation ; III. Vers une mondialisation fragmentée, à recomposer avec le Sud.
\item \textbf{Introduction} : En mars 2021, le blocage du canal de Suez par l'\emph{Ever Given} a révélé la fragilité d'une mondialisation organisée par quelques foyers. Les foyers moteurs sont les pôles qui concentrent richesse, IDE, innovation et capacité à fixer les règles du commerce mondial ; la mondialisation est le processus d'intensification des échanges et d'interdépendance à l'échelle planétaire. Depuis les années 1980, la Triade (Amérique du Nord, Europe occidentale, Asie de l'Est) structure ce système. Mais l'émergence chinoise, la guerre en Ukraine et la puissance des firmes brouillent ce contrôle. Les foyers moteurs sont-ils en train de perdre le contrôle de la mondialisation ? Nous verrons que leur domination reste forte (I), qu'elle est contestée par la Chine et la fragmentation géopolitique (II), et que la mondialisation se recompose plutôt qu'elle ne s'effondre, ouvrant des marges au Sud (III).
\item \textbf{Développement (exemples précis)} :
\textbf{I.} Les États-Unis conservent la première place financière (Wall Street, dollar), technologique (GAFA de la Silicon Valley) et le premier PIB mondial (\num{27000} Mds USD) ; l'UE impose ses normes par l'\cle{effet Bruxelles} (RGPD repris mondialement) et reste le premier partenaire commercial du Cameroun ; le Japon et la Corée dominent des segments critiques (écrans, mémoires Samsung/SK Hynix).
\textbf{II.} Mais la Chine, premier exportateur mondial (14 \%), déploie la \cle{BRI} (ports du Pirée, de Gwadar, financements à Kribi) et capte la conception de produits (Huawei, BYD) ; Taïwan concentre via TSMC l'essentiel des puces avancées — un \emph{choke point} technologique ; la guerre en Ukraine (2022) et les sanctions ont réorienté les flux énergétiques russes vers l'Asie ; la rivalité sino-américaine impose découplage, \emph{friend-shoring} et guerre des semi-conducteurs ; les FMN (Apple, TotalEnergies) décident parfois contre les États.
\textbf{III.} La mondialisation devient \og fragmentée \fg{} : blocs régionaux (USMCA, UE, RCEP asiatique), sanctions, relocalisations. Le Sud cherche sa place : la \cle{ZLECAf} (2021) veut créer un marché africain intégré ; le Cameroun mise sur son port de Kribi et sa position de corridor pour le Tchad et la Centrafrique. Le contrôle ne disparaît pas : il se partage et se négocie.
\item \textbf{Conclusion} : Les foyers moteurs n'ont pas perdu le contrôle de la mondialisation, mais ils ne l'exercent plus seuls : la puissance s'est diffusée vers l'Asie orientale et se dispute dans un monde fragmenté en blocs rivaux. Reste à savoir si cette mondialisation \og à la carte \fg{} offrira aux pays du Sud, comme le Cameroun, l'occasion de devenir acteurs plutôt que simples territoires de passage.
\end{enumerate}
\textbf{Barème indicatif (sur 20)} : introduction (accroche, définitions, problématique, plan) : 4 pts ; développement structuré avec au moins six exemples localisés : 10 pts ; cohérence de l'argumentation : 3 pts ; conclusion ouverte, langue : 3 pts.
\begin{attention}[Points de vigilance]
Définir les deux termes du sujet ; chaque partie doit comporter des exemples \emph{précis et localisés} (entreprises, ports, accords) ; éviter le catalogue sans fil directeur ; la conclusion doit répondre à la problématique avant de s'ouvrir.
\end{attention}
\end{corrige}

\begin{corrige}[Exercice 11 — Sujet type Bac : TP diagrammes]
\begin{enumerate}
\item \textbf{Graphique en courbes des PIB} :
\begin{popfigure}
\begin{center}
\begin{tikzpicture}
\begin{axis}[width=0.85\linewidth, height=6cm,
xlabel={Ann\'ee}, ylabel={PIB (milliards USD)},
xtick={2000,2010,2023}, xmin=1998, xmax=2025, ymin=0, ymax=29000,
legend style={at={(0.02,0.98)}, anchor=north west, font=\footnotesize},
grid=major, grid style={popline!50}]
\addplot[very thick, popblue, mark=*] coordinates {(2000,10200) (2010,15000) (2023,27000)};
\addlegendentry{\'Etats-Unis}
\addplot[very thick, poporange, mark=*] coordinates {(2000,1200) (2010,6100) (2023,17800)};
\addlegendentry{Chine}
\addplot[very thick, popgreen, mark=*] coordinates {(2000,7300) (2010,14500) (2023,18300)};
\addlegendentry{UE (27)}
\addplot[very thick, poppurple, mark=*] coordinates {(2000,4900) (2010,5700) (2023,4200)};
\addlegendentry{Japon}
\end{axis}
\end{tikzpicture}
\end{center}
\legende{\'Evolution des PIB de quatre foyers moteurs, 2000--2023 (ordres de grandeur ; sources : Banque mondiale, OCDE).}
\end{popfigure}
\item \textbf{Diagramme en barres des efforts de R\&D} :
\begin{popfigure}
\begin{center}
\begin{tikzpicture}
\begin{axis}[width=0.75\linewidth, height=5.5cm, ybar, bar width=0.5cm,
ylabel={R\&D (\% du PIB)}, symbolic x coords={2000,2010,2023}, xtick=data,
ymin=0, ymax=6, legend style={at={(0.02,0.98)}, anchor=north west, font=\footnotesize},
grid=major, grid style={popline!50}]
\addplot[fill=poppink, draw=poppink!70!black] coordinates {(2000,2.3) (2010,3.5) (2023,4.9)};
\addlegendentry{Cor\'ee du Sud}
\addplot[fill=popgold, draw=popgold!70!black] coordinates {(2000,0.9) (2010,1.7) (2023,2.4)};
\addlegendentry{Chine}
\end{axis}
\end{tikzpicture}
\end{center}
\legende{Effort de recherche-d\'eveloppement de la Cor\'ee du Sud et de la Chine, 2000--2023 (en \% du PIB).}
\end{popfigure}
\item \textbf{Taux de croissance} :
\begin{align*}
\text{Chine : } & \frac{\num{17800} - \num{1200}}{\num{1200}} \times 100 = \frac{\num{16600}}{\num{1200}} \times 100 \approx \textbf{+\num{1383}\ \%} \quad (\text{PIB multipli\'e par } \approx 14{,}8). \\
\text{Japon : } & \frac{\num{4200} - \num{4900}}{\num{4900}} \times 100 = \frac{\num{-700}}{\num{4900}} \times 100 \approx \textbf{-14{,}3\ \%}.
\end{align*}
Sur vingt-trois ans, le PIB chinois a été presque multiplié par quinze, tandis que celui du Japon a reculé d'environ un septième : le contraste est saisissant.
\item \textbf{Commentaire (5 lignes)} : \og Rattrapage asiatique et recompositions de la hiérarchie mondiale \fg{}. Entre 2000 et 2023, la Chine a accompli un rattrapage spectaculaire (+\num{1383} \%), dépassant l'UE en 2010 et se rapprochant des États-Unis, tandis que le Japon déclinait (-14,3 \%). Les courbes des PIB montrent un basculement du centre de gravité économique vers l'Asie de l'Est. Le diagramme en barres révèle le moteur de ce rattrapage : l'effort de R\&D coréen (4,9 \%, record mondial relatif) et chinois (2,4 \%, en hausse continue) prépare la domination technologique de demain. La hiérarchie mondiale se recompose : la Triade historique cède du terrain face à un pôle asiatique désormais innovant et non plus seulement manufacturier.
\item \textbf{Troisième graphique proposé} : un diagramme en barres groupées des stocks d'IDE entrants et sortants des principaux foyers (États-Unis, UE, Chine) à deux dates (2000 et 2023). \textbf{Justification} : les IDE mesurent l'attractivité (IDE entrants) et la capacité d'investissement à l'étranger (IDE sortants) ; les comparer montrerait que la Chine est devenue non seulement un foyer récepteur mais aussi un investisseur mondial (BRI), confirmant le rattrapage asiatique par un autre indicateur que le PIB.
\end{enumerate}
\textbf{Barème indicatif (sur 20)} : courbes correctes avec axes légendés (grandeurs, unités) : 5 pts ; diagramme en barres : 4 pts ; calculs exacts et comparés : 4 pts ; commentaire structuré : 5 pts ; proposition justifiée : 2 pts.
\begin{attention}[Points de vigilance]
Toujours légender les axes avec la grandeur \emph{et} l'unité ; un taux de croissance négatif (Japon) doit être signalé comme un recul, pas une \og croissance négative \fg{} maladroite ; vérifier l'échelle verticale pour que la courbe chinoise reste lisible.
\end{attention}
\end{corrige}

\newpage

\coursec{Fiche bilan}

\begin{popfigure}
\begin{tikzpicture}[mindmap, grow cyclic, every node/.style={concept, text=white, font=\small\bfseries},
root concept/.append style={concept color=popdark, font=\large\bfseries, minimum size=3.2cm},
level 1/.append style={level distance=4.6cm, sibling angle=51.5},
level 2/.append style={level distance=2.6cm, sibling angle=40, font=\scriptsize\bfseries, minimum size=1.5cm}]
\node[root concept]{Les foyers moteurs de l'économie mondiale}
  child[concept color=popblue] { node {La Triade}
    child { node {3 pôles : Amérique du Nord, UE, Asie de l'Est} }
    child { node {$\approx$ 60--70\,\% des richesses mondiales} }
  }
  child[concept color=popgold] { node {Union européenne}
    child { node {27 États, marché unique, euro (20 pays)} }
    child { node {1\textsuperscript{er} foyer intégré par le droit} }
  }
  child[concept color=poporange] { node {Amérique du Nord}
    child { node {USA : 1\textsuperscript{re} puissance (PIB, dollar, GAFAM)} }
    child { node {USMCA : USA--Canada--Mexique} }
  }
  child[concept color=poppink] { node {Asie de l'Est}
    child { node {Chine : atelier du monde, 2\textsuperscript{e} PIB} }
    child { node {Japon et Corée du Sud : high-tech} }
  }
  child[concept color=popgreen] { node {Russie}
    child { node {Énergie : gaz, pétrole} }
    child { node {Puissance militaire, pivot eurasiatique} }
  }
  child[concept color=popmuted] { node {Interdépendances}
    child { node {Flux commerciaux et financiers} }
    child { node {Rivalités et coopérations} }
  };
\end{tikzpicture}
\legende{Carte mentale de synthèse : les foyers moteurs de l'économie mondiale (Triade, quatre grands pôles étudiés et leurs interdépendances).}
\end{popfigure}

\begin{bilan}
\textbf{Définitions clés à connaître par cœur}
\begin{itemize}
  \item \cle{Triade} : les trois pôles dominants de l'économie mondiale (Amérique du Nord, Union européenne, Asie de l'Est), qui concentrent l'essentiel de la production, des échanges et des flux financiers mondiaux.
  \item \cle{Foyer moteur} : espace qui impulse la dynamique de l'économie mondiale par sa production, sa consommation, ses innovations et ses investissements.
  \item \cle{Marché unique} (UE) : espace de libre circulation des biens, des services, des capitaux et des personnes entre les 27 États membres.
  \item \cle{USMCA} (ex-ALENA) : accord de libre-échange entre les États-Unis, le Mexique et le Canada, entré en vigueur en 2020.
  \item \cle{Atelier du monde} : surnom de la Chine, premier exportateur mondial de biens manufacturés.
  \item \cle{Mondialisation} : mise en interdépendance croissante des économies par les flux de biens, de capitaux, d'informations et d'hommes.
\end{itemize}

\textbf{Repères chiffrés et faits majeurs (ordres de grandeur à retenir)}
\begin{center}
\begin{tabularx}{\linewidth}{|l|X|}
\hline
\rowcolor{popblueL} \textbf{Foyer} & \textbf{Faits saillants} \\
\hline
Union européenne & 27 membres ; environ 450 millions d'habitants ; euro partagé par 20 pays ; premier ensemble commercial mondial à parité avec les USA. \\
\hline
États-Unis & 1\textsuperscript{re} puissance économique (PIB de l'ordre de 25 à 27 000 milliards de dollars) ; dollar, monnaie de référence ; GAFAM, Wall Street. \\
\hline
Chine & 2\textsuperscript{e} PIB mondial ; 1\textsuperscript{er} exportateur de marchandises ; nouvelles routes de la soie (BRI, 2013). \\
\hline
Japon / Corée du Sud & 3\textsuperscript{e} et environ 13\textsuperscript{e} économies mondiales ; leaders en électronique, automobile, robotique (Toyota, Sony, Samsung, Hyundai). \\
\hline
Russie & Parmi les premiers producteurs mondiaux de gaz et de pétrole ; arsenal nucléaire majeur ; pont entre Europe et Asie. \\
\hline
\end{tabularx}
\end{center}

\textbf{Méthodes express}
\begin{itemize}
  \item \emph{Croquis de synthèse} : fond de planisphère simplifié, colorier les trois pôles de la Triade, flèches proportionnelles pour les flux, légende ordonnée (pôles, flux, interfaces : ports, détroits).
  \item \emph{Commentaire de document} : localiser $\rightarrow$ décrire les faits chiffrés $\rightarrow$ expliquer (atouts, politiques, histoire) $\rightarrow$ ouvrir (limites, rivalités).
  \item \emph{Comparaison de foyers} : toujours croiser trois critères : richesse (PIB), échanges (part du commerce mondial), innovation (brevets, firmes transnationales).
\end{itemize}
\end{bilan}

\begin{aretenir}[Checklist avant le Bac]
\begin{itemize}
  \item $\square$ Je sais définir la Triade et citer ses trois pôles.
  \item $\square$ Je connais les étapes clés de la construction européenne (1957 Rome, 1992 Maastricht, élargissements, Brexit 2020).
  \item $\square$ Je sais expliquer ce qu'est le marché unique et la zone euro.
  \item $\square$ Je connais les fondements de la puissance américaine (PIB, dollar, GAFAM, FTN, universités).
  \item $\square$ Je sais ce qu'est l'USMCA et je sais citer ses trois membres.
  \item $\square$ Je sais expliquer l'essor de la Chine (réformes de 1978, ZES, BRI) et ses limites.
  \item $\square$ Je connais les spécialisations du Japon et de la Corée du Sud.
  \item $\square$ Je sais présenter les atouts de la Russie (énergie, armée, position eurasiatique) et ses faiblesses (démographie, dépendance aux hydrocarbures).
  \item $\square$ Je sais construire un croquis mondial avec légende, flèches de flux et titre.
  \item $\square$ Je maîtrise quelques ordres de grandeur fiables (PIB, parts du commerce mondial) sans inventer de chiffres précis.
  \item $\square$ Je sais situer le Cameroun et la CEMAC face à ces foyers (partenaires commerciaux : Chine, UE).
\end{itemize}
\end{aretenir}

\findecours

\end{document}
